\documentclass[12pt]{report}
\usepackage[utf8x]{inputenc}

%Warning: tipa declares many non-standard macros used by utf8x to
%interpret utf8 characters but extra packages might have to be added
%such as "textgreek" for Greek letters not already in tipa
%or "stmaryrd" for mathematical symbols.
%Utf8 codes missing a LaTeX interpretation can be defined by using
%\DeclareUnicodeCharacter{code}{interpretation}.
%Use coqdoc's option -p to add new packages or declarations.
\usepackage{tipa}

\usepackage[T1]{fontenc}
\usepackage{fullpage}
\usepackage{coqdoc}
\usepackage{amsmath,amssymb}
\usepackage{url}
\begin{document}
%%%%%%%%%%%%%%%%%%%%%%%%%%%%%%%%%%%%%%%%%%%%%%%%%%%%%%%%%%%%%%%%%
%% This file has been automatically generated with the command
%% /nix/store/5nhkcrnn3hzlx1rsgy6zxlg2lrla4f5y-coq-with-packages-9.1.1/bin-orig/coqdoc -latex -toc -interpolate -utf8 -Q . SECF Maps.v Imp.v Equiv.v Hoare.v Hoare2.v Preface.v Noninterference.v StaticIFC.v SpecCT.v Postscript.v Bib.v MapsTest.v ImpTest.v EquivTest.v HoareTest.v Hoare2Test.v PrefaceTest.v NoninterferenceTest.v StaticIFCTest.v SpecCTTest.v PostscriptTest.v BibTest.v 
%%%%%%%%%%%%%%%%%%%%%%%%%%%%%%%%%%%%%%%%%%%%%%%%%%%%%%%%%%%%%%%%%
\coqlibrary{SECF.Imp}{Library }{SECF.Imp}

\begin{coqdoccode}
\end{coqdoccode}
\section{Imp: Simple Imperative Programs}



 In this chapter, we take a more serious look at how to use Rocq as
    a tool to study other things.  Our case study is a \textit{simple
    imperative programming language} called Imp, embodying a tiny core
    fragment of conventional mainstream languages such as C and Java.


    Here is a familiar mathematical function written in Imp.


       Z := X;
       Y := 1;
       while Z <> 0 do
         Y := Y * Z;
         Z := Z - 1
       end


 We concentrate here on defining the \textit{syntax} and \textit{semantics} of
    Imp; later, in \textit{Programming Language Foundations} (\textit{Software
    Foundations}, volume 2), we develop a theory of \textit{program
    equivalence} and introduce \textit{Hoare Logic}, a popular logic for
    reasoning about imperative programs. 
\begin{coqdoccode}
\coqdocemptyline
\coqdocnoindent
\coqdockw{Set} \coqdocvar{Warnings} "-notation-overridden".\coqdoceol
\coqdocnoindent
\coqdockw{From} \coqdocvar{Stdlib} \coqdockw{Require} \coqdockw{Import} \coqdoclibrary{Bool}.\coqdoceol
\coqdocnoindent
\coqdockw{From} \coqdocvar{Stdlib} \coqdockw{Require} \coqdockw{Import} \coqdoclibrary{Init.Nat}.\coqdoceol
\coqdocnoindent
\coqdockw{From} \coqdocvar{Stdlib} \coqdockw{Require} \coqdockw{Import} \coqdoclibrary{Arith}.\coqdoceol
\coqdocnoindent
\coqdockw{From} \coqdocvar{Stdlib} \coqdockw{Require} \coqdockw{Import} \coqdoclibrary{EqNat}. \coqdockw{Import} \coqdocmodule{Nat}.\coqdoceol
\coqdocnoindent
\coqdockw{From} \coqdocvar{Stdlib} \coqdockw{Require} \coqdockw{Import} \coqdoclibrary{Lia}.\coqdoceol
\coqdocnoindent
\coqdockw{From} \coqdocvar{Stdlib} \coqdockw{Require} \coqdockw{Import} \coqdoclibrary{List}. \coqdockw{Import} \coqdocmodule{ListNotations}.\coqdoceol
\coqdocnoindent
\coqdockw{From} \coqdocvar{Stdlib} \coqdockw{Require} \coqdockw{Import} \coqdoclibrary{Strings.String}.\coqdoceol
\coqdocnoindent
\coqdockw{From} \coqdocvar{SECF} \coqdockw{Require} \coqdockw{Import} \coqref{SECF.Maps}{\coqdoclibrary{Maps}}.\coqdoceol
\coqdocemptyline
\end{coqdoccode}
\section{Arithmetic and Boolean Expressions}



 We'll present Imp in three parts: first a core language of
    \textit{arithmetic and boolean expressions}, then an extension of these
    with \textit{variables}, and finally a language of \textit{commands} including
    assignment, conditionals, sequencing, and loops. 
\begin{coqdoccode}
\coqdocemptyline
\end{coqdoccode}
\subsection{Syntax}


\begin{coqdoccode}
\coqdocemptyline
\coqdocnoindent
\coqdockw{Module} \coqdef{SECF.Imp.AExp}{AExp}{\coqdocmodule{AExp}}.\coqdoceol
\coqdocemptyline
\end{coqdoccode}
These two definitions specify the \textit{abstract syntax} of
    arithmetic and boolean expressions. 
\begin{coqdoccode}
\coqdocemptyline
\coqdocnoindent
\coqdockw{Inductive} \coqdef{SECF.Imp.AExp.aexp}{aexp}{\coqdocinductive{aexp}} : \coqdockw{Type} :=\coqdoceol
\coqdocindent{1.00em}
\ensuremath{|} \coqdef{SECF.Imp.AExp.ANum}{ANum}{\coqdocconstructor{ANum}} (\coqdef{SECF.Imp.n:3}{n}{\coqdocbinder{n}} : \coqexternalref{nat}{http://rocq-prover.org/doc/V9.1.1/stdlib//Corelib.Init.Datatypes}{\coqdocinductive{nat}})\coqdoceol
\coqdocindent{1.00em}
\ensuremath{|} \coqdef{SECF.Imp.AExp.APlus}{APlus}{\coqdocconstructor{APlus}} (\coqdef{SECF.Imp.a1:4}{a1}{\coqdocbinder{a1}} \coqdef{SECF.Imp.a2:5}{a2}{\coqdocbinder{a2}} : \coqref{SECF.Imp.aexp:1}{\coqdocinductive{aexp}})\coqdoceol
\coqdocindent{1.00em}
\ensuremath{|} \coqdef{SECF.Imp.AExp.AMinus}{AMinus}{\coqdocconstructor{AMinus}} (\coqdef{SECF.Imp.a1:6}{a1}{\coqdocbinder{a1}} \coqdef{SECF.Imp.a2:7}{a2}{\coqdocbinder{a2}} : \coqref{SECF.Imp.aexp:1}{\coqdocinductive{aexp}})\coqdoceol
\coqdocindent{1.00em}
\ensuremath{|} \coqdef{SECF.Imp.AExp.AMult}{AMult}{\coqdocconstructor{AMult}} (\coqdef{SECF.Imp.a1:8}{a1}{\coqdocbinder{a1}} \coqdef{SECF.Imp.a2:9}{a2}{\coqdocbinder{a2}} : \coqref{SECF.Imp.aexp:1}{\coqdocinductive{aexp}}).\coqdoceol
\coqdocemptyline
\coqdocnoindent
\coqdockw{Inductive} \coqdef{SECF.Imp.AExp.bexp}{bexp}{\coqdocinductive{bexp}} : \coqdockw{Type} :=\coqdoceol
\coqdocindent{1.00em}
\ensuremath{|} \coqdef{SECF.Imp.AExp.BTrue}{BTrue}{\coqdocconstructor{BTrue}}\coqdoceol
\coqdocindent{1.00em}
\ensuremath{|} \coqdef{SECF.Imp.AExp.BFalse}{BFalse}{\coqdocconstructor{BFalse}}\coqdoceol
\coqdocindent{1.00em}
\ensuremath{|} \coqdef{SECF.Imp.AExp.BEq}{BEq}{\coqdocconstructor{BEq}} (\coqdef{SECF.Imp.a1:12}{a1}{\coqdocbinder{a1}} \coqdef{SECF.Imp.a2:13}{a2}{\coqdocbinder{a2}} : \coqref{SECF.Imp.AExp.aexp}{\coqdocinductive{aexp}})\coqdoceol
\coqdocindent{1.00em}
\ensuremath{|} \coqdef{SECF.Imp.AExp.BNeq}{BNeq}{\coqdocconstructor{BNeq}} (\coqdef{SECF.Imp.a1:14}{a1}{\coqdocbinder{a1}} \coqdef{SECF.Imp.a2:15}{a2}{\coqdocbinder{a2}} : \coqref{SECF.Imp.AExp.aexp}{\coqdocinductive{aexp}})\coqdoceol
\coqdocindent{1.00em}
\ensuremath{|} \coqdef{SECF.Imp.AExp.BLe}{BLe}{\coqdocconstructor{BLe}} (\coqdef{SECF.Imp.a1:16}{a1}{\coqdocbinder{a1}} \coqdef{SECF.Imp.a2:17}{a2}{\coqdocbinder{a2}} : \coqref{SECF.Imp.AExp.aexp}{\coqdocinductive{aexp}})\coqdoceol
\coqdocindent{1.00em}
\ensuremath{|} \coqdef{SECF.Imp.AExp.BGt}{BGt}{\coqdocconstructor{BGt}} (\coqdef{SECF.Imp.a1:18}{a1}{\coqdocbinder{a1}} \coqdef{SECF.Imp.a2:19}{a2}{\coqdocbinder{a2}} : \coqref{SECF.Imp.AExp.aexp}{\coqdocinductive{aexp}})\coqdoceol
\coqdocindent{1.00em}
\ensuremath{|} \coqdef{SECF.Imp.AExp.BNot}{BNot}{\coqdocconstructor{BNot}} (\coqdef{SECF.Imp.b:20}{b}{\coqdocbinder{b}} : \coqref{SECF.Imp.bexp:10}{\coqdocinductive{bexp}})\coqdoceol
\coqdocindent{1.00em}
\ensuremath{|} \coqdef{SECF.Imp.AExp.BAnd}{BAnd}{\coqdocconstructor{BAnd}} (\coqdef{SECF.Imp.b1:21}{b1}{\coqdocbinder{b1}} \coqdef{SECF.Imp.b2:22}{b2}{\coqdocbinder{b2}} : \coqref{SECF.Imp.bexp:10}{\coqdocinductive{bexp}}).\coqdoceol
\coqdocemptyline
\end{coqdoccode}
In this chapter, we'll mostly elide the translation from the
    concrete syntax that a programmer would actually write to these
    abstract syntax trees -- the process that, for example, would
    translate the string "1 + 2 \ensuremath{\times} 3" to the AST


      APlus (ANum 1) (AMult (ANum 2) (ANum 3)).


    The optional chapter \coqdocvar{ImpParser} develops a simple lexical
    analyzer and parser that can perform this translation.  You do not
    need to understand that chapter to understand this one, but if you
    haven't already taken a course where these techniques are
    covered (e.g., a course on compilers) you may want to skim it. 

 For comparison, here's a conventional BNF (Backus-Naur Form)
    grammar defining the same abstract syntax:


    a := nat
        | a + a
        | a - a
        | a * a


    b := true
        | false
        | a = a
        | a <> a
        | a <= a
        | a > a
        | \~{} b
        | b \&\& b


 Compared to the Rocq version above...



\begin{itemize}
\item  The BNF is more informal -- for example, it gives some
         suggestions about the surface syntax of expressions (like the
         fact that the addition operation is written with an infix
         +) while leaving other aspects of lexical analysis and
         parsing (like the relative precedence of +, -, and \ensuremath{\times},
         the use of parens to group subexpressions, etc.)
         unspecified.  Some additional information -- and human
         intelligence -- would be required to turn this description
         into a formal definition, e.g., for implementing a compiler.


         The Rocq version consistently omits all this information and
         concentrates on the abstract syntax only.



\item  Conversely, the BNF version is lighter and easier to read.
         Its informality makes it flexible, a big advantage in
         situations like discussions at the blackboard, where
         conveying general ideas is more important than nailing down
         every detail precisely.


         Indeed, there are dozens of BNF-like notations and people
         switch freely among them -- usually without bothering to say
         which kind of BNF they're using, because there is no need to:
         a rough-and-ready informal understanding is all that's
         important.

\end{itemize}


    It's good to be comfortable with both sorts of notations: informal
    ones for communicating between humans and formal ones for carrying
    out implementations and proofs. 
\begin{coqdoccode}
\coqdocemptyline
\end{coqdoccode}
\subsection{Evaluation}



 \textit{Evaluating} an arithmetic expression produces a number. 
\begin{coqdoccode}
\coqdocemptyline
\coqdocnoindent
\coqdockw{Fixpoint} \coqdef{SECF.Imp.AExp.aeval}{aeval}{\coqdocdefinition{aeval}} (\coqdef{SECF.Imp.a:23}{a}{\coqdocbinder{a}} : \coqref{SECF.Imp.AExp.aexp}{\coqdocinductive{aexp}}) : \coqexternalref{nat}{http://rocq-prover.org/doc/V9.1.1/stdlib//Corelib.Init.Datatypes}{\coqdocinductive{nat}} :=\coqdoceol
\coqdocindent{1.00em}
\coqdockw{match} \coqref{SECF.Imp.a:23}{\coqdocvariable{a}} \coqdockw{with}\coqdoceol
\coqdocindent{1.00em}
\ensuremath{|} \coqref{SECF.Imp.AExp.ANum}{\coqdocconstructor{ANum}} \coqdocvar{n} \ensuremath{\Rightarrow} \coqdocvar{n}\coqdoceol
\coqdocindent{1.00em}
\ensuremath{|} \coqref{SECF.Imp.AExp.APlus}{\coqdocconstructor{APlus}}  \coqdocvar{a1} \coqdocvar{a2} \ensuremath{\Rightarrow} \coqexternalref{::nat scope:x '+' x}{http://rocq-prover.org/doc/V9.1.1/stdlib//Corelib.Init.Nat}{\coqdocnotation{(}}\coqref{SECF.Imp.aeval:24}{\coqdocdefinition{aeval}} \coqdocvar{a1}\coqexternalref{::nat scope:x '+' x}{http://rocq-prover.org/doc/V9.1.1/stdlib//Corelib.Init.Nat}{\coqdocnotation{)}} \coqexternalref{::nat scope:x '+' x}{http://rocq-prover.org/doc/V9.1.1/stdlib//Corelib.Init.Nat}{\coqdocnotation{+}} \coqexternalref{::nat scope:x '+' x}{http://rocq-prover.org/doc/V9.1.1/stdlib//Corelib.Init.Nat}{\coqdocnotation{(}}\coqref{SECF.Imp.aeval:24}{\coqdocdefinition{aeval}} \coqdocvar{a2}\coqexternalref{::nat scope:x '+' x}{http://rocq-prover.org/doc/V9.1.1/stdlib//Corelib.Init.Nat}{\coqdocnotation{)}}\coqdoceol
\coqdocindent{1.00em}
\ensuremath{|} \coqref{SECF.Imp.AExp.AMinus}{\coqdocconstructor{AMinus}} \coqdocvar{a1} \coqdocvar{a2} \ensuremath{\Rightarrow} \coqexternalref{::nat scope:x '-' x}{http://rocq-prover.org/doc/V9.1.1/stdlib//Corelib.Init.Nat}{\coqdocnotation{(}}\coqref{SECF.Imp.aeval:24}{\coqdocdefinition{aeval}} \coqdocvar{a1}\coqexternalref{::nat scope:x '-' x}{http://rocq-prover.org/doc/V9.1.1/stdlib//Corelib.Init.Nat}{\coqdocnotation{)}} \coqexternalref{::nat scope:x '-' x}{http://rocq-prover.org/doc/V9.1.1/stdlib//Corelib.Init.Nat}{\coqdocnotation{-}} \coqexternalref{::nat scope:x '-' x}{http://rocq-prover.org/doc/V9.1.1/stdlib//Corelib.Init.Nat}{\coqdocnotation{(}}\coqref{SECF.Imp.aeval:24}{\coqdocdefinition{aeval}} \coqdocvar{a2}\coqexternalref{::nat scope:x '-' x}{http://rocq-prover.org/doc/V9.1.1/stdlib//Corelib.Init.Nat}{\coqdocnotation{)}}\coqdoceol
\coqdocindent{1.00em}
\ensuremath{|} \coqref{SECF.Imp.AExp.AMult}{\coqdocconstructor{AMult}}  \coqdocvar{a1} \coqdocvar{a2} \ensuremath{\Rightarrow} \coqexternalref{::nat scope:x '*' x}{http://rocq-prover.org/doc/V9.1.1/stdlib//Corelib.Init.Nat}{\coqdocnotation{(}}\coqref{SECF.Imp.aeval:24}{\coqdocdefinition{aeval}} \coqdocvar{a1}\coqexternalref{::nat scope:x '*' x}{http://rocq-prover.org/doc/V9.1.1/stdlib//Corelib.Init.Nat}{\coqdocnotation{)}} \coqexternalref{::nat scope:x '*' x}{http://rocq-prover.org/doc/V9.1.1/stdlib//Corelib.Init.Nat}{\coqdocnotation{\ensuremath{\times}}} \coqexternalref{::nat scope:x '*' x}{http://rocq-prover.org/doc/V9.1.1/stdlib//Corelib.Init.Nat}{\coqdocnotation{(}}\coqref{SECF.Imp.aeval:24}{\coqdocdefinition{aeval}} \coqdocvar{a2}\coqexternalref{::nat scope:x '*' x}{http://rocq-prover.org/doc/V9.1.1/stdlib//Corelib.Init.Nat}{\coqdocnotation{)}}\coqdoceol
\coqdocindent{1.00em}
\coqdockw{end}.\coqdoceol
\coqdocemptyline
\coqdocnoindent
\coqdockw{Example} \coqdef{SECF.Imp.AExp.test aeval1}{test\_aeval1}{\coqdocdefinition{test\_aeval1}}:\coqdoceol
\coqdocindent{1.00em}
\coqref{SECF.Imp.AExp.aeval}{\coqdocdefinition{aeval}} (\coqref{SECF.Imp.AExp.APlus}{\coqdocconstructor{APlus}} (\coqref{SECF.Imp.AExp.ANum}{\coqdocconstructor{ANum}} 2) (\coqref{SECF.Imp.AExp.ANum}{\coqdocconstructor{ANum}} 2)) \coqexternalref{::type scope:x '=' x}{http://rocq-prover.org/doc/V9.1.1/stdlib//Corelib.Init.Logic}{\coqdocnotation{=}} 4.\coqdoceol
\coqdocnoindent
\coqdockw{Proof}. \coqdoctac{reflexivity}. \coqdockw{Qed}.\coqdoceol
\coqdocemptyline
\end{coqdoccode}
Similarly, evaluating a boolean expression yields a boolean. 
\begin{coqdoccode}
\coqdocemptyline
\coqdocnoindent
\coqdockw{Fixpoint} \coqdef{SECF.Imp.AExp.beval}{beval}{\coqdocdefinition{beval}} (\coqdef{SECF.Imp.b:26}{b}{\coqdocbinder{b}} : \coqref{SECF.Imp.AExp.bexp}{\coqdocinductive{bexp}}) : \coqexternalref{bool}{http://rocq-prover.org/doc/V9.1.1/stdlib//Corelib.Init.Datatypes}{\coqdocinductive{bool}} :=\coqdoceol
\coqdocindent{1.00em}
\coqdockw{match} \coqref{SECF.Imp.b:26}{\coqdocvariable{b}} \coqdockw{with}\coqdoceol
\coqdocindent{1.00em}
\ensuremath{|} \coqref{SECF.Imp.AExp.BTrue}{\coqdocconstructor{BTrue}}       \ensuremath{\Rightarrow} \coqexternalref{true}{http://rocq-prover.org/doc/V9.1.1/stdlib//Corelib.Init.Datatypes}{\coqdocconstructor{true}}\coqdoceol
\coqdocindent{1.00em}
\ensuremath{|} \coqref{SECF.Imp.AExp.BFalse}{\coqdocconstructor{BFalse}}      \ensuremath{\Rightarrow} \coqexternalref{false}{http://rocq-prover.org/doc/V9.1.1/stdlib//Corelib.Init.Datatypes}{\coqdocconstructor{false}}\coqdoceol
\coqdocindent{1.00em}
\ensuremath{|} \coqref{SECF.Imp.AExp.BEq}{\coqdocconstructor{BEq}} \coqdocvar{a1} \coqdocvar{a2}   \ensuremath{\Rightarrow} \coqexternalref{::nat scope:x '=?' x}{http://rocq-prover.org/doc/V9.1.1/stdlib//Corelib.Init.Nat}{\coqdocnotation{(}}\coqref{SECF.Imp.AExp.aeval}{\coqdocdefinition{aeval}} \coqdocvar{a1}\coqexternalref{::nat scope:x '=?' x}{http://rocq-prover.org/doc/V9.1.1/stdlib//Corelib.Init.Nat}{\coqdocnotation{)}} \coqexternalref{::nat scope:x '=?' x}{http://rocq-prover.org/doc/V9.1.1/stdlib//Corelib.Init.Nat}{\coqdocnotation{=?}} \coqexternalref{::nat scope:x '=?' x}{http://rocq-prover.org/doc/V9.1.1/stdlib//Corelib.Init.Nat}{\coqdocnotation{(}}\coqref{SECF.Imp.AExp.aeval}{\coqdocdefinition{aeval}} \coqdocvar{a2}\coqexternalref{::nat scope:x '=?' x}{http://rocq-prover.org/doc/V9.1.1/stdlib//Corelib.Init.Nat}{\coqdocnotation{)}}\coqdoceol
\coqdocindent{1.00em}
\ensuremath{|} \coqref{SECF.Imp.AExp.BNeq}{\coqdocconstructor{BNeq}} \coqdocvar{a1} \coqdocvar{a2}  \ensuremath{\Rightarrow} \coqexternalref{negb}{http://rocq-prover.org/doc/V9.1.1/stdlib//Corelib.Init.Datatypes}{\coqdocdefinition{negb}} (\coqexternalref{::nat scope:x '=?' x}{http://rocq-prover.org/doc/V9.1.1/stdlib//Corelib.Init.Nat}{\coqdocnotation{(}}\coqref{SECF.Imp.AExp.aeval}{\coqdocdefinition{aeval}} \coqdocvar{a1}\coqexternalref{::nat scope:x '=?' x}{http://rocq-prover.org/doc/V9.1.1/stdlib//Corelib.Init.Nat}{\coqdocnotation{)}} \coqexternalref{::nat scope:x '=?' x}{http://rocq-prover.org/doc/V9.1.1/stdlib//Corelib.Init.Nat}{\coqdocnotation{=?}} \coqexternalref{::nat scope:x '=?' x}{http://rocq-prover.org/doc/V9.1.1/stdlib//Corelib.Init.Nat}{\coqdocnotation{(}}\coqref{SECF.Imp.AExp.aeval}{\coqdocdefinition{aeval}} \coqdocvar{a2}\coqexternalref{::nat scope:x '=?' x}{http://rocq-prover.org/doc/V9.1.1/stdlib//Corelib.Init.Nat}{\coqdocnotation{)}})\coqdoceol
\coqdocindent{1.00em}
\ensuremath{|} \coqref{SECF.Imp.AExp.BLe}{\coqdocconstructor{BLe}} \coqdocvar{a1} \coqdocvar{a2}   \ensuremath{\Rightarrow} \coqexternalref{::nat scope:x '<=?' x}{http://rocq-prover.org/doc/V9.1.1/stdlib//Corelib.Init.Nat}{\coqdocnotation{(}}\coqref{SECF.Imp.AExp.aeval}{\coqdocdefinition{aeval}} \coqdocvar{a1}\coqexternalref{::nat scope:x '<=?' x}{http://rocq-prover.org/doc/V9.1.1/stdlib//Corelib.Init.Nat}{\coqdocnotation{)}} \coqexternalref{::nat scope:x '<=?' x}{http://rocq-prover.org/doc/V9.1.1/stdlib//Corelib.Init.Nat}{\coqdocnotation{<=?}} \coqexternalref{::nat scope:x '<=?' x}{http://rocq-prover.org/doc/V9.1.1/stdlib//Corelib.Init.Nat}{\coqdocnotation{(}}\coqref{SECF.Imp.AExp.aeval}{\coqdocdefinition{aeval}} \coqdocvar{a2}\coqexternalref{::nat scope:x '<=?' x}{http://rocq-prover.org/doc/V9.1.1/stdlib//Corelib.Init.Nat}{\coqdocnotation{)}}\coqdoceol
\coqdocindent{1.00em}
\ensuremath{|} \coqref{SECF.Imp.AExp.BGt}{\coqdocconstructor{BGt}} \coqdocvar{a1} \coqdocvar{a2}   \ensuremath{\Rightarrow} \coqexternalref{negb}{http://rocq-prover.org/doc/V9.1.1/stdlib//Corelib.Init.Datatypes}{\coqdocdefinition{negb}} (\coqexternalref{::nat scope:x '<=?' x}{http://rocq-prover.org/doc/V9.1.1/stdlib//Corelib.Init.Nat}{\coqdocnotation{(}}\coqref{SECF.Imp.AExp.aeval}{\coqdocdefinition{aeval}} \coqdocvar{a1}\coqexternalref{::nat scope:x '<=?' x}{http://rocq-prover.org/doc/V9.1.1/stdlib//Corelib.Init.Nat}{\coqdocnotation{)}} \coqexternalref{::nat scope:x '<=?' x}{http://rocq-prover.org/doc/V9.1.1/stdlib//Corelib.Init.Nat}{\coqdocnotation{<=?}} \coqexternalref{::nat scope:x '<=?' x}{http://rocq-prover.org/doc/V9.1.1/stdlib//Corelib.Init.Nat}{\coqdocnotation{(}}\coqref{SECF.Imp.AExp.aeval}{\coqdocdefinition{aeval}} \coqdocvar{a2}\coqexternalref{::nat scope:x '<=?' x}{http://rocq-prover.org/doc/V9.1.1/stdlib//Corelib.Init.Nat}{\coqdocnotation{)}})\coqdoceol
\coqdocindent{1.00em}
\ensuremath{|} \coqref{SECF.Imp.AExp.BNot}{\coqdocconstructor{BNot}} \coqdocvar{b1}     \ensuremath{\Rightarrow} \coqexternalref{negb}{http://rocq-prover.org/doc/V9.1.1/stdlib//Corelib.Init.Datatypes}{\coqdocdefinition{negb}} (\coqref{SECF.Imp.beval:27}{\coqdocdefinition{beval}} \coqdocvar{b1})\coqdoceol
\coqdocindent{1.00em}
\ensuremath{|} \coqref{SECF.Imp.AExp.BAnd}{\coqdocconstructor{BAnd}} \coqdocvar{b1} \coqdocvar{b2}  \ensuremath{\Rightarrow} \coqexternalref{andb}{http://rocq-prover.org/doc/V9.1.1/stdlib//Corelib.Init.Datatypes}{\coqdocdefinition{andb}} (\coqref{SECF.Imp.beval:27}{\coqdocdefinition{beval}} \coqdocvar{b1}) (\coqref{SECF.Imp.beval:27}{\coqdocdefinition{beval}} \coqdocvar{b2})\coqdoceol
\coqdocindent{1.00em}
\coqdockw{end}.\coqdoceol
\coqdocemptyline
\end{coqdoccode}
\subsection{Optimization}



 We haven't defined very much yet, but we can already get
    some mileage out of the definitions.  Suppose we define a function
    that takes an arithmetic expression and slightly simplifies it,
    changing every occurrence of 0 + \coqdocvar{e} (i.e., (\coqref{SECF.SpecCT.APlus}{\coqdocdefinition{APlus}} (\coqref{SECF.SpecCT.ANum}{\coqdocconstructor{ANum}} 0) \coqdocvar{e})
    into just \coqdocvar{e}. 
\begin{coqdoccode}
\coqdocemptyline
\coqdocnoindent
\coqdockw{Fixpoint} \coqdef{SECF.Imp.AExp.optimize 0plus}{optimize\_0plus}{\coqdocdefinition{optimize\_0plus}} (\coqdef{SECF.Imp.a:29}{a}{\coqdocbinder{a}}:\coqref{SECF.Imp.AExp.aexp}{\coqdocinductive{aexp}}) : \coqref{SECF.Imp.AExp.aexp}{\coqdocinductive{aexp}} :=\coqdoceol
\coqdocindent{1.00em}
\coqdockw{match} \coqref{SECF.Imp.a:29}{\coqdocvariable{a}} \coqdockw{with}\coqdoceol
\coqdocindent{1.00em}
\ensuremath{|} \coqref{SECF.Imp.AExp.ANum}{\coqdocconstructor{ANum}} \coqdocvar{n} \ensuremath{\Rightarrow} \coqref{SECF.Imp.AExp.ANum}{\coqdocconstructor{ANum}} \coqdocvar{n}\coqdoceol
\coqdocindent{1.00em}
\ensuremath{|} \coqref{SECF.Imp.AExp.APlus}{\coqdocconstructor{APlus}} (\coqref{SECF.Imp.AExp.ANum}{\coqdocconstructor{ANum}} 0) \coqdocvar{e2} \ensuremath{\Rightarrow} \coqref{SECF.Imp.optimize 0plus:30}{\coqdocdefinition{optimize\_0plus}} \coqdocvar{e2}\coqdoceol
\coqdocindent{1.00em}
\ensuremath{|} \coqref{SECF.Imp.AExp.APlus}{\coqdocconstructor{APlus}}  \coqdocvar{e1} \coqdocvar{e2} \ensuremath{\Rightarrow} \coqref{SECF.Imp.AExp.APlus}{\coqdocconstructor{APlus}}  (\coqref{SECF.Imp.optimize 0plus:30}{\coqdocdefinition{optimize\_0plus}} \coqdocvar{e1}) (\coqref{SECF.Imp.optimize 0plus:30}{\coqdocdefinition{optimize\_0plus}} \coqdocvar{e2})\coqdoceol
\coqdocindent{1.00em}
\ensuremath{|} \coqref{SECF.Imp.AExp.AMinus}{\coqdocconstructor{AMinus}} \coqdocvar{e1} \coqdocvar{e2} \ensuremath{\Rightarrow} \coqref{SECF.Imp.AExp.AMinus}{\coqdocconstructor{AMinus}} (\coqref{SECF.Imp.optimize 0plus:30}{\coqdocdefinition{optimize\_0plus}} \coqdocvar{e1}) (\coqref{SECF.Imp.optimize 0plus:30}{\coqdocdefinition{optimize\_0plus}} \coqdocvar{e2})\coqdoceol
\coqdocindent{1.00em}
\ensuremath{|} \coqref{SECF.Imp.AExp.AMult}{\coqdocconstructor{AMult}}  \coqdocvar{e1} \coqdocvar{e2} \ensuremath{\Rightarrow} \coqref{SECF.Imp.AExp.AMult}{\coqdocconstructor{AMult}}  (\coqref{SECF.Imp.optimize 0plus:30}{\coqdocdefinition{optimize\_0plus}} \coqdocvar{e1}) (\coqref{SECF.Imp.optimize 0plus:30}{\coqdocdefinition{optimize\_0plus}} \coqdocvar{e2})\coqdoceol
\coqdocindent{1.00em}
\coqdockw{end}.\coqdoceol
\coqdocemptyline
\end{coqdoccode}
To gain confidence that our optimization is doing the right
    thing we can test it on some examples and see if the output looks
    OK. 
\begin{coqdoccode}
\coqdocemptyline
\coqdocnoindent
\coqdockw{Example} \coqdef{SECF.Imp.AExp.test optimize 0plus}{test\_optimize\_0plus}{\coqdocdefinition{test\_optimize\_0plus}}:\coqdoceol
\coqdocindent{1.00em}
\coqref{SECF.Imp.AExp.optimize 0plus}{\coqdocdefinition{optimize\_0plus}} (\coqref{SECF.Imp.AExp.APlus}{\coqdocconstructor{APlus}} (\coqref{SECF.Imp.AExp.ANum}{\coqdocconstructor{ANum}} 2)\coqdoceol
\coqdocindent{12.00em}
(\coqref{SECF.Imp.AExp.APlus}{\coqdocconstructor{APlus}} (\coqref{SECF.Imp.AExp.ANum}{\coqdocconstructor{ANum}} 0)\coqdoceol
\coqdocindent{15.50em}
(\coqref{SECF.Imp.AExp.APlus}{\coqdocconstructor{APlus}} (\coqref{SECF.Imp.AExp.ANum}{\coqdocconstructor{ANum}} 0) (\coqref{SECF.Imp.AExp.ANum}{\coqdocconstructor{ANum}} 1))))\coqdoceol
\coqdocindent{1.00em}
\coqexternalref{::type scope:x '=' x}{http://rocq-prover.org/doc/V9.1.1/stdlib//Corelib.Init.Logic}{\coqdocnotation{=}} \coqref{SECF.Imp.AExp.APlus}{\coqdocconstructor{APlus}} (\coqref{SECF.Imp.AExp.ANum}{\coqdocconstructor{ANum}} 2) (\coqref{SECF.Imp.AExp.ANum}{\coqdocconstructor{ANum}} 1).\coqdoceol
\coqdocnoindent
\coqdockw{Proof}. \coqdoctac{reflexivity}. \coqdockw{Qed}.\coqdoceol
\coqdocemptyline
\end{coqdoccode}
But if we want to be certain the optimization is correct --
    that evaluating an optimized expression \textit{always} gives the same
    result as the original -- we should prove it! 
\begin{coqdoccode}
\coqdocemptyline
\coqdocnoindent
\coqdockw{Theorem} \coqdef{SECF.Imp.AExp.optimize 0plus sound}{optimize\_0plus\_sound}{\coqdoclemma{optimize\_0plus\_sound}}: \coqdockw{\ensuremath{\forall}} \coqdef{SECF.Imp.a:32}{a}{\coqdocbinder{a}},\coqdoceol
\coqdocindent{1.00em}
\coqref{SECF.Imp.AExp.aeval}{\coqdocdefinition{aeval}} (\coqref{SECF.Imp.AExp.optimize 0plus}{\coqdocdefinition{optimize\_0plus}} \coqref{SECF.Imp.a:32}{\coqdocvariable{a}}) \coqexternalref{::type scope:x '=' x}{http://rocq-prover.org/doc/V9.1.1/stdlib//Corelib.Init.Logic}{\coqdocnotation{=}} \coqref{SECF.Imp.AExp.aeval}{\coqdocdefinition{aeval}} \coqref{SECF.Imp.a:32}{\coqdocvariable{a}}.\coqdoceol
\coqdocnoindent
\coqdockw{Proof}.\coqdoceol
\coqdocindent{1.00em}
\coqdoctac{intros} \coqdocvar{a}. \coqdoctac{induction} \coqdocvar{a}.\coqdoceol
\coqdocindent{1.00em}
-  \coqdoctac{reflexivity}.\coqdoceol
\coqdocindent{1.00em}
-  \coqdoctac{destruct} \coqdocvar{a1} \coqdocvar{eqn}:\coqdocvar{Ea1}.\coqdoceol
\coqdocindent{2.00em}
+  \coqdoctac{destruct} \coqdocvar{n} \coqdocvar{eqn}:\coqdocvar{En}.\coqdoceol
\coqdocindent{3.00em}
\ensuremath{\times}   \coqdoctac{simpl}. \coqdoctac{apply} \coqdocvar{IHa2}.\coqdoceol
\coqdocindent{3.00em}
\ensuremath{\times}  \coqdoctac{simpl}. \coqdoctac{rewrite} \coqdocvar{IHa2}. \coqdoctac{reflexivity}.\coqdoceol
\coqdocindent{2.00em}
+ \coqdoceol
\coqdocindent{3.00em}
\coqdoctac{simpl}. \coqdoctac{simpl} \coqdoctac{in} \coqdocvar{IHa1}. \coqdoctac{rewrite} \coqdocvar{IHa1}.\coqdoceol
\coqdocindent{3.00em}
\coqdoctac{rewrite} \coqdocvar{IHa2}. \coqdoctac{reflexivity}.\coqdoceol
\coqdocindent{2.00em}
+ \coqdoceol
\coqdocindent{3.00em}
\coqdoctac{simpl}. \coqdoctac{simpl} \coqdoctac{in} \coqdocvar{IHa1}. \coqdoctac{rewrite} \coqdocvar{IHa1}.\coqdoceol
\coqdocindent{3.00em}
\coqdoctac{rewrite} \coqdocvar{IHa2}. \coqdoctac{reflexivity}.\coqdoceol
\coqdocindent{2.00em}
+ \coqdoceol
\coqdocindent{3.00em}
\coqdoctac{simpl}. \coqdoctac{simpl} \coqdoctac{in} \coqdocvar{IHa1}. \coqdoctac{rewrite} \coqdocvar{IHa1}.\coqdoceol
\coqdocindent{3.00em}
\coqdoctac{rewrite} \coqdocvar{IHa2}. \coqdoctac{reflexivity}.\coqdoceol
\coqdocindent{1.00em}
- \coqdoceol
\coqdocindent{2.00em}
\coqdoctac{simpl}. \coqdoctac{rewrite} \coqdocvar{IHa1}. \coqdoctac{rewrite} \coqdocvar{IHa2}. \coqdoctac{reflexivity}.\coqdoceol
\coqdocindent{1.00em}
- \coqdoceol
\coqdocindent{2.00em}
\coqdoctac{simpl}. \coqdoctac{rewrite} \coqdocvar{IHa1}. \coqdoctac{rewrite} \coqdocvar{IHa2}. \coqdoctac{reflexivity}. \coqdockw{Qed}.\coqdoceol
\coqdocemptyline
\end{coqdoccode}
\section{Rocq Automation}



 The amount of repetition in this last proof is a little
    annoying.  And if either the language of arithmetic expressions or
    the optimization being proved sound were significantly more
    complex, it would start to be a real problem.


    So far, we've been doing all our proofs using just a small handful
    of Rocq's tactics and completely ignoring its powerful facilities
    for constructing parts of proofs automatically.  This section
    introduces some of these facilities, and we will see more over the
    next several chapters.  Getting used to them will take some
    energy -- Rocq's automation is a power tool -- but it will allow us
    to scale up our efforts to more complex definitions and more
    interesting properties without becoming overwhelmed by boring,
    repetitive, low-level details. 
\begin{coqdoccode}
\coqdocemptyline
\end{coqdoccode}
\subsection{Tacticals}



 \textit{Tacticals} is Rocq's term for tactics that take other tactics as
    arguments -- ``higher-order tactics,'' if you will.  
\begin{coqdoccode}
\coqdocemptyline
\end{coqdoccode}
\subsubsection{The \texorpdfstring{\protect\coqdoctac{try}}{try} Tactical}



 If \coqref{SECF.Hoare2.T}{\coqdocdefinition{T}} is a tactic, then \coqdoctac{try} \coqref{SECF.Hoare2.T}{\coqdocdefinition{T}} is a tactic that is just like \coqref{SECF.Hoare2.T}{\coqdocdefinition{T}}
    except that, if \coqref{SECF.Hoare2.T}{\coqdocdefinition{T}} fails, \coqdoctac{try} \coqref{SECF.Hoare2.T}{\coqdocdefinition{T}} \textit{successfully} does nothing at
    all (rather than failing). 
\begin{coqdoccode}
\coqdocnoindent
\coqdockw{Theorem} \coqdef{SECF.Imp.AExp.silly1}{silly1}{\coqdoclemma{silly1}} : \coqdockw{\ensuremath{\forall}} (\coqdef{SECF.Imp.P:33}{P}{\coqdocbinder{P}} : \coqdockw{Prop}), \coqref{SECF.Imp.P:33}{\coqdocvariable{P}} \coqexternalref{::type scope:x '->' x}{http://rocq-prover.org/doc/V9.1.1/stdlib//Corelib.Init.Logic}{\coqdocnotation{\ensuremath{\rightarrow}}} \coqref{SECF.Imp.P:33}{\coqdocvariable{P}}.\coqdoceol
\coqdocnoindent
\coqdockw{Proof}.\coqdoceol
\coqdocindent{1.00em}
\coqdoctac{intros} \coqdocvar{P} \coqdocvar{HP}.\coqdoceol
\coqdocindent{1.00em}
\coqdoctac{try} \coqdoctac{reflexivity}. \coqdocindent{1.00em}
\coqdoctac{apply} \coqdocvar{HP}. \coqdocnoindent
\coqdockw{Qed}.\coqdoceol
\coqdocemptyline
\coqdocnoindent
\coqdockw{Theorem} \coqdef{SECF.Imp.AExp.silly2}{silly2}{\coqdoclemma{silly2}} : \coqdockw{\ensuremath{\forall}} \coqdef{SECF.Imp.ae:34}{ae}{\coqdocbinder{ae}}, \coqref{SECF.Imp.AExp.aeval}{\coqdocdefinition{aeval}} \coqref{SECF.Imp.ae:34}{\coqdocvariable{ae}} \coqexternalref{::type scope:x '=' x}{http://rocq-prover.org/doc/V9.1.1/stdlib//Corelib.Init.Logic}{\coqdocnotation{=}} \coqref{SECF.Imp.AExp.aeval}{\coqdocdefinition{aeval}} \coqref{SECF.Imp.ae:34}{\coqdocvariable{ae}}.\coqdoceol
\coqdocnoindent
\coqdockw{Proof}.\coqdoceol
\coqdocindent{2.00em}
\coqdoctac{try} \coqdoctac{reflexivity}. \coqdocnoindent
\coqdockw{Qed}.\coqdoceol
\coqdocemptyline
\end{coqdoccode}
There is not much reason to use \coqdoctac{try} in completely manual
    proofs like these, but it is very useful for doing automated
    proofs in conjunction with the ; tactical, which we show
    next. 
\begin{coqdoccode}
\coqdocemptyline
\end{coqdoccode}
\subsubsection{The ; Tactical (Simple Form)}



 In its most common form, the ; tactical takes two tactics as
    arguments.  The compound tactic \coqref{SECF.Hoare2.T}{\coqdocdefinition{T}};\coqdocvar{T'} first performs \coqref{SECF.Hoare2.T}{\coqdocdefinition{T}} and then
    performs \coqdocvar{T'} on \textit{each subgoal} generated by \coqref{SECF.Hoare2.T}{\coqdocdefinition{T}}. 

 For example, consider the following trivial lemma: 
\begin{coqdoccode}
\coqdocemptyline
\coqdocnoindent
\coqdockw{Lemma} \coqdef{SECF.Imp.AExp.foo}{foo}{\coqdoclemma{foo}} : \coqdockw{\ensuremath{\forall}} \coqdef{SECF.Imp.n:35}{n}{\coqdocbinder{n}}, 0 \coqexternalref{::nat scope:x '<=?' x}{http://rocq-prover.org/doc/V9.1.1/stdlib//Corelib.Init.Nat}{\coqdocnotation{<=?}} \coqref{SECF.Imp.n:35}{\coqdocvariable{n}} \coqexternalref{::type scope:x '=' x}{http://rocq-prover.org/doc/V9.1.1/stdlib//Corelib.Init.Logic}{\coqdocnotation{=}} \coqexternalref{true}{http://rocq-prover.org/doc/V9.1.1/stdlib//Corelib.Init.Datatypes}{\coqdocconstructor{true}}.\coqdoceol
\coqdocnoindent
\coqdockw{Proof}.\coqdoceol
\coqdocindent{1.00em}
\coqdoctac{intros}.\coqdoceol
\coqdocindent{1.00em}
\coqdoctac{destruct} \coqdocvar{n}.\coqdoceol
\coqdocindent{2.00em}
-  \coqdoctac{simpl}. \coqdoctac{reflexivity}.\coqdoceol
\coqdocindent{2.00em}
-  \coqdoctac{simpl}. \coqdoctac{reflexivity}.\coqdoceol
\coqdocnoindent
\coqdockw{Qed}.\coqdoceol
\coqdocemptyline
\end{coqdoccode}
We can simplify this proof using the ; tactical: 
\begin{coqdoccode}
\coqdocemptyline
\coqdocnoindent
\coqdockw{Lemma} \coqdef{SECF.Imp.AExp.foo'}{foo'}{\coqdoclemma{foo'}} : \coqdockw{\ensuremath{\forall}} \coqdef{SECF.Imp.n:36}{n}{\coqdocbinder{n}}, 0 \coqexternalref{::nat scope:x '<=?' x}{http://rocq-prover.org/doc/V9.1.1/stdlib//Corelib.Init.Nat}{\coqdocnotation{<=?}} \coqref{SECF.Imp.n:36}{\coqdocvariable{n}} \coqexternalref{::type scope:x '=' x}{http://rocq-prover.org/doc/V9.1.1/stdlib//Corelib.Init.Logic}{\coqdocnotation{=}} \coqexternalref{true}{http://rocq-prover.org/doc/V9.1.1/stdlib//Corelib.Init.Datatypes}{\coqdocconstructor{true}}.\coqdoceol
\coqdocnoindent
\coqdockw{Proof}.\coqdoceol
\coqdocindent{1.00em}
\coqdoctac{intros}.\coqdoceol
\coqdocindent{1.00em}
\coqdoctac{destruct} \coqdocvar{n};\coqdoceol
\coqdocindent{1.00em}
\coqdoceol
\coqdocindent{1.00em}
\coqdoctac{simpl};\coqdoceol
\coqdocindent{1.00em}
\coqdoceol
\coqdocindent{1.00em}
\coqdoctac{reflexivity}.\coqdoceol
\coqdocnoindent
\coqdockw{Qed}.\coqdoceol
\coqdocemptyline
\end{coqdoccode}
Using \coqdoctac{try} and ; together, we can get rid of the repetition in
    the proof that was bothering us a little while ago. 
\begin{coqdoccode}
\coqdocemptyline
\coqdocnoindent
\coqdockw{Theorem} \coqdef{SECF.Imp.AExp.optimize 0plus sound'}{optimize\_0plus\_sound'}{\coqdoclemma{optimize\_0plus\_sound'}}: \coqdockw{\ensuremath{\forall}} \coqdef{SECF.Imp.a:37}{a}{\coqdocbinder{a}},\coqdoceol
\coqdocindent{1.00em}
\coqref{SECF.Imp.AExp.aeval}{\coqdocdefinition{aeval}} (\coqref{SECF.Imp.AExp.optimize 0plus}{\coqdocdefinition{optimize\_0plus}} \coqref{SECF.Imp.a:37}{\coqdocvariable{a}}) \coqexternalref{::type scope:x '=' x}{http://rocq-prover.org/doc/V9.1.1/stdlib//Corelib.Init.Logic}{\coqdocnotation{=}} \coqref{SECF.Imp.AExp.aeval}{\coqdocdefinition{aeval}} \coqref{SECF.Imp.a:37}{\coqdocvariable{a}}.\coqdoceol
\coqdocnoindent
\coqdockw{Proof}.\coqdoceol
\coqdocindent{1.00em}
\coqdoctac{intros} \coqdocvar{a}.\coqdoceol
\coqdocindent{1.00em}
\coqdoctac{induction} \coqdocvar{a};\coqdoceol
\coqdocindent{2.00em}
\coqdoceol
\coqdocindent{2.00em}
\coqdoctac{try} (\coqdoctac{simpl}; \coqdoctac{rewrite} \coqdocvar{IHa1}; \coqdoctac{rewrite} \coqdocvar{IHa2}; \coqdoctac{reflexivity}).\coqdoceol
\coqdocindent{1.00em}
-  \coqdoctac{reflexivity}.\coqdoceol
\coqdocindent{1.00em}
- \coqdoceol
\coqdocindent{2.00em}
\coqdoctac{destruct} \coqdocvar{a1} \coqdocvar{eqn}:\coqdocvar{Ea1};\coqdoceol
\coqdocindent{3.00em}
\coqdoceol
\coqdocindent{3.00em}
\coqdoctac{try} (\coqdoctac{simpl}; \coqdoctac{simpl} \coqdoctac{in} \coqdocvar{IHa1}; \coqdoctac{rewrite} \coqdocvar{IHa1};\coqdoceol
\coqdocindent{5.50em}
\coqdoctac{rewrite} \coqdocvar{IHa2}; \coqdoctac{reflexivity}).\coqdoceol
\coqdocindent{2.00em}
+  \coqdoctac{destruct} \coqdocvar{n} \coqdocvar{eqn}:\coqdocvar{En};\coqdoceol
\coqdocindent{3.00em}
\coqdoctac{simpl}; \coqdoctac{rewrite} \coqdocvar{IHa2}; \coqdoctac{reflexivity}. \coqdockw{Qed}.\coqdoceol
\coqdocemptyline
\end{coqdoccode}
Rocq experts often use this ``...; \coqdoctac{try}... '' idiom after a tactic
    like \coqdoctac{induction} to take care of many similar cases all at once.
    Indeed, this practice has an analog in informal proofs.  For
    example, here is an informal proof of the optimization theorem
    that matches the structure of the formal one:


    \textit{Theorem}: For all arithmetic expressions \coqdocvar{a},


       aeval (optimize\_0plus a) = aeval a.


    \textit{Proof}: By induction on \coqdocvar{a}.  Most cases follow directly from the
    IH.  The remaining cases are as follows:



\begin{itemize}
\item  Suppose \coqdocvar{a} = \coqref{SECF.SpecCT.ANum}{\coqdocconstructor{ANum}} \coqdocvar{n} for some \coqdocvar{n}.  We must show


          aeval (optimize\_0plus (ANum n)) = aeval (ANum n).


        This is immediate from the definition of \coqref{SECF.Imp.AExp.optimize 0plus}{\coqdocdefinition{optimize\_0plus}}.



\item  Suppose \coqdocvar{a} = \coqref{SECF.SpecCT.APlus}{\coqdocdefinition{APlus}} \coqdocvar{a1} \coqdocvar{a2} for some \coqdocvar{a1} and \coqdocvar{a2}.  We must
        show


          aeval (optimize\_0plus (APlus a1 a2)) = aeval (APlus a1 a2).


        Consider the possible forms of \coqdocvar{a1}.  For most of them,
        \coqref{SECF.Imp.AExp.optimize 0plus}{\coqdocdefinition{optimize\_0plus}} simply calls itself recursively for the
        subexpressions and rebuilds a new expression of the same form
        as \coqdocvar{a1}; in these cases, the result follows directly from the
        IH.


        The interesting case is when \coqdocvar{a1} = \coqref{SECF.SpecCT.ANum}{\coqdocconstructor{ANum}} \coqdocvar{n} for some \coqdocvar{n}.  If
        \coqdocvar{n} = 0, then


          optimize\_0plus (APlus a1 a2) = optimize\_0plus a2


        and the IH for \coqdocvar{a2} is exactly what we need.  On the other
        hand, if \coqdocvar{n} = \coqexternalref{S}{http://rocq-prover.org/doc/V9.1.1/stdlib//Corelib.Init.Datatypes}{\coqdocconstructor{S}} \coqdocvar{n'} for some \coqdocvar{n'}, then again \coqref{SECF.Imp.AExp.optimize 0plus}{\coqdocdefinition{optimize\_0plus}}
        simply calls itself recursively, and the result follows from
        the IH.  \ensuremath{\Box} 
\end{itemize}


 However, this proof can still be improved: the first case (for
    \coqdocvar{a} = \coqref{SECF.SpecCT.ANum}{\coqdocconstructor{ANum}} \coqdocvar{n}) is very trivial -- even more trivial than the cases
    that we said simply followed from the IH -- yet we have chosen to
    write it out in full.  It would be better and clearer to drop it
    and just say, at the top, ``Most cases are either immediate or
    direct from the IH.  The only interesting case is the one for
    \coqref{SECF.SpecCT.APlus}{\coqdocdefinition{APlus}}...''  We can make the same improvement in our formal proof
    too.  Here's how it looks: 
\begin{coqdoccode}
\coqdocemptyline
\coqdocnoindent
\coqdockw{Theorem} \coqdef{SECF.Imp.AExp.optimize 0plus sound''}{optimize\_0plus\_sound'{}'}{\coqdoclemma{optimize\_0plus\_sound'{}'}}: \coqdockw{\ensuremath{\forall}} \coqdef{SECF.Imp.a:38}{a}{\coqdocbinder{a}},\coqdoceol
\coqdocindent{1.00em}
\coqref{SECF.Imp.AExp.aeval}{\coqdocdefinition{aeval}} (\coqref{SECF.Imp.AExp.optimize 0plus}{\coqdocdefinition{optimize\_0plus}} \coqref{SECF.Imp.a:38}{\coqdocvariable{a}}) \coqexternalref{::type scope:x '=' x}{http://rocq-prover.org/doc/V9.1.1/stdlib//Corelib.Init.Logic}{\coqdocnotation{=}} \coqref{SECF.Imp.AExp.aeval}{\coqdocdefinition{aeval}} \coqref{SECF.Imp.a:38}{\coqdocvariable{a}}.\coqdoceol
\coqdocnoindent
\coqdockw{Proof}.\coqdoceol
\coqdocindent{1.00em}
\coqdoctac{intros} \coqdocvar{a}.\coqdoceol
\coqdocindent{1.00em}
\coqdoctac{induction} \coqdocvar{a};\coqdoceol
\coqdocindent{2.00em}
\coqdoceol
\coqdocindent{2.00em}
\coqdoctac{try} (\coqdoctac{simpl}; \coqdoctac{rewrite} \coqdocvar{IHa1}; \coqdoctac{rewrite} \coqdocvar{IHa2}; \coqdoctac{reflexivity});\coqdoceol
\coqdocindent{2.00em}
\coqdoceol
\coqdocindent{2.00em}
\coqdoctac{try} \coqdoctac{reflexivity}.\coqdoceol
\coqdocindent{1.00em}
- \coqdoceol
\coqdocindent{2.00em}
\coqdoctac{destruct} \coqdocvar{a1}; \coqdoctac{try} (\coqdoctac{simpl}; \coqdoctac{simpl} \coqdoctac{in} \coqdocvar{IHa1}; \coqdoctac{rewrite} \coqdocvar{IHa1};\coqdoceol
\coqdocindent{11.00em}
\coqdoctac{rewrite} \coqdocvar{IHa2}; \coqdoctac{reflexivity}).\coqdoceol
\coqdocindent{2.00em}
+  \coqdoctac{destruct} \coqdocvar{n};\coqdoceol
\coqdocindent{3.00em}
\coqdoctac{simpl}; \coqdoctac{rewrite} \coqdocvar{IHa2}; \coqdoctac{reflexivity}. \coqdockw{Qed}.\coqdoceol
\coqdocemptyline
\end{coqdoccode}
\subsubsection{The ; Tactical (General Form)}



 The ; tactical also has a more general form than the simple
    \coqref{SECF.Hoare2.T}{\coqdocdefinition{T}};\coqdocvar{T'} we've seen above.  If \coqref{SECF.Hoare2.T}{\coqdocdefinition{T}}, \coqdocvar{T1}, ..., \coqdocvar{Tn} are tactics,
    then


      T; \coqdocvar{T1} \ensuremath{|} \coqdocvar{T2} \ensuremath{|} ... \ensuremath{|} \coqdocvar{Tn}


    is a tactic that first performs \coqref{SECF.Hoare2.T}{\coqdocdefinition{T}} and then performs \coqdocvar{T1} on the
    first subgoal generated by \coqref{SECF.Hoare2.T}{\coqdocdefinition{T}}, performs \coqdocvar{T2} on the second
    subgoal, etc.


    So \coqref{SECF.Hoare2.T}{\coqdocdefinition{T}};\coqdocvar{T'} is just special notation for the case when all of the
    \coqdocvar{Ti}'s are the same tactic; i.e., \coqref{SECF.Hoare2.T}{\coqdocdefinition{T}};\coqdocvar{T'} is shorthand for:


      T; \coqdocvar{T'} \ensuremath{|} \coqdocvar{T'} \ensuremath{|} ... \ensuremath{|} \coqdocvar{T'}

\begin{coqdoccode}
\coqdocemptyline
\end{coqdoccode}
\subsubsection{The \texorpdfstring{\protect\coqdoctac{repeat}}{repeat} Tactical}



 The \coqdoctac{repeat} tactical takes another tactic and keeps applying this
    tactic until it fails or until it succeeds but doesn't make any
    progress.


    Here is an example proving that 10 is in a long list using
    \coqdoctac{repeat}. 
\begin{coqdoccode}
\coqdocemptyline
\coqdocnoindent
\coqdockw{Theorem} \coqdef{SECF.Imp.AExp.In10}{In10}{\coqdoclemma{In10}} : \coqdocdefinition{In} 10 \coqdocnotation{[}1\coqdocnotation{;}2\coqdocnotation{;}3\coqdocnotation{;}4\coqdocnotation{;}5\coqdocnotation{;}6\coqdocnotation{;}7\coqdocnotation{;}8\coqdocnotation{;}9\coqdocnotation{;}10\coqdocnotation{]}.\coqdoceol
\coqdocnoindent
\coqdockw{Proof}.\coqdoceol
\coqdocindent{1.00em}
\coqdoctac{repeat} (\coqdoctac{try} (\coqdoctac{left}; \coqdoctac{reflexivity}); \coqdoctac{right}).\coqdoceol
\coqdocnoindent
\coqdockw{Qed}.\coqdoceol
\coqdocemptyline
\end{coqdoccode}
The tactic \coqdoctac{repeat} \coqref{SECF.Hoare2.T}{\coqdocdefinition{T}} never fails: if the tactic \coqref{SECF.Hoare2.T}{\coqdocdefinition{T}} doesn't apply
    to the original goal, then repeat \textit{succeeds} without changing the
    goal at all (i.e., it repeats zero times). 
\begin{coqdoccode}
\coqdocemptyline
\coqdocnoindent
\coqdockw{Theorem} \coqdef{SECF.Imp.AExp.In10'}{In10'}{\coqdoclemma{In10'}} : \coqdocdefinition{In} 10 \coqdocnotation{[}1\coqdocnotation{;}2\coqdocnotation{;}3\coqdocnotation{;}4\coqdocnotation{;}5\coqdocnotation{;}6\coqdocnotation{;}7\coqdocnotation{;}8\coqdocnotation{;}9\coqdocnotation{;}10\coqdocnotation{]}.\coqdoceol
\coqdocnoindent
\coqdockw{Proof}.\coqdoceol
\coqdocindent{1.00em}
\coqdoctac{repeat} \coqdoctac{simpl}.\coqdoceol
\coqdocindent{1.00em}
\coqdoctac{repeat} (\coqdoctac{left}; \coqdoctac{reflexivity}).\coqdoceol
\coqdocindent{1.00em}
\coqdoctac{repeat} (\coqdoctac{right}; \coqdoctac{try} (\coqdoctac{left}; \coqdoctac{reflexivity})).\coqdoceol
\coqdocnoindent
\coqdockw{Qed}.\coqdoceol
\coqdocemptyline
\end{coqdoccode}
The tactic \coqdoctac{repeat} \coqref{SECF.Hoare2.T}{\coqdocdefinition{T}} does not have any upper bound on the
    number of times it applies \coqref{SECF.Hoare2.T}{\coqdocdefinition{T}}.  If \coqref{SECF.Hoare2.T}{\coqdocdefinition{T}} is a tactic that \textit{always}
    succeeds (and makes progress), then repeat \coqref{SECF.Hoare2.T}{\coqdocdefinition{T}} will loop
    forever. 
\begin{coqdoccode}
\coqdocemptyline
\coqdocnoindent
\coqdockw{Theorem} \coqdef{SECF.Imp.AExp.repeat loop}{repeat\_loop}{\coqdoclemma{repeat\_loop}} : \coqdockw{\ensuremath{\forall}} (\coqdef{SECF.Imp.m:39}{m}{\coqdocbinder{m}} \coqdef{SECF.Imp.n:40}{n}{\coqdocbinder{n}} : \coqexternalref{nat}{http://rocq-prover.org/doc/V9.1.1/stdlib//Corelib.Init.Datatypes}{\coqdocinductive{nat}}),\coqdoceol
\coqdocindent{1.00em}
\coqref{SECF.Imp.m:39}{\coqdocvariable{m}} \coqexternalref{::nat scope:x '+' x}{http://rocq-prover.org/doc/V9.1.1/stdlib//Corelib.Init.Nat}{\coqdocnotation{+}} \coqref{SECF.Imp.n:40}{\coqdocvariable{n}} \coqexternalref{::type scope:x '=' x}{http://rocq-prover.org/doc/V9.1.1/stdlib//Corelib.Init.Logic}{\coqdocnotation{=}} \coqref{SECF.Imp.n:40}{\coqdocvariable{n}} \coqexternalref{::nat scope:x '+' x}{http://rocq-prover.org/doc/V9.1.1/stdlib//Corelib.Init.Nat}{\coqdocnotation{+}} \coqref{SECF.Imp.m:39}{\coqdocvariable{m}}.\coqdoceol
\coqdocnoindent
\coqdockw{Proof}.\coqdoceol
\coqdocindent{1.00em}
\coqdoctac{intros} \coqdocvar{m} \coqdocvar{n}.\coqdoceol
\coqdocnoindent
\coqdocvar{Admitted}.\coqdoceol
\coqdocemptyline
\end{coqdoccode}
Wait -- did we just write an infinite loop in Rocq?!?!


    Sort of.


    While evaluation in Rocq's term language, Gallina, is guaranteed to
    terminate, \textit{tactic} evaluation is not.  This does not affect Rocq's
    logical consistency, however, since the job of \coqdoctac{repeat} and other
    tactics is to guide Rocq in constructing proofs; if the
    construction process diverges (i.e., it does not terminate), this
    simply means that we have failed to construct a proof at all, not
    that we have constructed a bad proof. 

\paragraph{Exercise: 3 stars, standard (optimize\_0plus\_b\_sound)}




    Since the \coqref{SECF.Imp.AExp.optimize 0plus}{\coqdocdefinition{optimize\_0plus}} transformation doesn't change the value
    of \coqref{SECF.Imp.aexp}{\coqdocinductive{aexp}}s, we should be able to apply it to all the \coqref{SECF.Imp.aexp}{\coqdocinductive{aexp}}s that
    appear in a \coqref{SECF.Imp.bexp}{\coqdocinductive{bexp}} without changing the \coqref{SECF.Imp.bexp}{\coqdocinductive{bexp}}'s value.  Write a
    function that performs this transformation on \coqref{SECF.Imp.bexp}{\coqdocinductive{bexp}}s and prove
    it is sound.  Use the tacticals we've just seen to make the proof
    as short and elegant as possible. 
\begin{coqdoccode}
\coqdocemptyline
\coqdocnoindent
\coqdockw{Fixpoint} \coqdef{SECF.Imp.AExp.optimize 0plus b}{optimize\_0plus\_b}{\coqdocdefinition{optimize\_0plus\_b}} (\coqdef{SECF.Imp.b:41}{b}{\coqdocbinder{b}} : \coqref{SECF.Imp.AExp.bexp}{\coqdocinductive{bexp}}) : \coqref{SECF.Imp.AExp.bexp}{\coqdocinductive{bexp}}\coqdoceol
\coqdocindent{1.00em}
. \coqdocvar{Admitted}.\coqdoceol
\coqdocemptyline
\coqdocnoindent
\coqdockw{Example} \coqdef{SECF.Imp.AExp.optimize 0plus b test1}{optimize\_0plus\_b\_test1}{\coqdocdefinition{optimize\_0plus\_b\_test1}}:\coqdoceol
\coqdocindent{1.00em}
\coqref{SECF.Imp.AExp.optimize 0plus b}{\coqdocaxiom{optimize\_0plus\_b}} (\coqref{SECF.Imp.AExp.BNot}{\coqdocconstructor{BNot}} (\coqref{SECF.Imp.AExp.BGt}{\coqdocconstructor{BGt}} (\coqref{SECF.Imp.AExp.APlus}{\coqdocconstructor{APlus}} (\coqref{SECF.Imp.AExp.ANum}{\coqdocconstructor{ANum}} 0) (\coqref{SECF.Imp.AExp.ANum}{\coqdocconstructor{ANum}} 4)) (\coqref{SECF.Imp.AExp.ANum}{\coqdocconstructor{ANum}} 8))) \coqexternalref{::type scope:x '=' x}{http://rocq-prover.org/doc/V9.1.1/stdlib//Corelib.Init.Logic}{\coqdocnotation{=}}\coqdoceol
\coqdocindent{9.50em}
\coqexternalref{::type scope:x '=' x}{http://rocq-prover.org/doc/V9.1.1/stdlib//Corelib.Init.Logic}{\coqdocnotation{(}}\coqref{SECF.Imp.AExp.BNot}{\coqdocconstructor{BNot}} (\coqref{SECF.Imp.AExp.BGt}{\coqdocconstructor{BGt}} (\coqref{SECF.Imp.AExp.ANum}{\coqdocconstructor{ANum}} 4) (\coqref{SECF.Imp.AExp.ANum}{\coqdocconstructor{ANum}} 8))\coqexternalref{::type scope:x '=' x}{http://rocq-prover.org/doc/V9.1.1/stdlib//Corelib.Init.Logic}{\coqdocnotation{)}}.\coqdoceol
\coqdocnoindent
\coqdockw{Proof}.  \coqdocvar{Admitted}.\coqdoceol
\coqdocemptyline
\coqdocnoindent
\coqdockw{Example} \coqdef{SECF.Imp.AExp.optimize 0plus b test2}{optimize\_0plus\_b\_test2}{\coqdocdefinition{optimize\_0plus\_b\_test2}}:\coqdoceol
\coqdocindent{1.00em}
\coqref{SECF.Imp.AExp.optimize 0plus b}{\coqdocaxiom{optimize\_0plus\_b}} (\coqref{SECF.Imp.AExp.BAnd}{\coqdocconstructor{BAnd}} (\coqref{SECF.Imp.AExp.BLe}{\coqdocconstructor{BLe}} (\coqref{SECF.Imp.AExp.APlus}{\coqdocconstructor{APlus}} (\coqref{SECF.Imp.AExp.ANum}{\coqdocconstructor{ANum}} 0) (\coqref{SECF.Imp.AExp.ANum}{\coqdocconstructor{ANum}} 4)) (\coqref{SECF.Imp.AExp.ANum}{\coqdocconstructor{ANum}} 5)) \coqref{SECF.Imp.AExp.BTrue}{\coqdocconstructor{BTrue}}) \coqexternalref{::type scope:x '=' x}{http://rocq-prover.org/doc/V9.1.1/stdlib//Corelib.Init.Logic}{\coqdocnotation{=}}\coqdoceol
\coqdocindent{9.50em}
\coqexternalref{::type scope:x '=' x}{http://rocq-prover.org/doc/V9.1.1/stdlib//Corelib.Init.Logic}{\coqdocnotation{(}}\coqref{SECF.Imp.AExp.BAnd}{\coqdocconstructor{BAnd}} (\coqref{SECF.Imp.AExp.BLe}{\coqdocconstructor{BLe}} (\coqref{SECF.Imp.AExp.ANum}{\coqdocconstructor{ANum}} 4) (\coqref{SECF.Imp.AExp.ANum}{\coqdocconstructor{ANum}} 5)) \coqref{SECF.Imp.AExp.BTrue}{\coqdocconstructor{BTrue}}\coqexternalref{::type scope:x '=' x}{http://rocq-prover.org/doc/V9.1.1/stdlib//Corelib.Init.Logic}{\coqdocnotation{)}}.\coqdoceol
\coqdocnoindent
\coqdockw{Proof}.  \coqdocvar{Admitted}.\coqdoceol
\coqdocemptyline
\coqdocnoindent
\coqdockw{Theorem} \coqdef{SECF.Imp.AExp.optimize 0plus b sound}{optimize\_0plus\_b\_sound}{\coqdoclemma{optimize\_0plus\_b\_sound}} : \coqdockw{\ensuremath{\forall}} \coqdef{SECF.Imp.b:43}{b}{\coqdocbinder{b}},\coqdoceol
\coqdocindent{1.00em}
\coqref{SECF.Imp.AExp.beval}{\coqdocdefinition{beval}} (\coqref{SECF.Imp.AExp.optimize 0plus b}{\coqdocaxiom{optimize\_0plus\_b}} \coqref{SECF.Imp.b:43}{\coqdocvariable{b}}) \coqexternalref{::type scope:x '=' x}{http://rocq-prover.org/doc/V9.1.1/stdlib//Corelib.Init.Logic}{\coqdocnotation{=}} \coqref{SECF.Imp.AExp.beval}{\coqdocdefinition{beval}} \coqref{SECF.Imp.b:43}{\coqdocvariable{b}}.\coqdoceol
\coqdocnoindent
\coqdockw{Proof}.\coqdoceol
 \coqdocvar{Admitted}.\coqdoceol
\end{coqdoccode}
\ensuremath{\Box} 

\paragraph{Exercise: 4 stars, standard, optional (optimize)}




    \textit{Design exercise}: The optimization implemented by our
    \coqref{SECF.Imp.AExp.optimize 0plus}{\coqdocdefinition{optimize\_0plus}} function is only one of many possible
    optimizations on arithmetic and boolean expressions.  Write a more
    sophisticated optimizer and prove it correct.  (You will probably
    find it easiest to start small -- add just a single, simple
    optimization and its correctness proof -- and build up
    incrementally to something more interesting.)  
\begin{coqdoccode}
\coqdocemptyline
\coqdocemptyline
\end{coqdoccode}
\subsection{Defining New Tactics}



 Rocq also provides facilities for ``programming'' in tactic
    scripts.


    The \coqdockw{Ltac} idiom illustrated below gives a handy way to define
    ``shorthand tactics'' that bundle several tactics into a single
    command.


    Ltac also includes syntactic pattern-matching on the goal and
    context, as well as general programming facilities.


    It is useful for proof automation and there are several idioms for
    programming with Ltac. Because it is a language style you might
    not have seen before, a good reference is the textbook ``Certified
    Programming with dependent types'' \coqdocvar{CPDT}, which is more advanced
    that what we will need in this course, but is considered by many
    the best reference for Ltac programming.


    Just for future reference: Rocq provides two other ways of defining
    new tactics.  There is a \coqdockw{Tactic} \coqdockw{Notation} command that allows
    defining new tactics with custom control over their concrete
    syntax. And there is also a low-level API that can be used to
    build tactics that directly manipulate Rocq's internal structures.
    We will not need either of these for present purposes.


    Here's an example \coqdockw{Ltac} script called \coqdocvar{invert}. 
\begin{coqdoccode}
\coqdocemptyline
\coqdocnoindent
\coqdockw{Ltac} \coqdocvar{invert} \coqdocvar{H} :=\coqdoceol
\coqdocindent{1.00em}
\coqdoctac{inversion} \coqdocvar{H}; \coqdoctac{subst}; \coqdoctac{clear} \coqdocvar{H}.\coqdoceol
\coqdocemptyline
\end{coqdoccode}
This defines a new tactic called \coqdocvar{invert} that takes a hypothesis
    \coqdocvar{H} as an argument and performs the sequence of commands
    \coqdoctac{inversion} \coqdocvar{H}; \coqdoctac{subst}; \coqdoctac{clear} \coqdocvar{H}. This gives us quick way to do
    inversion on evidence and constructors, rewrite with the generated
    equations, and remove the redundant hypothesis at the end. 
\begin{coqdoccode}
\coqdocemptyline
\coqdocnoindent
\coqdockw{Lemma} \coqdef{SECF.Imp.AExp.invert example1}{invert\_example1}{\coqdoclemma{invert\_example1}}: \coqdockw{\ensuremath{\forall}} \{\coqdef{SECF.Imp.a:44}{a}{\coqdocbinder{a}} \coqdef{SECF.Imp.b:45}{b}{\coqdocbinder{b}} \coqdef{SECF.Imp.c:46}{c}{\coqdocbinder{c}}: \coqexternalref{nat}{http://rocq-prover.org/doc/V9.1.1/stdlib//Corelib.Init.Datatypes}{\coqdocinductive{nat}}\}, \coqdocnotation{[}\coqref{SECF.Imp.a:44}{\coqdocvariable{a}} \coqdocnotation{;}\coqref{SECF.Imp.b:45}{\coqdocvariable{b}}\coqdocnotation{]} \coqexternalref{::type scope:x '=' x}{http://rocq-prover.org/doc/V9.1.1/stdlib//Corelib.Init.Logic}{\coqdocnotation{=}} \coqdocnotation{[}\coqref{SECF.Imp.a:44}{\coqdocvariable{a}}\coqdocnotation{;}\coqref{SECF.Imp.c:46}{\coqdocvariable{c}}\coqdocnotation{]} \coqexternalref{::type scope:x '->' x}{http://rocq-prover.org/doc/V9.1.1/stdlib//Corelib.Init.Logic}{\coqdocnotation{\ensuremath{\rightarrow}}} \coqref{SECF.Imp.b:45}{\coqdocvariable{b}} \coqexternalref{::type scope:x '=' x}{http://rocq-prover.org/doc/V9.1.1/stdlib//Corelib.Init.Logic}{\coqdocnotation{=}} \coqref{SECF.Imp.c:46}{\coqdocvariable{c}}.\coqdoceol
\coqdocindent{1.00em}
\coqdoctac{intros}.\coqdoceol
\coqdocindent{1.00em}
\coqdocvar{invert} \coqdocvar{H}.\coqdoceol
\coqdocindent{1.00em}
\coqdoctac{reflexivity}.\coqdoceol
\coqdocnoindent
\coqdockw{Qed}.\coqdoceol
\coqdocemptyline
\end{coqdoccode}
\subsection{The \texorpdfstring{\protect}{lia}\coqdocvar{lia} Tactic}



 The \coqdocvar{lia} tactic implements a decision procedure for integer linear
    arithmetic, a subset of propositional logic and arithmetic.


    If the goal is a universally quantified formula made out of



\begin{itemize}
\item  numeric constants, addition (+ and \coqexternalref{S}{http://rocq-prover.org/doc/V9.1.1/stdlib//Corelib.Init.Datatypes}{\coqdocconstructor{S}}), subtraction (-
        and \coqdocdefinition{pred}), and multiplication by constants (this is what
        makes it Presburger arithmetic),



\item  equality (= and \ensuremath{\not=}) and ordering (\ensuremath{\le} and >), and



\item  the logical connectives \ensuremath{\land}, \ensuremath{\lor}, \ensuremath{\lnot}, and \ensuremath{\rightarrow},

\end{itemize}


    then invoking \coqdocvar{lia} will either solve the goal or fail, meaning
    that the goal is actually false.  (If the goal is \textit{not} of this
    form, \coqdocvar{lia} will fail.) 
\begin{coqdoccode}
\coqdocemptyline
\coqdocnoindent
\coqdockw{Example} \coqdef{SECF.Imp.AExp.silly presburger example}{silly\_presburger\_example}{\coqdocdefinition{silly\_presburger\_example}} : \coqdockw{\ensuremath{\forall}} \coqdef{SECF.Imp.m:47}{m}{\coqdocbinder{m}} \coqdef{SECF.Imp.n:48}{n}{\coqdocbinder{n}} \coqdef{SECF.Imp.o:49}{o}{\coqdocbinder{o}} \coqdef{SECF.Imp.p:50}{p}{\coqdocbinder{p}},\coqdoceol
\coqdocindent{1.00em}
\coqref{SECF.Imp.m:47}{\coqdocvariable{m}} \coqexternalref{::nat scope:x '+' x}{http://rocq-prover.org/doc/V9.1.1/stdlib//Corelib.Init.Nat}{\coqdocnotation{+}} \coqref{SECF.Imp.n:48}{\coqdocvariable{n}} \coqexternalref{::nat scope:x '<=' x}{http://rocq-prover.org/doc/V9.1.1/stdlib//Corelib.Init.Peano}{\coqdocnotation{\ensuremath{\le}}} \coqref{SECF.Imp.n:48}{\coqdocvariable{n}} \coqexternalref{::nat scope:x '+' x}{http://rocq-prover.org/doc/V9.1.1/stdlib//Corelib.Init.Nat}{\coqdocnotation{+}} \coqref{SECF.Imp.o:49}{\coqdocvariable{o}} \coqexternalref{::type scope:x '/x5C' x}{http://rocq-prover.org/doc/V9.1.1/stdlib//Corelib.Init.Logic}{\coqdocnotation{\ensuremath{\land}}} \coqref{SECF.Imp.o:49}{\coqdocvariable{o}} \coqexternalref{::nat scope:x '+' x}{http://rocq-prover.org/doc/V9.1.1/stdlib//Corelib.Init.Nat}{\coqdocnotation{+}} 3 \coqexternalref{::type scope:x '=' x}{http://rocq-prover.org/doc/V9.1.1/stdlib//Corelib.Init.Logic}{\coqdocnotation{=}} \coqref{SECF.Imp.p:50}{\coqdocvariable{p}} \coqexternalref{::nat scope:x '+' x}{http://rocq-prover.org/doc/V9.1.1/stdlib//Corelib.Init.Nat}{\coqdocnotation{+}} 3 \coqexternalref{::type scope:x '->' x}{http://rocq-prover.org/doc/V9.1.1/stdlib//Corelib.Init.Logic}{\coqdocnotation{\ensuremath{\rightarrow}}}\coqdoceol
\coqdocindent{1.00em}
\coqref{SECF.Imp.m:47}{\coqdocvariable{m}} \coqexternalref{::nat scope:x '<=' x}{http://rocq-prover.org/doc/V9.1.1/stdlib//Corelib.Init.Peano}{\coqdocnotation{\ensuremath{\le}}} \coqref{SECF.Imp.p:50}{\coqdocvariable{p}}.\coqdoceol
\coqdocnoindent
\coqdockw{Proof}.\coqdoceol
\coqdocindent{1.00em}
\coqdoctac{intros}. \coqdocvar{lia}.\coqdoceol
\coqdocnoindent
\coqdockw{Qed}.\coqdoceol
\coqdocemptyline
\coqdocnoindent
\coqdockw{Example} \coqdef{SECF.Imp.AExp.add comm  lia}{add\_comm\_\_lia}{\coqdocdefinition{add\_comm\_\_lia}} : \coqdockw{\ensuremath{\forall}} \coqdef{SECF.Imp.m:51}{m}{\coqdocbinder{m}} \coqdef{SECF.Imp.n:52}{n}{\coqdocbinder{n}},\coqdoceol
\coqdocindent{2.00em}
\coqref{SECF.Imp.m:51}{\coqdocvariable{m}} \coqexternalref{::nat scope:x '+' x}{http://rocq-prover.org/doc/V9.1.1/stdlib//Corelib.Init.Nat}{\coqdocnotation{+}} \coqref{SECF.Imp.n:52}{\coqdocvariable{n}} \coqexternalref{::type scope:x '=' x}{http://rocq-prover.org/doc/V9.1.1/stdlib//Corelib.Init.Logic}{\coqdocnotation{=}} \coqref{SECF.Imp.n:52}{\coqdocvariable{n}} \coqexternalref{::nat scope:x '+' x}{http://rocq-prover.org/doc/V9.1.1/stdlib//Corelib.Init.Nat}{\coqdocnotation{+}} \coqref{SECF.Imp.m:51}{\coqdocvariable{m}}.\coqdoceol
\coqdocnoindent
\coqdockw{Proof}.\coqdoceol
\coqdocindent{1.00em}
\coqdoctac{intros}. \coqdocvar{lia}.\coqdoceol
\coqdocnoindent
\coqdockw{Qed}.\coqdoceol
\coqdocemptyline
\coqdocnoindent
\coqdockw{Example} \coqdef{SECF.Imp.AExp.add assoc  lia}{add\_assoc\_\_lia}{\coqdocdefinition{add\_assoc\_\_lia}} : \coqdockw{\ensuremath{\forall}} \coqdef{SECF.Imp.m:53}{m}{\coqdocbinder{m}} \coqdef{SECF.Imp.n:54}{n}{\coqdocbinder{n}} \coqdef{SECF.Imp.p:55}{p}{\coqdocbinder{p}},\coqdoceol
\coqdocindent{2.00em}
\coqref{SECF.Imp.m:53}{\coqdocvariable{m}} \coqexternalref{::nat scope:x '+' x}{http://rocq-prover.org/doc/V9.1.1/stdlib//Corelib.Init.Nat}{\coqdocnotation{+}} \coqexternalref{::nat scope:x '+' x}{http://rocq-prover.org/doc/V9.1.1/stdlib//Corelib.Init.Nat}{\coqdocnotation{(}}\coqref{SECF.Imp.n:54}{\coqdocvariable{n}} \coqexternalref{::nat scope:x '+' x}{http://rocq-prover.org/doc/V9.1.1/stdlib//Corelib.Init.Nat}{\coqdocnotation{+}} \coqref{SECF.Imp.p:55}{\coqdocvariable{p}}\coqexternalref{::nat scope:x '+' x}{http://rocq-prover.org/doc/V9.1.1/stdlib//Corelib.Init.Nat}{\coqdocnotation{)}} \coqexternalref{::type scope:x '=' x}{http://rocq-prover.org/doc/V9.1.1/stdlib//Corelib.Init.Logic}{\coqdocnotation{=}} \coqref{SECF.Imp.m:53}{\coqdocvariable{m}} \coqexternalref{::nat scope:x '+' x}{http://rocq-prover.org/doc/V9.1.1/stdlib//Corelib.Init.Nat}{\coqdocnotation{+}} \coqref{SECF.Imp.n:54}{\coqdocvariable{n}} \coqexternalref{::nat scope:x '+' x}{http://rocq-prover.org/doc/V9.1.1/stdlib//Corelib.Init.Nat}{\coqdocnotation{+}} \coqref{SECF.Imp.p:55}{\coqdocvariable{p}}.\coqdoceol
\coqdocnoindent
\coqdockw{Proof}.\coqdoceol
\coqdocindent{1.00em}
\coqdoctac{intros}. \coqdocvar{lia}.\coqdoceol
\coqdocnoindent
\coqdockw{Qed}.\coqdoceol
\coqdocemptyline
\end{coqdoccode}
(Note the \coqdockw{From} \coqdocvar{Stdlib} \coqdockw{Require} \coqdockw{Import} \coqdoclibrary{Lia}. at the top of
    this file, which makes \coqdocvar{lia} available.) 
\begin{coqdoccode}
\coqdocemptyline
\end{coqdoccode}
\subsection{A Few More Handy Tactics}



 Finally, here are some miscellaneous tactics that you may find
    convenient.



\begin{itemize}
\item  \coqdoctac{clear} \coqdocvar{H}: Delete hypothesis \coqdocvar{H} from the context.



\item  \coqdoctac{subst} \coqdocvar{x}: Given a variable \coqdocvar{x}, find an assumption \coqdocvar{x} = \coqdocvar{e} or
       \coqdocvar{e} = \coqdocvar{x} in the context, replace \coqdocvar{x} with \coqdocvar{e} throughout the
       context and current goal, and clear the assumption.



\item  \coqdoctac{subst}: Substitute away \textit{all} assumptions of the form \coqdocvar{x} = \coqdocvar{e}
       or \coqdocvar{e} = \coqdocvar{x} (where \coqdocvar{x} is a variable).



\item  \coqdoctac{rename}... \coqdocvar{into}...: Change the name of a hypothesis in the
       proof context.  For example, if the context includes a variable
       named \coqdocvar{x}, then \coqdoctac{rename} \coqdocvar{x} \coqdocvar{into} \coqdocvar{y} will change all occurrences
       of \coqdocvar{x} to \coqdocvar{y}.



\item  \coqdoctac{assumption}: Try to find a hypothesis \coqdocvar{H} in the context that
       exactly matches the goal; if one is found, solve the goal.



\item  \coqdocvar{contradiction}: Try to find a hypothesis \coqdocvar{H} in the context
       that is logically equivalent to \coqexternalref{False}{http://rocq-prover.org/doc/V9.1.1/stdlib//Corelib.Init.Logic}{\coqdocinductive{False}}.  If one is found,
       solve the goal.



\item  \coqdoctac{constructor}: Try to find a constructor \coqdocvar{c} (from some
       \coqdockw{Inductive} definition in the current environment) that can be
       applied to solve the current goal.  If one is found, behave
       like \coqdoctac{apply} \coqdocvar{c}.

\end{itemize}


    We'll see examples of all of these as we go along. 
\begin{coqdoccode}
\coqdocemptyline
\end{coqdoccode}
\section{Evaluation as a Relation}



 We have presented \coqref{SECF.Imp.aeval}{\coqdocdefinition{aeval}} and \coqref{SECF.Imp.beval}{\coqdocdefinition{beval}} as functions defined by
    \coqdockw{Fixpoint}s.  Another way to think about evaluation -- one that is
    often more flexible -- is as a \textit{relation} between expressions and
    their values.  This perspective leads to \coqdockw{Inductive} definitions
    like the following... 
\begin{coqdoccode}
\coqdocemptyline
\coqdocnoindent
\coqdockw{Module} \coqdef{SECF.Imp.AExp.aevalR first try}{aevalR\_first\_try}{\coqdocmodule{aevalR\_first\_try}}.\coqdoceol
\coqdocemptyline
\coqdocnoindent
\coqdockw{Inductive} \coqdef{SECF.Imp.AExp.aevalR first try.aevalR}{aevalR}{\coqdocinductive{aevalR}} : \coqref{SECF.Imp.AExp.aexp}{\coqdocinductive{aexp}} \coqexternalref{::type scope:x '->' x}{http://rocq-prover.org/doc/V9.1.1/stdlib//Corelib.Init.Logic}{\coqdocnotation{\ensuremath{\rightarrow}}} \coqexternalref{nat}{http://rocq-prover.org/doc/V9.1.1/stdlib//Corelib.Init.Datatypes}{\coqdocinductive{nat}} \coqexternalref{::type scope:x '->' x}{http://rocq-prover.org/doc/V9.1.1/stdlib//Corelib.Init.Logic}{\coqdocnotation{\ensuremath{\rightarrow}}} \coqdockw{Prop} :=\coqdoceol
\coqdocindent{1.00em}
\ensuremath{|} \coqdef{SECF.Imp.AExp.aevalR first try.E ANum}{E\_ANum}{\coqdocconstructor{E\_ANum}} (\coqdef{SECF.Imp.n:58}{n}{\coqdocbinder{n}} : \coqexternalref{nat}{http://rocq-prover.org/doc/V9.1.1/stdlib//Corelib.Init.Datatypes}{\coqdocinductive{nat}}) :\coqdoceol
\coqdocindent{3.00em}
\coqref{SECF.Imp.aevalR:56}{\coqdocinductive{aevalR}} (\coqref{SECF.Imp.AExp.ANum}{\coqdocconstructor{ANum}} \coqref{SECF.Imp.n:58}{\coqdocvariable{n}}) \coqref{SECF.Imp.n:58}{\coqdocvariable{n}}\coqdoceol
\coqdocindent{1.00em}
\ensuremath{|} \coqdef{SECF.Imp.AExp.aevalR first try.E APlus}{E\_APlus}{\coqdocconstructor{E\_APlus}} (\coqdef{SECF.Imp.e1:59}{e1}{\coqdocbinder{e1}} \coqdef{SECF.Imp.e2:60}{e2}{\coqdocbinder{e2}} : \coqref{SECF.Imp.AExp.aexp}{\coqdocinductive{aexp}}) (\coqdef{SECF.Imp.n1:61}{n1}{\coqdocbinder{n1}} \coqdef{SECF.Imp.n2:62}{n2}{\coqdocbinder{n2}} : \coqexternalref{nat}{http://rocq-prover.org/doc/V9.1.1/stdlib//Corelib.Init.Datatypes}{\coqdocinductive{nat}}) :\coqdoceol
\coqdocindent{3.00em}
\coqref{SECF.Imp.aevalR:56}{\coqdocinductive{aevalR}} \coqref{SECF.Imp.e1:59}{\coqdocvariable{e1}} \coqref{SECF.Imp.n1:61}{\coqdocvariable{n1}} \coqexternalref{::type scope:x '->' x}{http://rocq-prover.org/doc/V9.1.1/stdlib//Corelib.Init.Logic}{\coqdocnotation{\ensuremath{\rightarrow}}}\coqdoceol
\coqdocindent{3.00em}
\coqref{SECF.Imp.aevalR:56}{\coqdocinductive{aevalR}} \coqref{SECF.Imp.e2:60}{\coqdocvariable{e2}} \coqref{SECF.Imp.n2:62}{\coqdocvariable{n2}} \coqexternalref{::type scope:x '->' x}{http://rocq-prover.org/doc/V9.1.1/stdlib//Corelib.Init.Logic}{\coqdocnotation{\ensuremath{\rightarrow}}}\coqdoceol
\coqdocindent{3.00em}
\coqref{SECF.Imp.aevalR:56}{\coqdocinductive{aevalR}} (\coqref{SECF.Imp.AExp.APlus}{\coqdocconstructor{APlus}} \coqref{SECF.Imp.e1:59}{\coqdocvariable{e1}} \coqref{SECF.Imp.e2:60}{\coqdocvariable{e2}}) (\coqref{SECF.Imp.n1:61}{\coqdocvariable{n1}} \coqexternalref{::nat scope:x '+' x}{http://rocq-prover.org/doc/V9.1.1/stdlib//Corelib.Init.Nat}{\coqdocnotation{+}} \coqref{SECF.Imp.n2:62}{\coqdocvariable{n2}})\coqdoceol
\coqdocindent{1.00em}
\ensuremath{|} \coqdef{SECF.Imp.AExp.aevalR first try.E AMinus}{E\_AMinus}{\coqdocconstructor{E\_AMinus}} (\coqdef{SECF.Imp.e1:63}{e1}{\coqdocbinder{e1}} \coqdef{SECF.Imp.e2:64}{e2}{\coqdocbinder{e2}} : \coqref{SECF.Imp.AExp.aexp}{\coqdocinductive{aexp}}) (\coqdef{SECF.Imp.n1:65}{n1}{\coqdocbinder{n1}} \coqdef{SECF.Imp.n2:66}{n2}{\coqdocbinder{n2}} : \coqexternalref{nat}{http://rocq-prover.org/doc/V9.1.1/stdlib//Corelib.Init.Datatypes}{\coqdocinductive{nat}}) :\coqdoceol
\coqdocindent{3.00em}
\coqref{SECF.Imp.aevalR:56}{\coqdocinductive{aevalR}} \coqref{SECF.Imp.e1:63}{\coqdocvariable{e1}} \coqref{SECF.Imp.n1:65}{\coqdocvariable{n1}} \coqexternalref{::type scope:x '->' x}{http://rocq-prover.org/doc/V9.1.1/stdlib//Corelib.Init.Logic}{\coqdocnotation{\ensuremath{\rightarrow}}}\coqdoceol
\coqdocindent{3.00em}
\coqref{SECF.Imp.aevalR:56}{\coqdocinductive{aevalR}} \coqref{SECF.Imp.e2:64}{\coqdocvariable{e2}} \coqref{SECF.Imp.n2:66}{\coqdocvariable{n2}} \coqexternalref{::type scope:x '->' x}{http://rocq-prover.org/doc/V9.1.1/stdlib//Corelib.Init.Logic}{\coqdocnotation{\ensuremath{\rightarrow}}}\coqdoceol
\coqdocindent{3.00em}
\coqref{SECF.Imp.aevalR:56}{\coqdocinductive{aevalR}} (\coqref{SECF.Imp.AExp.AMinus}{\coqdocconstructor{AMinus}} \coqref{SECF.Imp.e1:63}{\coqdocvariable{e1}} \coqref{SECF.Imp.e2:64}{\coqdocvariable{e2}}) (\coqref{SECF.Imp.n1:65}{\coqdocvariable{n1}} \coqexternalref{::nat scope:x '-' x}{http://rocq-prover.org/doc/V9.1.1/stdlib//Corelib.Init.Nat}{\coqdocnotation{-}} \coqref{SECF.Imp.n2:66}{\coqdocvariable{n2}})\coqdoceol
\coqdocindent{1.00em}
\ensuremath{|} \coqdef{SECF.Imp.AExp.aevalR first try.E AMult}{E\_AMult}{\coqdocconstructor{E\_AMult}} (\coqdef{SECF.Imp.e1:67}{e1}{\coqdocbinder{e1}} \coqdef{SECF.Imp.e2:68}{e2}{\coqdocbinder{e2}} : \coqref{SECF.Imp.AExp.aexp}{\coqdocinductive{aexp}}) (\coqdef{SECF.Imp.n1:69}{n1}{\coqdocbinder{n1}} \coqdef{SECF.Imp.n2:70}{n2}{\coqdocbinder{n2}} : \coqexternalref{nat}{http://rocq-prover.org/doc/V9.1.1/stdlib//Corelib.Init.Datatypes}{\coqdocinductive{nat}}) :\coqdoceol
\coqdocindent{3.00em}
\coqref{SECF.Imp.aevalR:56}{\coqdocinductive{aevalR}} \coqref{SECF.Imp.e1:67}{\coqdocvariable{e1}} \coqref{SECF.Imp.n1:69}{\coqdocvariable{n1}} \coqexternalref{::type scope:x '->' x}{http://rocq-prover.org/doc/V9.1.1/stdlib//Corelib.Init.Logic}{\coqdocnotation{\ensuremath{\rightarrow}}}\coqdoceol
\coqdocindent{3.00em}
\coqref{SECF.Imp.aevalR:56}{\coqdocinductive{aevalR}} \coqref{SECF.Imp.e2:68}{\coqdocvariable{e2}} \coqref{SECF.Imp.n2:70}{\coqdocvariable{n2}} \coqexternalref{::type scope:x '->' x}{http://rocq-prover.org/doc/V9.1.1/stdlib//Corelib.Init.Logic}{\coqdocnotation{\ensuremath{\rightarrow}}}\coqdoceol
\coqdocindent{3.00em}
\coqref{SECF.Imp.aevalR:56}{\coqdocinductive{aevalR}} (\coqref{SECF.Imp.AExp.AMult}{\coqdocconstructor{AMult}} \coqref{SECF.Imp.e1:67}{\coqdocvariable{e1}} \coqref{SECF.Imp.e2:68}{\coqdocvariable{e2}}) (\coqref{SECF.Imp.n1:69}{\coqdocvariable{n1}} \coqexternalref{::nat scope:x '*' x}{http://rocq-prover.org/doc/V9.1.1/stdlib//Corelib.Init.Nat}{\coqdocnotation{\ensuremath{\times}}} \coqref{SECF.Imp.n2:70}{\coqdocvariable{n2}}).\coqdoceol
\coqdocemptyline
\coqdocnoindent
\coqdockw{Module} \coqdef{SECF.Imp.AExp.aevalR first try.HypothesisNames}{HypothesisNames}{\coqdocmodule{HypothesisNames}}.\coqdoceol
\coqdocemptyline
\end{coqdoccode}
A small notational aside. We could also write the definition of
    \coqref{SECF.Imp.aevalR extended.aevalR}{\coqdocinductive{aevalR}} as follow, with explicit names for the hypotheses in each
    case: 
\begin{coqdoccode}
\coqdocemptyline
\coqdocnoindent
\coqdockw{Inductive} \coqdef{SECF.Imp.AExp.aevalR first try.HypothesisNames.aevalR}{aevalR}{\coqdocinductive{aevalR}} : \coqref{SECF.Imp.AExp.aexp}{\coqdocinductive{aexp}} \coqexternalref{::type scope:x '->' x}{http://rocq-prover.org/doc/V9.1.1/stdlib//Corelib.Init.Logic}{\coqdocnotation{\ensuremath{\rightarrow}}} \coqexternalref{nat}{http://rocq-prover.org/doc/V9.1.1/stdlib//Corelib.Init.Datatypes}{\coqdocinductive{nat}} \coqexternalref{::type scope:x '->' x}{http://rocq-prover.org/doc/V9.1.1/stdlib//Corelib.Init.Logic}{\coqdocnotation{\ensuremath{\rightarrow}}} \coqdockw{Prop} :=\coqdoceol
\coqdocindent{1.00em}
\ensuremath{|} \coqdef{SECF.Imp.AExp.aevalR first try.HypothesisNames.E ANum}{E\_ANum}{\coqdocconstructor{E\_ANum}} (\coqdef{SECF.Imp.n:73}{n}{\coqdocbinder{n}} : \coqexternalref{nat}{http://rocq-prover.org/doc/V9.1.1/stdlib//Corelib.Init.Datatypes}{\coqdocinductive{nat}}) :\coqdoceol
\coqdocindent{3.00em}
\coqref{SECF.Imp.aevalR:71}{\coqdocinductive{aevalR}} (\coqref{SECF.Imp.AExp.ANum}{\coqdocconstructor{ANum}} \coqref{SECF.Imp.n:73}{\coqdocvariable{n}}) \coqref{SECF.Imp.n:73}{\coqdocvariable{n}}\coqdoceol
\coqdocindent{1.00em}
\ensuremath{|} \coqdef{SECF.Imp.AExp.aevalR first try.HypothesisNames.E APlus}{E\_APlus}{\coqdocconstructor{E\_APlus}} (\coqdef{SECF.Imp.e1:74}{e1}{\coqdocbinder{e1}} \coqdef{SECF.Imp.e2:75}{e2}{\coqdocbinder{e2}} : \coqref{SECF.Imp.AExp.aexp}{\coqdocinductive{aexp}}) (\coqdef{SECF.Imp.n1:76}{n1}{\coqdocbinder{n1}} \coqdef{SECF.Imp.n2:77}{n2}{\coqdocbinder{n2}} : \coqexternalref{nat}{http://rocq-prover.org/doc/V9.1.1/stdlib//Corelib.Init.Datatypes}{\coqdocinductive{nat}})\coqdoceol
\coqdocindent{3.00em}
(\coqdef{SECF.Imp.H1:78}{H1}{\coqdocbinder{H1}} : \coqref{SECF.Imp.aevalR:71}{\coqdocinductive{aevalR}} \coqref{SECF.Imp.e1:74}{\coqdocvariable{e1}} \coqref{SECF.Imp.n1:76}{\coqdocvariable{n1}})\coqdoceol
\coqdocindent{3.00em}
(\coqdef{SECF.Imp.H2:79}{H2}{\coqdocbinder{H2}} : \coqref{SECF.Imp.aevalR:71}{\coqdocinductive{aevalR}} \coqref{SECF.Imp.e2:75}{\coqdocvariable{e2}} \coqref{SECF.Imp.n2:77}{\coqdocvariable{n2}}) :\coqdoceol
\coqdocindent{3.00em}
\coqref{SECF.Imp.aevalR:71}{\coqdocinductive{aevalR}} (\coqref{SECF.Imp.AExp.APlus}{\coqdocconstructor{APlus}} \coqref{SECF.Imp.e1:74}{\coqdocvariable{e1}} \coqref{SECF.Imp.e2:75}{\coqdocvariable{e2}}) (\coqref{SECF.Imp.n1:76}{\coqdocvariable{n1}} \coqexternalref{::nat scope:x '+' x}{http://rocq-prover.org/doc/V9.1.1/stdlib//Corelib.Init.Nat}{\coqdocnotation{+}} \coqref{SECF.Imp.n2:77}{\coqdocvariable{n2}})\coqdoceol
\coqdocindent{1.00em}
\ensuremath{|} \coqdef{SECF.Imp.AExp.aevalR first try.HypothesisNames.E AMinus}{E\_AMinus}{\coqdocconstructor{E\_AMinus}} (\coqdef{SECF.Imp.e1:80}{e1}{\coqdocbinder{e1}} \coqdef{SECF.Imp.e2:81}{e2}{\coqdocbinder{e2}} : \coqref{SECF.Imp.AExp.aexp}{\coqdocinductive{aexp}}) (\coqdef{SECF.Imp.n1:82}{n1}{\coqdocbinder{n1}} \coqdef{SECF.Imp.n2:83}{n2}{\coqdocbinder{n2}} : \coqexternalref{nat}{http://rocq-prover.org/doc/V9.1.1/stdlib//Corelib.Init.Datatypes}{\coqdocinductive{nat}})\coqdoceol
\coqdocindent{3.00em}
(\coqdef{SECF.Imp.H1:84}{H1}{\coqdocbinder{H1}} : \coqref{SECF.Imp.aevalR:71}{\coqdocinductive{aevalR}} \coqref{SECF.Imp.e1:80}{\coqdocvariable{e1}} \coqref{SECF.Imp.n1:82}{\coqdocvariable{n1}})\coqdoceol
\coqdocindent{3.00em}
(\coqdef{SECF.Imp.H2:85}{H2}{\coqdocbinder{H2}} : \coqref{SECF.Imp.aevalR:71}{\coqdocinductive{aevalR}} \coqref{SECF.Imp.e2:81}{\coqdocvariable{e2}} \coqref{SECF.Imp.n2:83}{\coqdocvariable{n2}}) :\coqdoceol
\coqdocindent{3.00em}
\coqref{SECF.Imp.aevalR:71}{\coqdocinductive{aevalR}} (\coqref{SECF.Imp.AExp.AMinus}{\coqdocconstructor{AMinus}} \coqref{SECF.Imp.e1:80}{\coqdocvariable{e1}} \coqref{SECF.Imp.e2:81}{\coqdocvariable{e2}}) (\coqref{SECF.Imp.n1:82}{\coqdocvariable{n1}} \coqexternalref{::nat scope:x '-' x}{http://rocq-prover.org/doc/V9.1.1/stdlib//Corelib.Init.Nat}{\coqdocnotation{-}} \coqref{SECF.Imp.n2:83}{\coqdocvariable{n2}})\coqdoceol
\coqdocindent{1.00em}
\ensuremath{|} \coqdef{SECF.Imp.AExp.aevalR first try.HypothesisNames.E AMult}{E\_AMult}{\coqdocconstructor{E\_AMult}} (\coqdef{SECF.Imp.e1:86}{e1}{\coqdocbinder{e1}} \coqdef{SECF.Imp.e2:87}{e2}{\coqdocbinder{e2}} : \coqref{SECF.Imp.AExp.aexp}{\coqdocinductive{aexp}}) (\coqdef{SECF.Imp.n1:88}{n1}{\coqdocbinder{n1}} \coqdef{SECF.Imp.n2:89}{n2}{\coqdocbinder{n2}} : \coqexternalref{nat}{http://rocq-prover.org/doc/V9.1.1/stdlib//Corelib.Init.Datatypes}{\coqdocinductive{nat}})\coqdoceol
\coqdocindent{3.00em}
(\coqdef{SECF.Imp.H1:90}{H1}{\coqdocbinder{H1}} : \coqref{SECF.Imp.aevalR:71}{\coqdocinductive{aevalR}} \coqref{SECF.Imp.e1:86}{\coqdocvariable{e1}} \coqref{SECF.Imp.n1:88}{\coqdocvariable{n1}})\coqdoceol
\coqdocindent{3.00em}
(\coqdef{SECF.Imp.H2:91}{H2}{\coqdocbinder{H2}} : \coqref{SECF.Imp.aevalR:71}{\coqdocinductive{aevalR}} \coqref{SECF.Imp.e2:87}{\coqdocvariable{e2}} \coqref{SECF.Imp.n2:89}{\coqdocvariable{n2}}) :\coqdoceol
\coqdocindent{3.00em}
\coqref{SECF.Imp.aevalR:71}{\coqdocinductive{aevalR}} (\coqref{SECF.Imp.AExp.AMult}{\coqdocconstructor{AMult}} \coqref{SECF.Imp.e1:86}{\coqdocvariable{e1}} \coqref{SECF.Imp.e2:87}{\coqdocvariable{e2}}) (\coqref{SECF.Imp.n1:88}{\coqdocvariable{n1}} \coqexternalref{::nat scope:x '*' x}{http://rocq-prover.org/doc/V9.1.1/stdlib//Corelib.Init.Nat}{\coqdocnotation{\ensuremath{\times}}} \coqref{SECF.Imp.n2:89}{\coqdocvariable{n2}}).\coqdoceol
\coqdocemptyline
\end{coqdoccode}
This style gives us more control over the names that Rocq chooses
    during proofs involving \coqref{SECF.Imp.aevalR extended.aevalR}{\coqdocinductive{aevalR}}, at the cost of making the
    definition a little more verbose. 
\begin{coqdoccode}
\coqdocemptyline
\coqdocnoindent
\coqdockw{End} \coqref{SECF.Imp.AExp.aevalR first try.HypothesisNames}{\coqdocmodule{HypothesisNames}}.\coqdoceol
\coqdocemptyline
\end{coqdoccode}
It will be convenient to have an infix notation for
    \coqref{SECF.Imp.aevalR extended.aevalR}{\coqdocinductive{aevalR}}.  We'll write \coqdocvar{e} ==> \coqdocvar{n} to mean that arithmetic expression
    \coqdocvar{e} evaluates to value \coqdocvar{n}. 
\begin{coqdoccode}
\coqdocemptyline
\coqdocnoindent
\coqdockw{Notation} \coqdef{SECF.Imp.AExp.aevalR first try.::type scope:x '==>' x}{"}{"}e '==>' n"\coqdoceol
\coqdocindent{4.50em}
:= (\coqref{SECF.Imp.AExp.aevalR first try.aevalR}{\coqdocinductive{aevalR}} \coqdocvar{e} \coqdocvar{n})\coqdoceol
\coqdocindent{6.00em}
(\coqdoctac{at} \coqdockw{level} 90, \coqdoctac{left} \coqdockw{associativity})\coqdoceol
\coqdocindent{4.50em}
: \coqdocvar{type\_scope}.\coqdoceol
\coqdocemptyline
\coqdocnoindent
\coqdockw{End} \coqref{SECF.Imp.AExp.aevalR first try}{\coqdocmodule{aevalR\_first\_try}}.\coqdoceol
\coqdocemptyline
\end{coqdoccode}
As we saw in our case study of regular expressions in
    chapter \coqdocvar{IndProp}, Rocq provides a way to use this notation in
    the definition of \coqref{SECF.Imp.aevalR extended.aevalR}{\coqdocinductive{aevalR}} itself. 
\begin{coqdoccode}
\coqdocemptyline
\coqdocnoindent
\coqdockw{Reserved Notation} "e '==>' n" (\coqdoctac{at} \coqdockw{level} 90, \coqdoctac{left} \coqdockw{associativity}).\coqdoceol
\coqdocemptyline
\coqdocnoindent
\coqdockw{Inductive} \coqdef{SECF.Imp.AExp.aevalR}{aevalR}{\coqdocinductive{aevalR}} : \coqref{SECF.Imp.AExp.aexp}{\coqdocinductive{aexp}} \coqexternalref{::type scope:x '->' x}{http://rocq-prover.org/doc/V9.1.1/stdlib//Corelib.Init.Logic}{\coqdocnotation{\ensuremath{\rightarrow}}} \coqexternalref{nat}{http://rocq-prover.org/doc/V9.1.1/stdlib//Corelib.Init.Datatypes}{\coqdocinductive{nat}} \coqexternalref{::type scope:x '->' x}{http://rocq-prover.org/doc/V9.1.1/stdlib//Corelib.Init.Logic}{\coqdocnotation{\ensuremath{\rightarrow}}} \coqdockw{Prop} :=\coqdoceol
\coqdocindent{1.00em}
\ensuremath{|} \coqdef{SECF.Imp.AExp.E ANum}{E\_ANum}{\coqdocconstructor{E\_ANum}} (\coqdef{SECF.Imp.n:94}{n}{\coqdocbinder{n}} : \coqexternalref{nat}{http://rocq-prover.org/doc/V9.1.1/stdlib//Corelib.Init.Datatypes}{\coqdocinductive{nat}}) :\coqdoceol
\coqdocindent{3.00em}
\coqref{SECF.Imp.AExp.::type scope:x '==>' x}{\coqdocnotation{(}}\coqref{SECF.Imp.AExp.ANum}{\coqdocconstructor{ANum}} \coqref{SECF.Imp.n:94}{\coqdocvariable{n}}\coqref{SECF.Imp.AExp.::type scope:x '==>' x}{\coqdocnotation{)}} \coqref{SECF.Imp.AExp.::type scope:x '==>' x}{\coqdocnotation{==>}} \coqref{SECF.Imp.n:94}{\coqdocvariable{n}}\coqdoceol
\coqdocindent{1.00em}
\ensuremath{|} \coqdef{SECF.Imp.AExp.E APlus}{E\_APlus}{\coqdocconstructor{E\_APlus}} (\coqdef{SECF.Imp.e1:95}{e1}{\coqdocbinder{e1}} \coqdef{SECF.Imp.e2:96}{e2}{\coqdocbinder{e2}} : \coqref{SECF.Imp.AExp.aexp}{\coqdocinductive{aexp}}) (\coqdef{SECF.Imp.n1:97}{n1}{\coqdocbinder{n1}} \coqdef{SECF.Imp.n2:98}{n2}{\coqdocbinder{n2}} : \coqexternalref{nat}{http://rocq-prover.org/doc/V9.1.1/stdlib//Corelib.Init.Datatypes}{\coqdocinductive{nat}}) :\coqdoceol
\coqdocindent{3.00em}
\coqexternalref{::type scope:x '->' x}{http://rocq-prover.org/doc/V9.1.1/stdlib//Corelib.Init.Logic}{\coqdocnotation{(}}\coqref{SECF.Imp.e1:95}{\coqdocvariable{e1}} \coqref{SECF.Imp.AExp.::type scope:x '==>' x}{\coqdocnotation{==>}} \coqref{SECF.Imp.n1:97}{\coqdocvariable{n1}}\coqexternalref{::type scope:x '->' x}{http://rocq-prover.org/doc/V9.1.1/stdlib//Corelib.Init.Logic}{\coqdocnotation{)}} \coqexternalref{::type scope:x '->' x}{http://rocq-prover.org/doc/V9.1.1/stdlib//Corelib.Init.Logic}{\coqdocnotation{\ensuremath{\rightarrow}}}\coqdoceol
\coqdocindent{3.00em}
\coqexternalref{::type scope:x '->' x}{http://rocq-prover.org/doc/V9.1.1/stdlib//Corelib.Init.Logic}{\coqdocnotation{(}}\coqref{SECF.Imp.e2:96}{\coqdocvariable{e2}} \coqref{SECF.Imp.AExp.::type scope:x '==>' x}{\coqdocnotation{==>}} \coqref{SECF.Imp.n2:98}{\coqdocvariable{n2}}\coqexternalref{::type scope:x '->' x}{http://rocq-prover.org/doc/V9.1.1/stdlib//Corelib.Init.Logic}{\coqdocnotation{)}} \coqexternalref{::type scope:x '->' x}{http://rocq-prover.org/doc/V9.1.1/stdlib//Corelib.Init.Logic}{\coqdocnotation{\ensuremath{\rightarrow}}}\coqdoceol
\coqdocindent{3.00em}
\coqref{SECF.Imp.AExp.::type scope:x '==>' x}{\coqdocnotation{(}}\coqref{SECF.Imp.AExp.APlus}{\coqdocconstructor{APlus}} \coqref{SECF.Imp.e1:95}{\coqdocvariable{e1}} \coqref{SECF.Imp.e2:96}{\coqdocvariable{e2}}\coqref{SECF.Imp.AExp.::type scope:x '==>' x}{\coqdocnotation{)}}  \coqref{SECF.Imp.AExp.::type scope:x '==>' x}{\coqdocnotation{==>}} \coqref{SECF.Imp.AExp.::type scope:x '==>' x}{\coqdocnotation{(}}\coqref{SECF.Imp.n1:97}{\coqdocvariable{n1}} \coqexternalref{::nat scope:x '+' x}{http://rocq-prover.org/doc/V9.1.1/stdlib//Corelib.Init.Nat}{\coqdocnotation{+}} \coqref{SECF.Imp.n2:98}{\coqdocvariable{n2}}\coqref{SECF.Imp.AExp.::type scope:x '==>' x}{\coqdocnotation{)}}\coqdoceol
\coqdocindent{1.00em}
\ensuremath{|} \coqdef{SECF.Imp.AExp.E AMinus}{E\_AMinus}{\coqdocconstructor{E\_AMinus}} (\coqdef{SECF.Imp.e1:99}{e1}{\coqdocbinder{e1}} \coqdef{SECF.Imp.e2:100}{e2}{\coqdocbinder{e2}} : \coqref{SECF.Imp.AExp.aexp}{\coqdocinductive{aexp}}) (\coqdef{SECF.Imp.n1:101}{n1}{\coqdocbinder{n1}} \coqdef{SECF.Imp.n2:102}{n2}{\coqdocbinder{n2}} : \coqexternalref{nat}{http://rocq-prover.org/doc/V9.1.1/stdlib//Corelib.Init.Datatypes}{\coqdocinductive{nat}}) :\coqdoceol
\coqdocindent{3.00em}
\coqexternalref{::type scope:x '->' x}{http://rocq-prover.org/doc/V9.1.1/stdlib//Corelib.Init.Logic}{\coqdocnotation{(}}\coqref{SECF.Imp.e1:99}{\coqdocvariable{e1}} \coqref{SECF.Imp.AExp.::type scope:x '==>' x}{\coqdocnotation{==>}} \coqref{SECF.Imp.n1:101}{\coqdocvariable{n1}}\coqexternalref{::type scope:x '->' x}{http://rocq-prover.org/doc/V9.1.1/stdlib//Corelib.Init.Logic}{\coqdocnotation{)}} \coqexternalref{::type scope:x '->' x}{http://rocq-prover.org/doc/V9.1.1/stdlib//Corelib.Init.Logic}{\coqdocnotation{\ensuremath{\rightarrow}}}\coqdoceol
\coqdocindent{3.00em}
\coqexternalref{::type scope:x '->' x}{http://rocq-prover.org/doc/V9.1.1/stdlib//Corelib.Init.Logic}{\coqdocnotation{(}}\coqref{SECF.Imp.e2:100}{\coqdocvariable{e2}} \coqref{SECF.Imp.AExp.::type scope:x '==>' x}{\coqdocnotation{==>}} \coqref{SECF.Imp.n2:102}{\coqdocvariable{n2}}\coqexternalref{::type scope:x '->' x}{http://rocq-prover.org/doc/V9.1.1/stdlib//Corelib.Init.Logic}{\coqdocnotation{)}} \coqexternalref{::type scope:x '->' x}{http://rocq-prover.org/doc/V9.1.1/stdlib//Corelib.Init.Logic}{\coqdocnotation{\ensuremath{\rightarrow}}}\coqdoceol
\coqdocindent{3.00em}
\coqref{SECF.Imp.AExp.::type scope:x '==>' x}{\coqdocnotation{(}}\coqref{SECF.Imp.AExp.AMinus}{\coqdocconstructor{AMinus}} \coqref{SECF.Imp.e1:99}{\coqdocvariable{e1}} \coqref{SECF.Imp.e2:100}{\coqdocvariable{e2}}\coqref{SECF.Imp.AExp.::type scope:x '==>' x}{\coqdocnotation{)}} \coqref{SECF.Imp.AExp.::type scope:x '==>' x}{\coqdocnotation{==>}} \coqref{SECF.Imp.AExp.::type scope:x '==>' x}{\coqdocnotation{(}}\coqref{SECF.Imp.n1:101}{\coqdocvariable{n1}} \coqexternalref{::nat scope:x '-' x}{http://rocq-prover.org/doc/V9.1.1/stdlib//Corelib.Init.Nat}{\coqdocnotation{-}} \coqref{SECF.Imp.n2:102}{\coqdocvariable{n2}}\coqref{SECF.Imp.AExp.::type scope:x '==>' x}{\coqdocnotation{)}}\coqdoceol
\coqdocindent{1.00em}
\ensuremath{|} \coqdef{SECF.Imp.AExp.E AMult}{E\_AMult}{\coqdocconstructor{E\_AMult}} (\coqdef{SECF.Imp.e1:103}{e1}{\coqdocbinder{e1}} \coqdef{SECF.Imp.e2:104}{e2}{\coqdocbinder{e2}} : \coqref{SECF.Imp.AExp.aexp}{\coqdocinductive{aexp}}) (\coqdef{SECF.Imp.n1:105}{n1}{\coqdocbinder{n1}} \coqdef{SECF.Imp.n2:106}{n2}{\coqdocbinder{n2}} : \coqexternalref{nat}{http://rocq-prover.org/doc/V9.1.1/stdlib//Corelib.Init.Datatypes}{\coqdocinductive{nat}}) :\coqdoceol
\coqdocindent{3.00em}
\coqexternalref{::type scope:x '->' x}{http://rocq-prover.org/doc/V9.1.1/stdlib//Corelib.Init.Logic}{\coqdocnotation{(}}\coqref{SECF.Imp.e1:103}{\coqdocvariable{e1}} \coqref{SECF.Imp.AExp.::type scope:x '==>' x}{\coqdocnotation{==>}} \coqref{SECF.Imp.n1:105}{\coqdocvariable{n1}}\coqexternalref{::type scope:x '->' x}{http://rocq-prover.org/doc/V9.1.1/stdlib//Corelib.Init.Logic}{\coqdocnotation{)}} \coqexternalref{::type scope:x '->' x}{http://rocq-prover.org/doc/V9.1.1/stdlib//Corelib.Init.Logic}{\coqdocnotation{\ensuremath{\rightarrow}}}\coqdoceol
\coqdocindent{3.00em}
\coqexternalref{::type scope:x '->' x}{http://rocq-prover.org/doc/V9.1.1/stdlib//Corelib.Init.Logic}{\coqdocnotation{(}}\coqref{SECF.Imp.e2:104}{\coqdocvariable{e2}} \coqref{SECF.Imp.AExp.::type scope:x '==>' x}{\coqdocnotation{==>}} \coqref{SECF.Imp.n2:106}{\coqdocvariable{n2}}\coqexternalref{::type scope:x '->' x}{http://rocq-prover.org/doc/V9.1.1/stdlib//Corelib.Init.Logic}{\coqdocnotation{)}} \coqexternalref{::type scope:x '->' x}{http://rocq-prover.org/doc/V9.1.1/stdlib//Corelib.Init.Logic}{\coqdocnotation{\ensuremath{\rightarrow}}}\coqdoceol
\coqdocindent{3.00em}
\coqref{SECF.Imp.AExp.::type scope:x '==>' x}{\coqdocnotation{(}}\coqref{SECF.Imp.AExp.AMult}{\coqdocconstructor{AMult}} \coqref{SECF.Imp.e1:103}{\coqdocvariable{e1}} \coqref{SECF.Imp.e2:104}{\coqdocvariable{e2}}\coqref{SECF.Imp.AExp.::type scope:x '==>' x}{\coqdocnotation{)}}  \coqref{SECF.Imp.AExp.::type scope:x '==>' x}{\coqdocnotation{==>}} \coqref{SECF.Imp.AExp.::type scope:x '==>' x}{\coqdocnotation{(}}\coqref{SECF.Imp.n1:105}{\coqdocvariable{n1}} \coqexternalref{::nat scope:x '*' x}{http://rocq-prover.org/doc/V9.1.1/stdlib//Corelib.Init.Nat}{\coqdocnotation{\ensuremath{\times}}} \coqref{SECF.Imp.n2:106}{\coqdocvariable{n2}}\coqref{SECF.Imp.AExp.::type scope:x '==>' x}{\coqdocnotation{)}}\coqdoceol
\coqdocnoindent
\coqdoceol
\coqdocindent{1.00em}
\coqdockw{where} \coqdef{SECF.Imp.AExp.::type scope:x '==>' x}{"}{"}e '==>' n" := (\coqref{SECF.Imp.aevalR:93}{\coqdocinductive{aevalR}} \coqdocvar{e} \coqdocvar{n}) : \coqdocvar{type\_scope}.\coqdoceol
\coqdocemptyline
\end{coqdoccode}
\subsection{Inference Rule Notation}



 In informal discussions, it is convenient to write the rules
    for \coqref{SECF.Imp.aevalR extended.aevalR}{\coqdocinductive{aevalR}} and similar relations in the more readable graphical
    form of \textit{inference rules}, where the premises above the line
    justify the conclusion below the line.


    For example, the constructor \coqref{SECF.Imp.aevalR extended.E APlus}{\coqdocconstructor{E\_APlus}}...


      | E\_APlus : forall (e1 e2 : aexp) (n1 n2 : nat),
          aevalR e1 n1 ->
          aevalR e2 n2 ->
          aevalR (APlus e1 e2) (n1 + n2)


    ...can be written like this as an inference rule:


                               e1 ==> n1
                               e2 ==> n2
\par
\noindent\hrulefill\par
\noindent{}                (E\_APlus)
                         APlus e1 e2 ==> n1+n2


 Formally, there is nothing deep about inference rules: they
    are just implications.


    You can read the rule name on the right as the name of the
    constructor and read each of the linebreaks between the premises
    above the line (as well as the line itself) as \ensuremath{\rightarrow}.


    All the variables mentioned in the rule (\coqdocvar{e1}, \coqdocvar{n1}, etc.) are
    implicitly bound by universal quantifiers at the beginning. (Such
    variables are often called \textit{metavariables} to distinguish them
    from the variables of the language we are defining.  At the
    moment, our arithmetic expressions don't include variables, but
    we'll soon be adding them.)


    The whole collection of rules is understood as being wrapped in an
    \coqdockw{Inductive} declaration.  In informal prose, this is sometimes
    indicated by saying something like ``Let \coqref{SECF.Imp.aevalR extended.aevalR}{\coqdocinductive{aevalR}} be the smallest
    relation closed under the following rules...''. 

 For example, we could define ==> as the smallest relation
    closed under these rules:
\par
\noindent\hrulefill\par
\noindent{}                               (E\_ANum)
                             ANum n ==> n


                               e1 ==> n1
                               e2 ==> n2
\par
\noindent\hrulefill\par
\noindent{}                         (E\_APlus)
                         APlus e1 e2 ==> n1+n2


                               e1 ==> n1
                               e2 ==> n2
\par
\noindent\hrulefill\par
\noindent{}                        (E\_AMinus)
                        AMinus e1 e2 ==> n1-n2


                               e1 ==> n1
                               e2 ==> n2
\par
\noindent\hrulefill\par
\noindent{}                         (E\_AMult)
                         AMult e1 e2 ==> n1*n2


\paragraph{Exercise: 1 star, standard, optional (beval\_rules)}




    Here, again, is the Rocq definition of the \coqref{SECF.Imp.beval}{\coqdocdefinition{beval}} function:


  Fixpoint beval (e : bexp) : bool :=
    match e with
    | BTrue       => true
    | BFalse      => false
    | BEq a1 a2   => (aeval a1) =? (aeval a2)
    | BNeq a1 a2  => negb ((aeval a1) =? (aeval a2))
    | BLe a1 a2   => (aeval a1) <=? (aeval a2)
    | BGt a1 a2   => \~{}((aeval a1) <=? (aeval a2))
    | BNot b      => negb (beval b)
    | BAnd b1 b2  => andb (beval b1) (beval b2)
    end.


    Write out a corresponding definition of boolean evaluation as a
    relation (in inference rule notation). 
\begin{coqdoccode}
\coqdocemptyline
\coqdocnoindent
\coqdockw{Definition} \coqdef{SECF.Imp.AExp.manual grade for beval rules}{manual\_grade\_for\_beval\_rules}{\coqdocdefinition{manual\_grade\_for\_beval\_rules}} : \coqexternalref{option}{http://rocq-prover.org/doc/V9.1.1/stdlib//Corelib.Init.Datatypes}{\coqdocinductive{option}} (\coqexternalref{nat}{http://rocq-prover.org/doc/V9.1.1/stdlib//Corelib.Init.Datatypes}{\coqdocinductive{nat}}\coqexternalref{::type scope:x '*' x}{http://rocq-prover.org/doc/V9.1.1/stdlib//Corelib.Init.Datatypes}{\coqdocnotation{\ensuremath{\times}}}\coqdocinductive{string}) := \coqexternalref{None}{http://rocq-prover.org/doc/V9.1.1/stdlib//Corelib.Init.Datatypes}{\coqdocconstructor{None}}.\coqdoceol
\end{coqdoccode}
\ensuremath{\Box} 
\begin{coqdoccode}
\coqdocemptyline
\end{coqdoccode}
\subsection{Equivalence of the Definitions}



 It is straightforward to prove that the relational and functional
    definitions of evaluation agree: 
\begin{coqdoccode}
\coqdocemptyline
\coqdocnoindent
\coqdockw{Theorem} \coqdef{SECF.Imp.AExp.aevalR iff aeval}{aevalR\_iff\_aeval}{\coqdoclemma{aevalR\_iff\_aeval}} : \coqdockw{\ensuremath{\forall}} \coqdef{SECF.Imp.a:107}{a}{\coqdocbinder{a}} \coqdef{SECF.Imp.n:108}{n}{\coqdocbinder{n}},\coqdoceol
\coqdocindent{1.00em}
\coqexternalref{::type scope:x '<->' x}{http://rocq-prover.org/doc/V9.1.1/stdlib//Corelib.Init.Logic}{\coqdocnotation{(}}\coqref{SECF.Imp.a:107}{\coqdocvariable{a}} \coqref{SECF.Imp.AExp.::type scope:x '==>' x}{\coqdocnotation{==>}} \coqref{SECF.Imp.n:108}{\coqdocvariable{n}}\coqexternalref{::type scope:x '<->' x}{http://rocq-prover.org/doc/V9.1.1/stdlib//Corelib.Init.Logic}{\coqdocnotation{)}} \coqexternalref{::type scope:x '<->' x}{http://rocq-prover.org/doc/V9.1.1/stdlib//Corelib.Init.Logic}{\coqdocnotation{\ensuremath{\leftrightarrow}}} \coqref{SECF.Imp.AExp.aeval}{\coqdocdefinition{aeval}} \coqref{SECF.Imp.a:107}{\coqdocvariable{a}} \coqexternalref{::type scope:x '=' x}{http://rocq-prover.org/doc/V9.1.1/stdlib//Corelib.Init.Logic}{\coqdocnotation{=}} \coqref{SECF.Imp.n:108}{\coqdocvariable{n}}.\coqdoceol
\coqdocnoindent
\coqdockw{Proof}.\coqdoceol
\coqdocindent{1.00em}
\coqdoctac{split}.\coqdoceol
\coqdocindent{1.00em}
- \coqdoceol
\coqdocindent{2.00em}
\coqdoctac{intros} \coqdocvar{H}.\coqdoceol
\coqdocindent{2.00em}
\coqdoctac{induction} \coqdocvar{H}; \coqdoctac{simpl}.\coqdoceol
\coqdocindent{2.00em}
+ \coqdoceol
\coqdocindent{3.00em}
\coqdoctac{reflexivity}.\coqdoceol
\coqdocindent{2.00em}
+ \coqdoceol
\coqdocindent{3.00em}
\coqdoctac{rewrite} \coqdocvar{IHaevalR1}. \coqdoctac{rewrite} \coqdocvar{IHaevalR2}. \coqdoctac{reflexivity}.\coqdoceol
\coqdocindent{2.00em}
+ \coqdoceol
\coqdocindent{3.00em}
\coqdoctac{rewrite} \coqdocvar{IHaevalR1}. \coqdoctac{rewrite} \coqdocvar{IHaevalR2}. \coqdoctac{reflexivity}.\coqdoceol
\coqdocindent{2.00em}
+ \coqdoceol
\coqdocindent{3.00em}
\coqdoctac{rewrite} \coqdocvar{IHaevalR1}. \coqdoctac{rewrite} \coqdocvar{IHaevalR2}. \coqdoctac{reflexivity}.\coqdoceol
\coqdocindent{1.00em}
- \coqdoceol
\coqdocindent{2.00em}
\coqdoctac{generalize} \coqdoctac{dependent} \coqdocvar{n}.\coqdoceol
\coqdocindent{2.00em}
\coqdoctac{induction} \coqdocvar{a};\coqdoceol
\coqdocindent{3.50em}
\coqdoctac{simpl}; \coqdoctac{intros}; \coqdoctac{subst}.\coqdoceol
\coqdocindent{2.00em}
+ \coqdoceol
\coqdocindent{3.00em}
\coqdoctac{apply} \coqref{SECF.Imp.AExp.E ANum}{\coqdocconstructor{E\_ANum}}.\coqdoceol
\coqdocindent{2.00em}
+ \coqdoceol
\coqdocindent{3.00em}
\coqdoctac{apply} \coqref{SECF.Imp.AExp.E APlus}{\coqdocconstructor{E\_APlus}}.\coqdoceol
\coqdocindent{3.00em}
\ensuremath{\times} \coqdoctac{apply} \coqdocvar{IHa1}. \coqdoctac{reflexivity}.\coqdoceol
\coqdocindent{3.00em}
\ensuremath{\times} \coqdoctac{apply} \coqdocvar{IHa2}. \coqdoctac{reflexivity}.\coqdoceol
\coqdocindent{2.00em}
+ \coqdoceol
\coqdocindent{3.00em}
\coqdoctac{apply} \coqref{SECF.Imp.AExp.E AMinus}{\coqdocconstructor{E\_AMinus}}.\coqdoceol
\coqdocindent{3.00em}
\ensuremath{\times} \coqdoctac{apply} \coqdocvar{IHa1}. \coqdoctac{reflexivity}.\coqdoceol
\coqdocindent{3.00em}
\ensuremath{\times} \coqdoctac{apply} \coqdocvar{IHa2}. \coqdoctac{reflexivity}.\coqdoceol
\coqdocindent{2.00em}
+ \coqdoceol
\coqdocindent{3.00em}
\coqdoctac{apply} \coqref{SECF.Imp.AExp.E AMult}{\coqdocconstructor{E\_AMult}}.\coqdoceol
\coqdocindent{3.00em}
\ensuremath{\times} \coqdoctac{apply} \coqdocvar{IHa1}. \coqdoctac{reflexivity}.\coqdoceol
\coqdocindent{3.00em}
\ensuremath{\times} \coqdoctac{apply} \coqdocvar{IHa2}. \coqdoctac{reflexivity}.\coqdoceol
\coqdocnoindent
\coqdockw{Qed}.\coqdoceol
\coqdocemptyline
\end{coqdoccode}
Again, we can make the proof quite a bit shorter using some
    tacticals. 
\begin{coqdoccode}
\coqdocemptyline
\coqdocnoindent
\coqdockw{Theorem} \coqdef{SECF.Imp.AExp.aevalR iff aeval'}{aevalR\_iff\_aeval'}{\coqdoclemma{aevalR\_iff\_aeval'}} : \coqdockw{\ensuremath{\forall}} \coqdef{SECF.Imp.a:109}{a}{\coqdocbinder{a}} \coqdef{SECF.Imp.n:110}{n}{\coqdocbinder{n}},\coqdoceol
\coqdocindent{1.00em}
\coqexternalref{::type scope:x '<->' x}{http://rocq-prover.org/doc/V9.1.1/stdlib//Corelib.Init.Logic}{\coqdocnotation{(}}\coqref{SECF.Imp.a:109}{\coqdocvariable{a}} \coqref{SECF.Imp.AExp.::type scope:x '==>' x}{\coqdocnotation{==>}} \coqref{SECF.Imp.n:110}{\coqdocvariable{n}}\coqexternalref{::type scope:x '<->' x}{http://rocq-prover.org/doc/V9.1.1/stdlib//Corelib.Init.Logic}{\coqdocnotation{)}} \coqexternalref{::type scope:x '<->' x}{http://rocq-prover.org/doc/V9.1.1/stdlib//Corelib.Init.Logic}{\coqdocnotation{\ensuremath{\leftrightarrow}}} \coqref{SECF.Imp.AExp.aeval}{\coqdocdefinition{aeval}} \coqref{SECF.Imp.a:109}{\coqdocvariable{a}} \coqexternalref{::type scope:x '=' x}{http://rocq-prover.org/doc/V9.1.1/stdlib//Corelib.Init.Logic}{\coqdocnotation{=}} \coqref{SECF.Imp.n:110}{\coqdocvariable{n}}.\coqdoceol
\coqdocnoindent
\coqdockw{Proof}.\coqdoceol
\coqdocindent{1.00em}
\coqdoctac{split}.\coqdoceol
\coqdocindent{1.00em}
- \coqdoceol
\coqdocindent{2.00em}
\coqdoctac{intros} \coqdocvar{H}; \coqdoctac{induction} \coqdocvar{H}; \coqdoctac{subst}; \coqdoctac{reflexivity}.\coqdoceol
\coqdocindent{1.00em}
- \coqdoceol
\coqdocindent{2.00em}
\coqdoctac{generalize} \coqdoctac{dependent} \coqdocvar{n}.\coqdoceol
\coqdocindent{2.00em}
\coqdoctac{induction} \coqdocvar{a}; \coqdoctac{simpl}; \coqdoctac{intros}; \coqdoctac{subst}; \coqdoctac{constructor};\coqdoceol
\coqdocindent{3.50em}
\coqdoctac{try} \coqdoctac{apply} \coqdocvar{IHa1}; \coqdoctac{try} \coqdoctac{apply} \coqdocvar{IHa2}; \coqdoctac{reflexivity}.\coqdoceol
\coqdocnoindent
\coqdockw{Qed}.\coqdoceol
\coqdocemptyline
\end{coqdoccode}
\paragraph{Exercise: 3 stars, standard (bevalR)}




    Write a relation \coqref{SECF.Imp.AExp.bevalR}{\coqdocinductive{bevalR}} in the same style as
    \coqref{SECF.Imp.aevalR extended.aevalR}{\coqdocinductive{aevalR}}, and prove that it is equivalent to \coqref{SECF.Imp.beval}{\coqdocdefinition{beval}}. 
\begin{coqdoccode}
\coqdocemptyline
\coqdocnoindent
\coqdockw{Reserved Notation} "e '==>b' b" (\coqdoctac{at} \coqdockw{level} 90, \coqdoctac{left} \coqdockw{associativity}).\coqdoceol
\coqdocnoindent
\coqdockw{Inductive} \coqdef{SECF.Imp.AExp.bevalR}{bevalR}{\coqdocinductive{bevalR}}: \coqref{SECF.Imp.AExp.bexp}{\coqdocinductive{bexp}} \coqexternalref{::type scope:x '->' x}{http://rocq-prover.org/doc/V9.1.1/stdlib//Corelib.Init.Logic}{\coqdocnotation{\ensuremath{\rightarrow}}} \coqexternalref{bool}{http://rocq-prover.org/doc/V9.1.1/stdlib//Corelib.Init.Datatypes}{\coqdocinductive{bool}} \coqexternalref{::type scope:x '->' x}{http://rocq-prover.org/doc/V9.1.1/stdlib//Corelib.Init.Logic}{\coqdocnotation{\ensuremath{\rightarrow}}} \coqdockw{Prop} :=\coqdoceol
\coqdocnoindent
\coqdoceol
\coqdocnoindent
\coqdockw{where} \coqdef{SECF.Imp.AExp.::type scope:x '==>b' x}{"}{"}e '==>b' b" := (\coqref{SECF.Imp.bevalR:112}{\coqdocinductive{bevalR}} \coqdocvar{e} \coqdocvar{b}) : \coqdocvar{type\_scope}\coqdoceol
\coqdocnoindent
.\coqdoceol
\coqdocemptyline
\coqdocnoindent
\coqdockw{Lemma} \coqdef{SECF.Imp.AExp.bevalR iff beval}{bevalR\_iff\_beval}{\coqdoclemma{bevalR\_iff\_beval}} : \coqdockw{\ensuremath{\forall}} \coqdef{SECF.Imp.b:113}{b}{\coqdocbinder{b}} \coqdef{SECF.Imp.bv:114}{bv}{\coqdocbinder{bv}},\coqdoceol
\coqdocindent{1.00em}
\coqref{SECF.Imp.b:113}{\coqdocvariable{b}} \coqref{SECF.Imp.AExp.::type scope:x '==>b' x}{\coqdocnotation{==>}}\coqref{SECF.Imp.AExp.::type scope:x '==>b' x}{\coqdocnotation{b}} \coqref{SECF.Imp.bv:114}{\coqdocvariable{bv}} \coqexternalref{::type scope:x '<->' x}{http://rocq-prover.org/doc/V9.1.1/stdlib//Corelib.Init.Logic}{\coqdocnotation{\ensuremath{\leftrightarrow}}} \coqref{SECF.Imp.AExp.beval}{\coqdocdefinition{beval}} \coqref{SECF.Imp.b:113}{\coqdocvariable{b}} \coqexternalref{::type scope:x '=' x}{http://rocq-prover.org/doc/V9.1.1/stdlib//Corelib.Init.Logic}{\coqdocnotation{=}} \coqref{SECF.Imp.bv:114}{\coqdocvariable{bv}}.\coqdoceol
\coqdocnoindent
\coqdockw{Proof}.\coqdoceol
 \coqdocvar{Admitted}.\coqdoceol
\end{coqdoccode}
\ensuremath{\Box} 
\begin{coqdoccode}
\coqdocemptyline
\coqdocnoindent
\coqdockw{End} \coqref{SECF.Imp.AExp}{\coqdocmodule{AExp}}.\coqdoceol
\coqdocemptyline
\end{coqdoccode}
\subsection{Computational vs. Relational Definitions}



 For the definitions of evaluation for arithmetic and boolean
    expressions, the choice of whether to use functional or relational
    definitions is mainly a matter of taste: either way works fine.


    However, there are many situations where relational definitions of
    evaluation work much better than functional ones.  
\begin{coqdoccode}
\coqdocemptyline
\coqdocnoindent
\coqdockw{Module} \coqdef{SECF.Imp.aevalR division}{aevalR\_division}{\coqdocmodule{aevalR\_division}}.\coqdoceol
\coqdocemptyline
\end{coqdoccode}
For example, suppose that we wanted to extend the arithmetic
    operations with division: 
\begin{coqdoccode}
\coqdocemptyline
\coqdocnoindent
\coqdockw{Inductive} \coqdef{SECF.Imp.aevalR division.aexp}{aexp}{\coqdocinductive{aexp}} : \coqdockw{Type} :=\coqdoceol
\coqdocindent{1.00em}
\ensuremath{|} \coqdef{SECF.Imp.aevalR division.ANum}{ANum}{\coqdocconstructor{ANum}} (\coqdef{SECF.Imp.n:117}{n}{\coqdocbinder{n}} : \coqexternalref{nat}{http://rocq-prover.org/doc/V9.1.1/stdlib//Corelib.Init.Datatypes}{\coqdocinductive{nat}})\coqdoceol
\coqdocindent{1.00em}
\ensuremath{|} \coqdef{SECF.Imp.aevalR division.APlus}{APlus}{\coqdocconstructor{APlus}} (\coqdef{SECF.Imp.a1:118}{a1}{\coqdocbinder{a1}} \coqdef{SECF.Imp.a2:119}{a2}{\coqdocbinder{a2}} : \coqref{SECF.Imp.aexp:115}{\coqdocinductive{aexp}})\coqdoceol
\coqdocindent{1.00em}
\ensuremath{|} \coqdef{SECF.Imp.aevalR division.AMinus}{AMinus}{\coqdocconstructor{AMinus}} (\coqdef{SECF.Imp.a1:120}{a1}{\coqdocbinder{a1}} \coqdef{SECF.Imp.a2:121}{a2}{\coqdocbinder{a2}} : \coqref{SECF.Imp.aexp:115}{\coqdocinductive{aexp}})\coqdoceol
\coqdocindent{1.00em}
\ensuremath{|} \coqdef{SECF.Imp.aevalR division.AMult}{AMult}{\coqdocconstructor{AMult}} (\coqdef{SECF.Imp.a1:122}{a1}{\coqdocbinder{a1}} \coqdef{SECF.Imp.a2:123}{a2}{\coqdocbinder{a2}} : \coqref{SECF.Imp.aexp:115}{\coqdocinductive{aexp}})\coqdoceol
\coqdocindent{1.00em}
\ensuremath{|} \coqdef{SECF.Imp.aevalR division.ADiv}{ADiv}{\coqdocconstructor{ADiv}} (\coqdef{SECF.Imp.a1:124}{a1}{\coqdocbinder{a1}} \coqdef{SECF.Imp.a2:125}{a2}{\coqdocbinder{a2}} : \coqref{SECF.Imp.aexp:115}{\coqdocinductive{aexp}}). \coqdocemptyline
\end{coqdoccode}
Extending the definition of \coqref{SECF.Imp.aeval}{\coqdocdefinition{aeval}} to handle this new
    operation would not be straightforward (what should we return as
    the result of \coqref{SECF.Imp.aevalR division.ADiv}{\coqdocconstructor{ADiv}} (\coqref{SECF.SpecCT.ANum}{\coqdocconstructor{ANum}} 5) (\coqref{SECF.SpecCT.ANum}{\coqdocconstructor{ANum}} 0)?).  But extending \coqref{SECF.Imp.aevalR extended.aevalR}{\coqdocinductive{aevalR}}
    is very easy. 
\begin{coqdoccode}
\coqdocemptyline
\coqdocnoindent
\coqdockw{Reserved Notation} "e '==>' n"\coqdoceol
\coqdocindent{9.00em}
(\coqdoctac{at} \coqdockw{level} 90, \coqdoctac{left} \coqdockw{associativity}).\coqdoceol
\coqdocemptyline
\coqdocnoindent
\coqdockw{Inductive} \coqdef{SECF.Imp.aevalR division.aevalR}{aevalR}{\coqdocinductive{aevalR}} : \coqref{SECF.Imp.aevalR division.aexp}{\coqdocinductive{aexp}} \coqexternalref{::type scope:x '->' x}{http://rocq-prover.org/doc/V9.1.1/stdlib//Corelib.Init.Logic}{\coqdocnotation{\ensuremath{\rightarrow}}} \coqexternalref{nat}{http://rocq-prover.org/doc/V9.1.1/stdlib//Corelib.Init.Datatypes}{\coqdocinductive{nat}} \coqexternalref{::type scope:x '->' x}{http://rocq-prover.org/doc/V9.1.1/stdlib//Corelib.Init.Logic}{\coqdocnotation{\ensuremath{\rightarrow}}} \coqdockw{Prop} :=\coqdoceol
\coqdocindent{1.00em}
\ensuremath{|} \coqdef{SECF.Imp.aevalR division.E ANum}{E\_ANum}{\coqdocconstructor{E\_ANum}} (\coqdef{SECF.Imp.n:128}{n}{\coqdocbinder{n}} : \coqexternalref{nat}{http://rocq-prover.org/doc/V9.1.1/stdlib//Corelib.Init.Datatypes}{\coqdocinductive{nat}}) :\coqdoceol
\coqdocindent{3.00em}
\coqref{SECF.Imp.aevalR division.::type scope:x '==>' x}{\coqdocnotation{(}}\coqref{SECF.Imp.aevalR division.ANum}{\coqdocconstructor{ANum}} \coqref{SECF.Imp.n:128}{\coqdocvariable{n}}\coqref{SECF.Imp.aevalR division.::type scope:x '==>' x}{\coqdocnotation{)}} \coqref{SECF.Imp.aevalR division.::type scope:x '==>' x}{\coqdocnotation{==>}} \coqref{SECF.Imp.n:128}{\coqdocvariable{n}}\coqdoceol
\coqdocindent{1.00em}
\ensuremath{|} \coqdef{SECF.Imp.aevalR division.E APlus}{E\_APlus}{\coqdocconstructor{E\_APlus}} (\coqdef{SECF.Imp.a1:129}{a1}{\coqdocbinder{a1}} \coqdef{SECF.Imp.a2:130}{a2}{\coqdocbinder{a2}} : \coqref{SECF.Imp.aevalR division.aexp}{\coqdocinductive{aexp}}) (\coqdef{SECF.Imp.n1:131}{n1}{\coqdocbinder{n1}} \coqdef{SECF.Imp.n2:132}{n2}{\coqdocbinder{n2}} : \coqexternalref{nat}{http://rocq-prover.org/doc/V9.1.1/stdlib//Corelib.Init.Datatypes}{\coqdocinductive{nat}}) :\coqdoceol
\coqdocindent{3.00em}
\coqexternalref{::type scope:x '->' x}{http://rocq-prover.org/doc/V9.1.1/stdlib//Corelib.Init.Logic}{\coqdocnotation{(}}\coqref{SECF.Imp.a1:129}{\coqdocvariable{a1}} \coqref{SECF.Imp.aevalR division.::type scope:x '==>' x}{\coqdocnotation{==>}} \coqref{SECF.Imp.n1:131}{\coqdocvariable{n1}}\coqexternalref{::type scope:x '->' x}{http://rocq-prover.org/doc/V9.1.1/stdlib//Corelib.Init.Logic}{\coqdocnotation{)}} \coqexternalref{::type scope:x '->' x}{http://rocq-prover.org/doc/V9.1.1/stdlib//Corelib.Init.Logic}{\coqdocnotation{\ensuremath{\rightarrow}}} \coqexternalref{::type scope:x '->' x}{http://rocq-prover.org/doc/V9.1.1/stdlib//Corelib.Init.Logic}{\coqdocnotation{(}}\coqref{SECF.Imp.a2:130}{\coqdocvariable{a2}} \coqref{SECF.Imp.aevalR division.::type scope:x '==>' x}{\coqdocnotation{==>}} \coqref{SECF.Imp.n2:132}{\coqdocvariable{n2}}\coqexternalref{::type scope:x '->' x}{http://rocq-prover.org/doc/V9.1.1/stdlib//Corelib.Init.Logic}{\coqdocnotation{)}} \coqexternalref{::type scope:x '->' x}{http://rocq-prover.org/doc/V9.1.1/stdlib//Corelib.Init.Logic}{\coqdocnotation{\ensuremath{\rightarrow}}} \coqref{SECF.Imp.aevalR division.::type scope:x '==>' x}{\coqdocnotation{(}}\coqref{SECF.Imp.aevalR division.APlus}{\coqdocconstructor{APlus}} \coqref{SECF.Imp.a1:129}{\coqdocvariable{a1}} \coqref{SECF.Imp.a2:130}{\coqdocvariable{a2}}\coqref{SECF.Imp.aevalR division.::type scope:x '==>' x}{\coqdocnotation{)}} \coqref{SECF.Imp.aevalR division.::type scope:x '==>' x}{\coqdocnotation{==>}} \coqref{SECF.Imp.aevalR division.::type scope:x '==>' x}{\coqdocnotation{(}}\coqref{SECF.Imp.n1:131}{\coqdocvariable{n1}} \coqexternalref{::nat scope:x '+' x}{http://rocq-prover.org/doc/V9.1.1/stdlib//Corelib.Init.Nat}{\coqdocnotation{+}} \coqref{SECF.Imp.n2:132}{\coqdocvariable{n2}}\coqref{SECF.Imp.aevalR division.::type scope:x '==>' x}{\coqdocnotation{)}}\coqdoceol
\coqdocindent{1.00em}
\ensuremath{|} \coqdef{SECF.Imp.aevalR division.E AMinus}{E\_AMinus}{\coqdocconstructor{E\_AMinus}} (\coqdef{SECF.Imp.a1:133}{a1}{\coqdocbinder{a1}} \coqdef{SECF.Imp.a2:134}{a2}{\coqdocbinder{a2}} : \coqref{SECF.Imp.aevalR division.aexp}{\coqdocinductive{aexp}}) (\coqdef{SECF.Imp.n1:135}{n1}{\coqdocbinder{n1}} \coqdef{SECF.Imp.n2:136}{n2}{\coqdocbinder{n2}} : \coqexternalref{nat}{http://rocq-prover.org/doc/V9.1.1/stdlib//Corelib.Init.Datatypes}{\coqdocinductive{nat}}) :\coqdoceol
\coqdocindent{3.00em}
\coqexternalref{::type scope:x '->' x}{http://rocq-prover.org/doc/V9.1.1/stdlib//Corelib.Init.Logic}{\coqdocnotation{(}}\coqref{SECF.Imp.a1:133}{\coqdocvariable{a1}} \coqref{SECF.Imp.aevalR division.::type scope:x '==>' x}{\coqdocnotation{==>}} \coqref{SECF.Imp.n1:135}{\coqdocvariable{n1}}\coqexternalref{::type scope:x '->' x}{http://rocq-prover.org/doc/V9.1.1/stdlib//Corelib.Init.Logic}{\coqdocnotation{)}} \coqexternalref{::type scope:x '->' x}{http://rocq-prover.org/doc/V9.1.1/stdlib//Corelib.Init.Logic}{\coqdocnotation{\ensuremath{\rightarrow}}} \coqexternalref{::type scope:x '->' x}{http://rocq-prover.org/doc/V9.1.1/stdlib//Corelib.Init.Logic}{\coqdocnotation{(}}\coqref{SECF.Imp.a2:134}{\coqdocvariable{a2}} \coqref{SECF.Imp.aevalR division.::type scope:x '==>' x}{\coqdocnotation{==>}} \coqref{SECF.Imp.n2:136}{\coqdocvariable{n2}}\coqexternalref{::type scope:x '->' x}{http://rocq-prover.org/doc/V9.1.1/stdlib//Corelib.Init.Logic}{\coqdocnotation{)}} \coqexternalref{::type scope:x '->' x}{http://rocq-prover.org/doc/V9.1.1/stdlib//Corelib.Init.Logic}{\coqdocnotation{\ensuremath{\rightarrow}}} \coqref{SECF.Imp.aevalR division.::type scope:x '==>' x}{\coqdocnotation{(}}\coqref{SECF.Imp.aevalR division.AMinus}{\coqdocconstructor{AMinus}} \coqref{SECF.Imp.a1:133}{\coqdocvariable{a1}} \coqref{SECF.Imp.a2:134}{\coqdocvariable{a2}}\coqref{SECF.Imp.aevalR division.::type scope:x '==>' x}{\coqdocnotation{)}} \coqref{SECF.Imp.aevalR division.::type scope:x '==>' x}{\coqdocnotation{==>}} \coqref{SECF.Imp.aevalR division.::type scope:x '==>' x}{\coqdocnotation{(}}\coqref{SECF.Imp.n1:135}{\coqdocvariable{n1}} \coqexternalref{::nat scope:x '-' x}{http://rocq-prover.org/doc/V9.1.1/stdlib//Corelib.Init.Nat}{\coqdocnotation{-}} \coqref{SECF.Imp.n2:136}{\coqdocvariable{n2}}\coqref{SECF.Imp.aevalR division.::type scope:x '==>' x}{\coqdocnotation{)}}\coqdoceol
\coqdocindent{1.00em}
\ensuremath{|} \coqdef{SECF.Imp.aevalR division.E AMult}{E\_AMult}{\coqdocconstructor{E\_AMult}} (\coqdef{SECF.Imp.a1:137}{a1}{\coqdocbinder{a1}} \coqdef{SECF.Imp.a2:138}{a2}{\coqdocbinder{a2}} : \coqref{SECF.Imp.aevalR division.aexp}{\coqdocinductive{aexp}}) (\coqdef{SECF.Imp.n1:139}{n1}{\coqdocbinder{n1}} \coqdef{SECF.Imp.n2:140}{n2}{\coqdocbinder{n2}} : \coqexternalref{nat}{http://rocq-prover.org/doc/V9.1.1/stdlib//Corelib.Init.Datatypes}{\coqdocinductive{nat}}) :\coqdoceol
\coqdocindent{3.00em}
\coqexternalref{::type scope:x '->' x}{http://rocq-prover.org/doc/V9.1.1/stdlib//Corelib.Init.Logic}{\coqdocnotation{(}}\coqref{SECF.Imp.a1:137}{\coqdocvariable{a1}} \coqref{SECF.Imp.aevalR division.::type scope:x '==>' x}{\coqdocnotation{==>}} \coqref{SECF.Imp.n1:139}{\coqdocvariable{n1}}\coqexternalref{::type scope:x '->' x}{http://rocq-prover.org/doc/V9.1.1/stdlib//Corelib.Init.Logic}{\coqdocnotation{)}} \coqexternalref{::type scope:x '->' x}{http://rocq-prover.org/doc/V9.1.1/stdlib//Corelib.Init.Logic}{\coqdocnotation{\ensuremath{\rightarrow}}} \coqexternalref{::type scope:x '->' x}{http://rocq-prover.org/doc/V9.1.1/stdlib//Corelib.Init.Logic}{\coqdocnotation{(}}\coqref{SECF.Imp.a2:138}{\coqdocvariable{a2}} \coqref{SECF.Imp.aevalR division.::type scope:x '==>' x}{\coqdocnotation{==>}} \coqref{SECF.Imp.n2:140}{\coqdocvariable{n2}}\coqexternalref{::type scope:x '->' x}{http://rocq-prover.org/doc/V9.1.1/stdlib//Corelib.Init.Logic}{\coqdocnotation{)}} \coqexternalref{::type scope:x '->' x}{http://rocq-prover.org/doc/V9.1.1/stdlib//Corelib.Init.Logic}{\coqdocnotation{\ensuremath{\rightarrow}}} \coqref{SECF.Imp.aevalR division.::type scope:x '==>' x}{\coqdocnotation{(}}\coqref{SECF.Imp.aevalR division.AMult}{\coqdocconstructor{AMult}} \coqref{SECF.Imp.a1:137}{\coqdocvariable{a1}} \coqref{SECF.Imp.a2:138}{\coqdocvariable{a2}}\coqref{SECF.Imp.aevalR division.::type scope:x '==>' x}{\coqdocnotation{)}} \coqref{SECF.Imp.aevalR division.::type scope:x '==>' x}{\coqdocnotation{==>}} \coqref{SECF.Imp.aevalR division.::type scope:x '==>' x}{\coqdocnotation{(}}\coqref{SECF.Imp.n1:139}{\coqdocvariable{n1}} \coqexternalref{::nat scope:x '*' x}{http://rocq-prover.org/doc/V9.1.1/stdlib//Corelib.Init.Nat}{\coqdocnotation{\ensuremath{\times}}} \coqref{SECF.Imp.n2:140}{\coqdocvariable{n2}}\coqref{SECF.Imp.aevalR division.::type scope:x '==>' x}{\coqdocnotation{)}}\coqdoceol
\coqdocindent{1.00em}
\ensuremath{|} \coqdef{SECF.Imp.aevalR division.E ADiv}{E\_ADiv}{\coqdocconstructor{E\_ADiv}} (\coqdef{SECF.Imp.a1:141}{a1}{\coqdocbinder{a1}} \coqdef{SECF.Imp.a2:142}{a2}{\coqdocbinder{a2}} : \coqref{SECF.Imp.aevalR division.aexp}{\coqdocinductive{aexp}}) (\coqdef{SECF.Imp.n1:143}{n1}{\coqdocbinder{n1}} \coqdef{SECF.Imp.n2:144}{n2}{\coqdocbinder{n2}} \coqdef{SECF.Imp.n3:145}{n3}{\coqdocbinder{n3}} : \coqexternalref{nat}{http://rocq-prover.org/doc/V9.1.1/stdlib//Corelib.Init.Datatypes}{\coqdocinductive{nat}}) :          \coqdoceol
\coqdocindent{3.00em}
\coqexternalref{::type scope:x '->' x}{http://rocq-prover.org/doc/V9.1.1/stdlib//Corelib.Init.Logic}{\coqdocnotation{(}}\coqref{SECF.Imp.a1:141}{\coqdocvariable{a1}} \coqref{SECF.Imp.aevalR division.::type scope:x '==>' x}{\coqdocnotation{==>}} \coqref{SECF.Imp.n1:143}{\coqdocvariable{n1}}\coqexternalref{::type scope:x '->' x}{http://rocq-prover.org/doc/V9.1.1/stdlib//Corelib.Init.Logic}{\coqdocnotation{)}} \coqexternalref{::type scope:x '->' x}{http://rocq-prover.org/doc/V9.1.1/stdlib//Corelib.Init.Logic}{\coqdocnotation{\ensuremath{\rightarrow}}} \coqexternalref{::type scope:x '->' x}{http://rocq-prover.org/doc/V9.1.1/stdlib//Corelib.Init.Logic}{\coqdocnotation{(}}\coqref{SECF.Imp.a2:142}{\coqdocvariable{a2}} \coqref{SECF.Imp.aevalR division.::type scope:x '==>' x}{\coqdocnotation{==>}} \coqref{SECF.Imp.n2:144}{\coqdocvariable{n2}}\coqexternalref{::type scope:x '->' x}{http://rocq-prover.org/doc/V9.1.1/stdlib//Corelib.Init.Logic}{\coqdocnotation{)}} \coqexternalref{::type scope:x '->' x}{http://rocq-prover.org/doc/V9.1.1/stdlib//Corelib.Init.Logic}{\coqdocnotation{\ensuremath{\rightarrow}}} \coqexternalref{::type scope:x '->' x}{http://rocq-prover.org/doc/V9.1.1/stdlib//Corelib.Init.Logic}{\coqdocnotation{(}}\coqref{SECF.Imp.n2:144}{\coqdocvariable{n2}} \coqexternalref{::nat scope:x '>' x}{http://rocq-prover.org/doc/V9.1.1/stdlib//Corelib.Init.Peano}{\coqdocnotation{>}} 0\coqexternalref{::type scope:x '->' x}{http://rocq-prover.org/doc/V9.1.1/stdlib//Corelib.Init.Logic}{\coqdocnotation{)}} \coqexternalref{::type scope:x '->' x}{http://rocq-prover.org/doc/V9.1.1/stdlib//Corelib.Init.Logic}{\coqdocnotation{\ensuremath{\rightarrow}}}\coqdoceol
\coqdocindent{3.00em}
\coqexternalref{::type scope:x '->' x}{http://rocq-prover.org/doc/V9.1.1/stdlib//Corelib.Init.Logic}{\coqdocnotation{(}}\coqexternalref{mult}{http://rocq-prover.org/doc/V9.1.1/stdlib//Corelib.Init.Peano}{\coqdocabbreviation{mult}} \coqref{SECF.Imp.n2:144}{\coqdocvariable{n2}} \coqref{SECF.Imp.n3:145}{\coqdocvariable{n3}} \coqexternalref{::type scope:x '=' x}{http://rocq-prover.org/doc/V9.1.1/stdlib//Corelib.Init.Logic}{\coqdocnotation{=}} \coqref{SECF.Imp.n1:143}{\coqdocvariable{n1}}\coqexternalref{::type scope:x '->' x}{http://rocq-prover.org/doc/V9.1.1/stdlib//Corelib.Init.Logic}{\coqdocnotation{)}} \coqexternalref{::type scope:x '->' x}{http://rocq-prover.org/doc/V9.1.1/stdlib//Corelib.Init.Logic}{\coqdocnotation{\ensuremath{\rightarrow}}} \coqref{SECF.Imp.aevalR division.::type scope:x '==>' x}{\coqdocnotation{(}}\coqref{SECF.Imp.aevalR division.ADiv}{\coqdocconstructor{ADiv}} \coqref{SECF.Imp.a1:141}{\coqdocvariable{a1}} \coqref{SECF.Imp.a2:142}{\coqdocvariable{a2}}\coqref{SECF.Imp.aevalR division.::type scope:x '==>' x}{\coqdocnotation{)}} \coqref{SECF.Imp.aevalR division.::type scope:x '==>' x}{\coqdocnotation{==>}} \coqref{SECF.Imp.n3:145}{\coqdocvariable{n3}}\coqdoceol
\coqdocnoindent
\coqdoceol
\coqdocnoindent
\coqdockw{where} \coqdef{SECF.Imp.aevalR division.::type scope:x '==>' x}{"}{"}a '==>' n" := (\coqref{SECF.Imp.aevalR:127}{\coqdocinductive{aevalR}} \coqdocvar{a} \coqdocvar{n}) : \coqdocvar{type\_scope}.\coqdoceol
\coqdocemptyline
\end{coqdoccode}
Notice that this evaluation relation corresponds to a \textit{partial}
    function: There are some inputs for which it does not specify an
    output. 
\begin{coqdoccode}
\coqdocemptyline
\coqdocnoindent
\coqdockw{End} \coqref{SECF.Imp.aevalR division}{\coqdocmodule{aevalR\_division}}.\coqdoceol
\coqdocemptyline
\coqdocnoindent
\coqdockw{Module} \coqdef{SECF.Imp.aevalR extended}{aevalR\_extended}{\coqdocmodule{aevalR\_extended}}.\coqdoceol
\coqdocemptyline
\end{coqdoccode}
Or suppose that we want to extend the arithmetic operations
    by a nondeterministic number generator \coqdocvar{any} that, when evaluated,
    may yield any number.


    (Note that this is not the same as making a \textit{probabilistic} choice
    among all possible numbers -- we're not specifying any particular
    probability distribution for the results, just saying what results
    are \textit{possible}.) 
\begin{coqdoccode}
\coqdocemptyline
\coqdocnoindent
\coqdockw{Reserved Notation} "e '==>' n" (\coqdoctac{at} \coqdockw{level} 90, \coqdoctac{left} \coqdockw{associativity}).\coqdoceol
\coqdocemptyline
\coqdocnoindent
\coqdockw{Inductive} \coqdef{SECF.Imp.aevalR extended.aexp}{aexp}{\coqdocinductive{aexp}} : \coqdockw{Type} :=\coqdoceol
\coqdocindent{1.00em}
\ensuremath{|} \coqdef{SECF.Imp.aevalR extended.AAny}{AAny}{\coqdocconstructor{AAny}}                           \coqdoceol
\coqdocindent{1.00em}
\ensuremath{|} \coqdef{SECF.Imp.aevalR extended.ANum}{ANum}{\coqdocconstructor{ANum}} (\coqdef{SECF.Imp.n:148}{n}{\coqdocbinder{n}} : \coqexternalref{nat}{http://rocq-prover.org/doc/V9.1.1/stdlib//Corelib.Init.Datatypes}{\coqdocinductive{nat}})\coqdoceol
\coqdocindent{1.00em}
\ensuremath{|} \coqdef{SECF.Imp.aevalR extended.APlus}{APlus}{\coqdocconstructor{APlus}} (\coqdef{SECF.Imp.a1:149}{a1}{\coqdocbinder{a1}} \coqdef{SECF.Imp.a2:150}{a2}{\coqdocbinder{a2}} : \coqref{SECF.Imp.aexp:146}{\coqdocinductive{aexp}})\coqdoceol
\coqdocindent{1.00em}
\ensuremath{|} \coqdef{SECF.Imp.aevalR extended.AMinus}{AMinus}{\coqdocconstructor{AMinus}} (\coqdef{SECF.Imp.a1:151}{a1}{\coqdocbinder{a1}} \coqdef{SECF.Imp.a2:152}{a2}{\coqdocbinder{a2}} : \coqref{SECF.Imp.aexp:146}{\coqdocinductive{aexp}})\coqdoceol
\coqdocindent{1.00em}
\ensuremath{|} \coqdef{SECF.Imp.aevalR extended.AMult}{AMult}{\coqdocconstructor{AMult}} (\coqdef{SECF.Imp.a1:153}{a1}{\coqdocbinder{a1}} \coqdef{SECF.Imp.a2:154}{a2}{\coqdocbinder{a2}} : \coqref{SECF.Imp.aexp:146}{\coqdocinductive{aexp}}).\coqdoceol
\coqdocemptyline
\end{coqdoccode}
Again, extending \coqref{SECF.Imp.aeval}{\coqdocdefinition{aeval}} would be tricky, since now
    evaluation is \textit{not} a deterministic function from expressions to
    numbers; but extending \coqref{SECF.Imp.aevalR extended.aevalR}{\coqdocinductive{aevalR}} is no problem... 
\begin{coqdoccode}
\coqdocemptyline
\coqdocnoindent
\coqdockw{Inductive} \coqdef{SECF.Imp.aevalR extended.aevalR}{aevalR}{\coqdocinductive{aevalR}} : \coqref{SECF.Imp.aevalR extended.aexp}{\coqdocinductive{aexp}} \coqexternalref{::type scope:x '->' x}{http://rocq-prover.org/doc/V9.1.1/stdlib//Corelib.Init.Logic}{\coqdocnotation{\ensuremath{\rightarrow}}} \coqexternalref{nat}{http://rocq-prover.org/doc/V9.1.1/stdlib//Corelib.Init.Datatypes}{\coqdocinductive{nat}} \coqexternalref{::type scope:x '->' x}{http://rocq-prover.org/doc/V9.1.1/stdlib//Corelib.Init.Logic}{\coqdocnotation{\ensuremath{\rightarrow}}} \coqdockw{Prop} :=\coqdoceol
\coqdocindent{1.00em}
\ensuremath{|} \coqdef{SECF.Imp.aevalR extended.E Any}{E\_Any}{\coqdocconstructor{E\_Any}} (\coqdef{SECF.Imp.n:157}{n}{\coqdocbinder{n}} : \coqexternalref{nat}{http://rocq-prover.org/doc/V9.1.1/stdlib//Corelib.Init.Datatypes}{\coqdocinductive{nat}}) :\coqdoceol
\coqdocindent{3.00em}
\coqref{SECF.Imp.aevalR extended.AAny}{\coqdocconstructor{AAny}} \coqref{SECF.Imp.aevalR extended.::type scope:x '==>' x}{\coqdocnotation{==>}} \coqref{SECF.Imp.n:157}{\coqdocvariable{n}}                        \coqdoceol
\coqdocindent{1.00em}
\ensuremath{|} \coqdef{SECF.Imp.aevalR extended.E ANum}{E\_ANum}{\coqdocconstructor{E\_ANum}} (\coqdef{SECF.Imp.n:158}{n}{\coqdocbinder{n}} : \coqexternalref{nat}{http://rocq-prover.org/doc/V9.1.1/stdlib//Corelib.Init.Datatypes}{\coqdocinductive{nat}}) :\coqdoceol
\coqdocindent{3.00em}
\coqref{SECF.Imp.aevalR extended.::type scope:x '==>' x}{\coqdocnotation{(}}\coqref{SECF.Imp.aevalR extended.ANum}{\coqdocconstructor{ANum}} \coqref{SECF.Imp.n:158}{\coqdocvariable{n}}\coqref{SECF.Imp.aevalR extended.::type scope:x '==>' x}{\coqdocnotation{)}} \coqref{SECF.Imp.aevalR extended.::type scope:x '==>' x}{\coqdocnotation{==>}} \coqref{SECF.Imp.n:158}{\coqdocvariable{n}}\coqdoceol
\coqdocindent{1.00em}
\ensuremath{|} \coqdef{SECF.Imp.aevalR extended.E APlus}{E\_APlus}{\coqdocconstructor{E\_APlus}} (\coqdef{SECF.Imp.a1:159}{a1}{\coqdocbinder{a1}} \coqdef{SECF.Imp.a2:160}{a2}{\coqdocbinder{a2}} : \coqref{SECF.Imp.aevalR extended.aexp}{\coqdocinductive{aexp}}) (\coqdef{SECF.Imp.n1:161}{n1}{\coqdocbinder{n1}} \coqdef{SECF.Imp.n2:162}{n2}{\coqdocbinder{n2}} : \coqexternalref{nat}{http://rocq-prover.org/doc/V9.1.1/stdlib//Corelib.Init.Datatypes}{\coqdocinductive{nat}}) :\coqdoceol
\coqdocindent{3.00em}
\coqexternalref{::type scope:x '->' x}{http://rocq-prover.org/doc/V9.1.1/stdlib//Corelib.Init.Logic}{\coqdocnotation{(}}\coqref{SECF.Imp.a1:159}{\coqdocvariable{a1}} \coqref{SECF.Imp.aevalR extended.::type scope:x '==>' x}{\coqdocnotation{==>}} \coqref{SECF.Imp.n1:161}{\coqdocvariable{n1}}\coqexternalref{::type scope:x '->' x}{http://rocq-prover.org/doc/V9.1.1/stdlib//Corelib.Init.Logic}{\coqdocnotation{)}} \coqexternalref{::type scope:x '->' x}{http://rocq-prover.org/doc/V9.1.1/stdlib//Corelib.Init.Logic}{\coqdocnotation{\ensuremath{\rightarrow}}} \coqexternalref{::type scope:x '->' x}{http://rocq-prover.org/doc/V9.1.1/stdlib//Corelib.Init.Logic}{\coqdocnotation{(}}\coqref{SECF.Imp.a2:160}{\coqdocvariable{a2}} \coqref{SECF.Imp.aevalR extended.::type scope:x '==>' x}{\coqdocnotation{==>}} \coqref{SECF.Imp.n2:162}{\coqdocvariable{n2}}\coqexternalref{::type scope:x '->' x}{http://rocq-prover.org/doc/V9.1.1/stdlib//Corelib.Init.Logic}{\coqdocnotation{)}} \coqexternalref{::type scope:x '->' x}{http://rocq-prover.org/doc/V9.1.1/stdlib//Corelib.Init.Logic}{\coqdocnotation{\ensuremath{\rightarrow}}} \coqref{SECF.Imp.aevalR extended.::type scope:x '==>' x}{\coqdocnotation{(}}\coqref{SECF.Imp.aevalR extended.APlus}{\coqdocconstructor{APlus}} \coqref{SECF.Imp.a1:159}{\coqdocvariable{a1}} \coqref{SECF.Imp.a2:160}{\coqdocvariable{a2}}\coqref{SECF.Imp.aevalR extended.::type scope:x '==>' x}{\coqdocnotation{)}} \coqref{SECF.Imp.aevalR extended.::type scope:x '==>' x}{\coqdocnotation{==>}} \coqref{SECF.Imp.aevalR extended.::type scope:x '==>' x}{\coqdocnotation{(}}\coqref{SECF.Imp.n1:161}{\coqdocvariable{n1}} \coqexternalref{::nat scope:x '+' x}{http://rocq-prover.org/doc/V9.1.1/stdlib//Corelib.Init.Nat}{\coqdocnotation{+}} \coqref{SECF.Imp.n2:162}{\coqdocvariable{n2}}\coqref{SECF.Imp.aevalR extended.::type scope:x '==>' x}{\coqdocnotation{)}}\coqdoceol
\coqdocindent{1.00em}
\ensuremath{|} \coqdef{SECF.Imp.aevalR extended.E AMinus}{E\_AMinus}{\coqdocconstructor{E\_AMinus}} (\coqdef{SECF.Imp.a1:163}{a1}{\coqdocbinder{a1}} \coqdef{SECF.Imp.a2:164}{a2}{\coqdocbinder{a2}} : \coqref{SECF.Imp.aevalR extended.aexp}{\coqdocinductive{aexp}}) (\coqdef{SECF.Imp.n1:165}{n1}{\coqdocbinder{n1}} \coqdef{SECF.Imp.n2:166}{n2}{\coqdocbinder{n2}} : \coqexternalref{nat}{http://rocq-prover.org/doc/V9.1.1/stdlib//Corelib.Init.Datatypes}{\coqdocinductive{nat}}) :\coqdoceol
\coqdocindent{3.00em}
\coqexternalref{::type scope:x '->' x}{http://rocq-prover.org/doc/V9.1.1/stdlib//Corelib.Init.Logic}{\coqdocnotation{(}}\coqref{SECF.Imp.a1:163}{\coqdocvariable{a1}} \coqref{SECF.Imp.aevalR extended.::type scope:x '==>' x}{\coqdocnotation{==>}} \coqref{SECF.Imp.n1:165}{\coqdocvariable{n1}}\coqexternalref{::type scope:x '->' x}{http://rocq-prover.org/doc/V9.1.1/stdlib//Corelib.Init.Logic}{\coqdocnotation{)}} \coqexternalref{::type scope:x '->' x}{http://rocq-prover.org/doc/V9.1.1/stdlib//Corelib.Init.Logic}{\coqdocnotation{\ensuremath{\rightarrow}}} \coqexternalref{::type scope:x '->' x}{http://rocq-prover.org/doc/V9.1.1/stdlib//Corelib.Init.Logic}{\coqdocnotation{(}}\coqref{SECF.Imp.a2:164}{\coqdocvariable{a2}} \coqref{SECF.Imp.aevalR extended.::type scope:x '==>' x}{\coqdocnotation{==>}} \coqref{SECF.Imp.n2:166}{\coqdocvariable{n2}}\coqexternalref{::type scope:x '->' x}{http://rocq-prover.org/doc/V9.1.1/stdlib//Corelib.Init.Logic}{\coqdocnotation{)}} \coqexternalref{::type scope:x '->' x}{http://rocq-prover.org/doc/V9.1.1/stdlib//Corelib.Init.Logic}{\coqdocnotation{\ensuremath{\rightarrow}}} \coqref{SECF.Imp.aevalR extended.::type scope:x '==>' x}{\coqdocnotation{(}}\coqref{SECF.Imp.aevalR extended.AMinus}{\coqdocconstructor{AMinus}} \coqref{SECF.Imp.a1:163}{\coqdocvariable{a1}} \coqref{SECF.Imp.a2:164}{\coqdocvariable{a2}}\coqref{SECF.Imp.aevalR extended.::type scope:x '==>' x}{\coqdocnotation{)}} \coqref{SECF.Imp.aevalR extended.::type scope:x '==>' x}{\coqdocnotation{==>}} \coqref{SECF.Imp.aevalR extended.::type scope:x '==>' x}{\coqdocnotation{(}}\coqref{SECF.Imp.n1:165}{\coqdocvariable{n1}} \coqexternalref{::nat scope:x '-' x}{http://rocq-prover.org/doc/V9.1.1/stdlib//Corelib.Init.Nat}{\coqdocnotation{-}} \coqref{SECF.Imp.n2:166}{\coqdocvariable{n2}}\coqref{SECF.Imp.aevalR extended.::type scope:x '==>' x}{\coqdocnotation{)}}\coqdoceol
\coqdocindent{1.00em}
\ensuremath{|} \coqdef{SECF.Imp.aevalR extended.E AMult}{E\_AMult}{\coqdocconstructor{E\_AMult}} (\coqdef{SECF.Imp.a1:167}{a1}{\coqdocbinder{a1}} \coqdef{SECF.Imp.a2:168}{a2}{\coqdocbinder{a2}} : \coqref{SECF.Imp.aevalR extended.aexp}{\coqdocinductive{aexp}}) (\coqdef{SECF.Imp.n1:169}{n1}{\coqdocbinder{n1}} \coqdef{SECF.Imp.n2:170}{n2}{\coqdocbinder{n2}} : \coqexternalref{nat}{http://rocq-prover.org/doc/V9.1.1/stdlib//Corelib.Init.Datatypes}{\coqdocinductive{nat}}) :\coqdoceol
\coqdocindent{3.00em}
\coqexternalref{::type scope:x '->' x}{http://rocq-prover.org/doc/V9.1.1/stdlib//Corelib.Init.Logic}{\coqdocnotation{(}}\coqref{SECF.Imp.a1:167}{\coqdocvariable{a1}} \coqref{SECF.Imp.aevalR extended.::type scope:x '==>' x}{\coqdocnotation{==>}} \coqref{SECF.Imp.n1:169}{\coqdocvariable{n1}}\coqexternalref{::type scope:x '->' x}{http://rocq-prover.org/doc/V9.1.1/stdlib//Corelib.Init.Logic}{\coqdocnotation{)}} \coqexternalref{::type scope:x '->' x}{http://rocq-prover.org/doc/V9.1.1/stdlib//Corelib.Init.Logic}{\coqdocnotation{\ensuremath{\rightarrow}}} \coqexternalref{::type scope:x '->' x}{http://rocq-prover.org/doc/V9.1.1/stdlib//Corelib.Init.Logic}{\coqdocnotation{(}}\coqref{SECF.Imp.a2:168}{\coqdocvariable{a2}} \coqref{SECF.Imp.aevalR extended.::type scope:x '==>' x}{\coqdocnotation{==>}} \coqref{SECF.Imp.n2:170}{\coqdocvariable{n2}}\coqexternalref{::type scope:x '->' x}{http://rocq-prover.org/doc/V9.1.1/stdlib//Corelib.Init.Logic}{\coqdocnotation{)}} \coqexternalref{::type scope:x '->' x}{http://rocq-prover.org/doc/V9.1.1/stdlib//Corelib.Init.Logic}{\coqdocnotation{\ensuremath{\rightarrow}}} \coqref{SECF.Imp.aevalR extended.::type scope:x '==>' x}{\coqdocnotation{(}}\coqref{SECF.Imp.aevalR extended.AMult}{\coqdocconstructor{AMult}} \coqref{SECF.Imp.a1:167}{\coqdocvariable{a1}} \coqref{SECF.Imp.a2:168}{\coqdocvariable{a2}}\coqref{SECF.Imp.aevalR extended.::type scope:x '==>' x}{\coqdocnotation{)}} \coqref{SECF.Imp.aevalR extended.::type scope:x '==>' x}{\coqdocnotation{==>}} \coqref{SECF.Imp.aevalR extended.::type scope:x '==>' x}{\coqdocnotation{(}}\coqref{SECF.Imp.n1:169}{\coqdocvariable{n1}} \coqexternalref{::nat scope:x '*' x}{http://rocq-prover.org/doc/V9.1.1/stdlib//Corelib.Init.Nat}{\coqdocnotation{\ensuremath{\times}}} \coqref{SECF.Imp.n2:170}{\coqdocvariable{n2}}\coqref{SECF.Imp.aevalR extended.::type scope:x '==>' x}{\coqdocnotation{)}}\coqdoceol
\coqdocnoindent
\coqdoceol
\coqdocnoindent
\coqdockw{where} \coqdef{SECF.Imp.aevalR extended.::type scope:x '==>' x}{"}{"}a '==>' n" := (\coqref{SECF.Imp.aevalR:156}{\coqdocinductive{aevalR}} \coqdocvar{a} \coqdocvar{n}) : \coqdocvar{type\_scope}.\coqdoceol
\coqdocemptyline
\coqdocnoindent
\coqdockw{End} \coqref{SECF.Imp.aevalR extended}{\coqdocmodule{aevalR\_extended}}.\coqdoceol
\coqdocemptyline
\end{coqdoccode}
At this point you maybe wondering: Which of these styles
    should I use by default?


    In the examples we've just seen, relational definitions turned out
    to be more useful than functional ones.  For situations like
    these, where the thing being defined is not easy to express as a
    function, or indeed where it is \textit{not} a function, there is no real
    choice.  But what about when both styles are workable?


    One point in favor of relational definitions is that they can be
    more elegant and easier to understand.


    Another is that Rocq automatically generates nice inversion and
    induction principles from \coqdockw{Inductive} definitions.


    On the other hand, functional definitions can often be more
    convenient:

\begin{itemize}
\item  Functions are automatically deterministic and total; for a
       relational definition, we have to \textit{prove} these properties
       explicitly if we need them.

\item  With functions we can also take advantage of Rocq's computation
       mechanism to simplify expressions during proofs.

\end{itemize}


    Furthermore, functions can be directly ``extracted'' from Gallina to
    executable code in OCaml or Haskell.


    Ultimately, the choice often comes down to either the specifics of
    a particular situation or simply a question of taste.  Indeed, in
    large Rocq developments it is common to see a definition given in
    \textit{both} functional and relational styles, plus a lemma stating that
    the two coincide, allowing further proofs to switch from one point
    of view to the other at will. 
\begin{coqdoccode}
\coqdocemptyline
\end{coqdoccode}
\section{Expressions With Variables}



 Let's return to defining Imp, where the next thing we need to do
    is to enrich our arithmetic and boolean expressions with
    variables.


    To keep things simple, we'll assume that all variables are global
    and that they only hold numbers. 
\begin{coqdoccode}
\coqdocemptyline
\end{coqdoccode}
\subsection{States}



 Since we'll want to look variables up to find out their current
    values, we'll use total maps from the \coqdoclibrary{Maps} chapter.


    A \textit{machine state} (or just \textit{state}) represents the current values
    of all variables at some point in the execution of a program. 

 For simplicity, we assume that the state is defined for
    \textit{all} variables, even though any given program is only able to
    mention a finite number of them.  Because each variable stores a
    natural number, we can represent the state as a total map from
    strings (variable names) to \coqexternalref{nat}{http://rocq-prover.org/doc/V9.1.1/stdlib//Corelib.Init.Datatypes}{\coqdocinductive{nat}}, and will use 0 as default
    value in the store. 
\begin{coqdoccode}
\coqdocemptyline
\coqdocnoindent
\coqdockw{Definition} \coqdef{SECF.Imp.state}{state}{\coqdocdefinition{state}} := \coqref{SECF.Maps.total map}{\coqdocdefinition{total\_map}} \coqexternalref{nat}{http://rocq-prover.org/doc/V9.1.1/stdlib//Corelib.Init.Datatypes}{\coqdocinductive{nat}}.\coqdoceol
\coqdocemptyline
\end{coqdoccode}
\subsection{Syntax}



 We can add variables to the arithmetic expressions we had before
    simply by including one more constructor: 
\begin{coqdoccode}
\coqdocemptyline
\coqdocnoindent
\coqdockw{Inductive} \coqdef{SECF.Imp.aexp}{aexp}{\coqdocinductive{aexp}} : \coqdockw{Type} :=\coqdoceol
\coqdocindent{1.00em}
\ensuremath{|} \coqdef{SECF.Imp.ANum}{ANum}{\coqdocconstructor{ANum}} (\coqdef{SECF.Imp.n:173}{n}{\coqdocbinder{n}} : \coqexternalref{nat}{http://rocq-prover.org/doc/V9.1.1/stdlib//Corelib.Init.Datatypes}{\coqdocinductive{nat}})\coqdoceol
\coqdocindent{1.00em}
\ensuremath{|} \coqdef{SECF.Imp.AId}{AId}{\coqdocconstructor{AId}} (\coqdef{SECF.Imp.x:174}{x}{\coqdocbinder{x}} : \coqdocinductive{string})              \coqdoceol
\coqdocindent{1.00em}
\ensuremath{|} \coqdef{SECF.Imp.APlus}{APlus}{\coqdocconstructor{APlus}} (\coqdef{SECF.Imp.a1:175}{a1}{\coqdocbinder{a1}} \coqdef{SECF.Imp.a2:176}{a2}{\coqdocbinder{a2}} : \coqref{SECF.Imp.aexp:171}{\coqdocinductive{aexp}})\coqdoceol
\coqdocindent{1.00em}
\ensuremath{|} \coqdef{SECF.Imp.AMinus}{AMinus}{\coqdocconstructor{AMinus}} (\coqdef{SECF.Imp.a1:177}{a1}{\coqdocbinder{a1}} \coqdef{SECF.Imp.a2:178}{a2}{\coqdocbinder{a2}} : \coqref{SECF.Imp.aexp:171}{\coqdocinductive{aexp}})\coqdoceol
\coqdocindent{1.00em}
\ensuremath{|} \coqdef{SECF.Imp.AMult}{AMult}{\coqdocconstructor{AMult}} (\coqdef{SECF.Imp.a1:179}{a1}{\coqdocbinder{a1}} \coqdef{SECF.Imp.a2:180}{a2}{\coqdocbinder{a2}} : \coqref{SECF.Imp.aexp:171}{\coqdocinductive{aexp}}).\coqdoceol
\coqdocemptyline
\end{coqdoccode}
Defining a few variable names as notational shorthands will make
    examples easier to read: 
\begin{coqdoccode}
\coqdocemptyline
\coqdocnoindent
\coqdockw{Definition} \coqdef{SECF.Imp.W}{W}{\coqdocdefinition{W}} : \coqdocinductive{string} := "W".\coqdoceol
\coqdocnoindent
\coqdockw{Definition} \coqdef{SECF.Imp.X}{X}{\coqdocdefinition{X}} : \coqdocinductive{string} := "X".\coqdoceol
\coqdocnoindent
\coqdockw{Definition} \coqdef{SECF.Imp.Y}{Y}{\coqdocdefinition{Y}} : \coqdocinductive{string} := "Y".\coqdoceol
\coqdocnoindent
\coqdockw{Definition} \coqdef{SECF.Imp.Z}{Z}{\coqdocdefinition{Z}} : \coqdocinductive{string} := "Z".\coqdoceol
\coqdocemptyline
\end{coqdoccode}
(This convention for naming program variables (\coqref{SECF.SpecCT.X}{\coqdocdefinition{X}}, \coqref{SECF.SpecCT.Y}{\coqdocdefinition{Y}},
    \coqref{SECF.SpecCT.Z}{\coqdocdefinition{Z}}) clashes a bit with our earlier use of uppercase letters for
    types.  Since we're not using polymorphism heavily in the chapters
    developed to Imp, this overloading should not cause confusion.) 

 The definition of \coqref{SECF.Imp.bexp}{\coqdocinductive{bexp}}s is unchanged (except that it now refers
    to the new \coqref{SECF.Imp.aexp}{\coqdocinductive{aexp}}s): 
\begin{coqdoccode}
\coqdocemptyline
\coqdocnoindent
\coqdockw{Inductive} \coqdef{SECF.Imp.bexp}{bexp}{\coqdocinductive{bexp}} : \coqdockw{Type} :=\coqdoceol
\coqdocindent{1.00em}
\ensuremath{|} \coqdef{SECF.Imp.BTrue}{BTrue}{\coqdocconstructor{BTrue}}\coqdoceol
\coqdocindent{1.00em}
\ensuremath{|} \coqdef{SECF.Imp.BFalse}{BFalse}{\coqdocconstructor{BFalse}}\coqdoceol
\coqdocindent{1.00em}
\ensuremath{|} \coqdef{SECF.Imp.BEq}{BEq}{\coqdocconstructor{BEq}} (\coqdef{SECF.Imp.a1:183}{a1}{\coqdocbinder{a1}} \coqdef{SECF.Imp.a2:184}{a2}{\coqdocbinder{a2}} : \coqref{SECF.Imp.aexp}{\coqdocinductive{aexp}})\coqdoceol
\coqdocindent{1.00em}
\ensuremath{|} \coqdef{SECF.Imp.BNeq}{BNeq}{\coqdocconstructor{BNeq}} (\coqdef{SECF.Imp.a1:185}{a1}{\coqdocbinder{a1}} \coqdef{SECF.Imp.a2:186}{a2}{\coqdocbinder{a2}} : \coqref{SECF.Imp.aexp}{\coqdocinductive{aexp}})\coqdoceol
\coqdocindent{1.00em}
\ensuremath{|} \coqdef{SECF.Imp.BLe}{BLe}{\coqdocconstructor{BLe}} (\coqdef{SECF.Imp.a1:187}{a1}{\coqdocbinder{a1}} \coqdef{SECF.Imp.a2:188}{a2}{\coqdocbinder{a2}} : \coqref{SECF.Imp.aexp}{\coqdocinductive{aexp}})\coqdoceol
\coqdocindent{1.00em}
\ensuremath{|} \coqdef{SECF.Imp.BGt}{BGt}{\coqdocconstructor{BGt}} (\coqdef{SECF.Imp.a1:189}{a1}{\coqdocbinder{a1}} \coqdef{SECF.Imp.a2:190}{a2}{\coqdocbinder{a2}} : \coqref{SECF.Imp.aexp}{\coqdocinductive{aexp}})\coqdoceol
\coqdocindent{1.00em}
\ensuremath{|} \coqdef{SECF.Imp.BNot}{BNot}{\coqdocconstructor{BNot}} (\coqdef{SECF.Imp.b:191}{b}{\coqdocbinder{b}} : \coqref{SECF.Imp.bexp:181}{\coqdocinductive{bexp}})\coqdoceol
\coqdocindent{1.00em}
\ensuremath{|} \coqdef{SECF.Imp.BAnd}{BAnd}{\coqdocconstructor{BAnd}} (\coqdef{SECF.Imp.b1:192}{b1}{\coqdocbinder{b1}} \coqdef{SECF.Imp.b2:193}{b2}{\coqdocbinder{b2}} : \coqref{SECF.Imp.bexp:181}{\coqdocinductive{bexp}}).\coqdoceol
\coqdocemptyline
\end{coqdoccode}
\subsection{Notations}



 To make Imp programs easier to read and write, we introduce some
    notations and implicit coercions.  

 You do not need to understand exactly what these declarations do.


    Briefly, though:

\begin{itemize}
\item  The \coqdockw{Coercion} declaration stipulates that a function (or
         constructor) can be implicitly used by the type system to
         coerce a value of the input type to a value of the output
         type.  For instance, the coercion declaration for \coqref{SECF.SpecCT.AId}{\coqdocconstructor{AId}}
         allows us to use plain strings when an \coqref{SECF.Imp.aexp}{\coqdocinductive{aexp}} is expected;
         the string will implicitly be wrapped with \coqref{SECF.SpecCT.AId}{\coqdocconstructor{AId}}.

\item  \coqdocvar{Declare} \coqdocvar{Custom} \coqdocvar{Entry} \coqref{SECF.SpecCT.Div.com}{\coqdocinductive{com}} tells Rocq to create a new ``custom
         grammar'' for parsing Imp expressions and programs. The first
         notation declaration after this tells Rocq that anything
         between <\{ and \}> should be parsed using the Imp
         grammar. Again, it is not necessary to understand the
         details, but it is important to recognize that we are
         defining \textit{new} interpretations for some familiar operators
         like +, -, \ensuremath{\times}, =, \ensuremath{\le}, etc., when they occur between
         <\{ and \}>. 

\end{itemize}
\begin{coqdoccode}
\coqdocemptyline
\coqdocnoindent
\coqdockw{Coercion} \coqref{SECF.Imp.AId}{\coqdocconstructor{AId}} \coqref{SECF.Imp.AId}{\coqdocconstructor{:}} \coqref{SECF.Imp.AId}{\coqdocconstructor{string}} \coqref{SECF.Imp.AId}{\coqdocconstructor{>->}} \coqref{SECF.Imp.AId}{\coqdocconstructor{aexp}}.\coqdoceol
\coqdocnoindent
\coqdockw{Coercion} \coqref{SECF.Imp.ANum}{\coqdocconstructor{ANum}} \coqref{SECF.Imp.ANum}{\coqdocconstructor{:}} \coqref{SECF.Imp.ANum}{\coqdocconstructor{nat}} \coqref{SECF.Imp.ANum}{\coqdocconstructor{>->}} \coqref{SECF.Imp.ANum}{\coqdocconstructor{aexp}}.\coqdoceol
\coqdocemptyline
\coqdocnoindent
\coqdocvar{Declare} \coqdocvar{Custom} \coqdocvar{Entry} \coqdocvar{com}.\coqdoceol
\coqdocnoindent
\coqdockw{Declare Scope} \coqdocvar{com\_scope}.\coqdoceol
\coqdocemptyline
\coqdocnoindent
\coqdockw{Notation} \coqdef{SECF.Imp.::com scope:'<x7B' x 'x7D>'}{"}{"}<\{ e \}>" := \coqdocvar{e}\coqdoceol
\coqdocindent{1.00em}
(\coqdocvar{e} \coqdocvar{custom} \coqdocvar{com}, \coqdocvar{format} "'[hv' <\{ '/  ' '[v' e ']' '/' \}> ']'") : \coqdocvar{com\_scope}.\coqdoceol
\coqdocemptyline
\coqdocnoindent
\coqdockw{Notation} \coqdef{SECF.Imp.:com::'(' x ')'}{"}{"}( x )" := \coqdocvar{x} (\coqdoctac{in} \coqdocvar{custom} \coqdocvar{com}, \coqdocvar{x} \coqdoctac{at} \coqdockw{level} 99).\coqdoceol
\coqdocnoindent
\coqdockw{Notation} \coqdef{SECF.Imp.:com::x}{"}{"}x" := \coqdocvar{x} (\coqdoctac{in} \coqdocvar{custom} \coqdocvar{com} \coqdoctac{at} \coqdockw{level} 0, \coqdocvar{x} \coqdockw{constr} \coqdoctac{at} \coqdockw{level} 0).\coqdoceol
\coqdocemptyline
\coqdocnoindent
\coqdockw{Notation} \coqdef{SECF.Imp.:com::x x '..' x}{"}{"}f x .. y" := (.. (\coqdocvar{f} \coqdocvar{x}) .. \coqdocvar{y})\coqdoceol
\coqdocindent{9.00em}
(\coqdoctac{in} \coqdocvar{custom} \coqdocvar{com} \coqdoctac{at} \coqdockw{level} 0, \coqdocvar{only} \coqdocvar{parsing},\coqdoceol
\coqdocindent{9.00em}
\coqdocvar{f} \coqdockw{constr} \coqdoctac{at} \coqdockw{level} 0, \coqdocvar{x} \coqdockw{constr} \coqdoctac{at} \coqdockw{level} 1,\coqdoceol
\coqdocindent{11.00em}
\coqdocvar{y} \coqdockw{constr} \coqdoctac{at} \coqdockw{level} 1).\coqdoceol
\coqdocnoindent
\coqdockw{Notation} \coqdef{SECF.Imp.:com::x '+' x}{"}{"}x + y"   := (\coqref{SECF.Imp.APlus}{\coqdocconstructor{APlus}} \coqdocvar{x} \coqdocvar{y}) (\coqdoctac{in} \coqdocvar{custom} \coqdocvar{com} \coqdoctac{at} \coqdockw{level} 50, \coqdoctac{left} \coqdockw{associativity}).\coqdoceol
\coqdocnoindent
\coqdockw{Notation} \coqdef{SECF.Imp.:com::x '-' x}{"}{"}x - y"   := (\coqref{SECF.Imp.AMinus}{\coqdocconstructor{AMinus}} \coqdocvar{x} \coqdocvar{y}) (\coqdoctac{in} \coqdocvar{custom} \coqdocvar{com} \coqdoctac{at} \coqdockw{level} 50, \coqdoctac{left} \coqdockw{associativity}).\coqdoceol
\coqdocnoindent
\coqdockw{Notation} \coqdef{SECF.Imp.:com::x '*' x}{"}{"}x * y"   := (\coqref{SECF.Imp.AMult}{\coqdocconstructor{AMult}} \coqdocvar{x} \coqdocvar{y}) (\coqdoctac{in} \coqdocvar{custom} \coqdocvar{com} \coqdoctac{at} \coqdockw{level} 40, \coqdoctac{left} \coqdockw{associativity}).\coqdoceol
\coqdocnoindent
\coqdockw{Notation} \coqdef{SECF.Imp.:::'true'}{"}{"}'true'"  := \coqexternalref{true}{http://rocq-prover.org/doc/V9.1.1/stdlib//Corelib.Init.Datatypes}{\coqdocconstructor{true}} (\coqdoctac{at} \coqdockw{level} 1).\coqdoceol
\coqdocnoindent
\coqdockw{Notation} \coqdef{SECF.Imp.:com::'true'}{"}{"}'true'"  := \coqref{SECF.Imp.BTrue}{\coqdocconstructor{BTrue}} (\coqdoctac{in} \coqdocvar{custom} \coqdocvar{com} \coqdoctac{at} \coqdockw{level} 0).\coqdoceol
\coqdocnoindent
\coqdockw{Notation} \coqdef{SECF.Imp.:::'false'}{"}{"}'false'" := \coqexternalref{false}{http://rocq-prover.org/doc/V9.1.1/stdlib//Corelib.Init.Datatypes}{\coqdocconstructor{false}} (\coqdoctac{at} \coqdockw{level} 1).\coqdoceol
\coqdocnoindent
\coqdockw{Notation} \coqdef{SECF.Imp.:com::'false'}{"}{"}'false'" := \coqref{SECF.Imp.BFalse}{\coqdocconstructor{BFalse}} (\coqdoctac{in} \coqdocvar{custom} \coqdocvar{com} \coqdoctac{at} \coqdockw{level} 0).\coqdoceol
\coqdocnoindent
\coqdockw{Notation} \coqdef{SECF.Imp.:com::x '<=' x}{"}{"}x <= y"  := (\coqref{SECF.Imp.BLe}{\coqdocconstructor{BLe}} \coqdocvar{x} \coqdocvar{y}) (\coqdoctac{in} \coqdocvar{custom} \coqdocvar{com} \coqdoctac{at} \coqdockw{level} 70, \coqdockw{no} \coqdockw{associativity}).\coqdoceol
\coqdocnoindent
\coqdockw{Notation} \coqdef{SECF.Imp.:com::x '>' x}{"}{"}x > y"   := (\coqref{SECF.Imp.BGt}{\coqdocconstructor{BGt}} \coqdocvar{x} \coqdocvar{y}) (\coqdoctac{in} \coqdocvar{custom} \coqdocvar{com} \coqdoctac{at} \coqdockw{level} 70, \coqdockw{no} \coqdockw{associativity}).\coqdoceol
\coqdocnoindent
\coqdockw{Notation} \coqdef{SECF.Imp.:com::x '=' x}{"}{"}x = y"   := (\coqref{SECF.Imp.BEq}{\coqdocconstructor{BEq}} \coqdocvar{x} \coqdocvar{y}) (\coqdoctac{in} \coqdocvar{custom} \coqdocvar{com} \coqdoctac{at} \coqdockw{level} 70, \coqdockw{no} \coqdockw{associativity}).\coqdoceol
\coqdocnoindent
\coqdockw{Notation} \coqdef{SECF.Imp.:com::x '<>' x}{"}{"}x <> y"  := (\coqref{SECF.Imp.BNeq}{\coqdocconstructor{BNeq}} \coqdocvar{x} \coqdocvar{y}) (\coqdoctac{in} \coqdocvar{custom} \coqdocvar{com} \coqdoctac{at} \coqdockw{level} 70, \coqdockw{no} \coqdockw{associativity}).\coqdoceol
\coqdocnoindent
\coqdockw{Notation} \coqdef{SECF.Imp.:com::x 'x26x26' x}{"}{"}x \&\& y"  := (\coqref{SECF.Imp.BAnd}{\coqdocconstructor{BAnd}} \coqdocvar{x} \coqdocvar{y}) (\coqdoctac{in} \coqdocvar{custom} \coqdocvar{com} \coqdoctac{at} \coqdockw{level} 80, \coqdoctac{left} \coqdockw{associativity}).\coqdoceol
\coqdocnoindent
\coqdockw{Notation} \coqdef{SECF.Imp.:com::'x7E' x}{"}{"}'\~{}' b"   := (\coqref{SECF.Imp.BNot}{\coqdocconstructor{BNot}} \coqdocvar{b}) (\coqdoctac{in} \coqdocvar{custom} \coqdocvar{com} \coqdoctac{at} \coqdockw{level} 75, \coqdoctac{right} \coqdockw{associativity}).\coqdoceol
\coqdocemptyline
\coqdocnoindent
\coqdockw{Open} \coqdockw{Scope} \coqdocvar{com\_scope}.\coqdoceol
\coqdocemptyline
\end{coqdoccode}
We can now write 3 + (\coqref{SECF.SpecCT.X}{\coqdocdefinition{X}} \ensuremath{\times} 2) instead  of \coqref{SECF.SpecCT.APlus}{\coqdocdefinition{APlus}} 3 (\coqref{SECF.SpecCT.AMult}{\coqdocdefinition{AMult}} \coqref{SECF.SpecCT.X}{\coqdocdefinition{X}} 2),
    and \coqexternalref{true}{http://rocq-prover.org/doc/V9.1.1/stdlib//Corelib.Init.Datatypes}{\coqdocconstructor{true}} \&\& \~{}(\coqref{SECF.SpecCT.X}{\coqdocdefinition{X}} \ensuremath{\le} 4) instead of \coqref{SECF.SpecCT.BAnd}{\coqdocdefinition{BAnd}} \coqexternalref{true}{http://rocq-prover.org/doc/V9.1.1/stdlib//Corelib.Init.Datatypes}{\coqdocconstructor{true}} (\coqref{SECF.SpecCT.BNot}{\coqdocdefinition{BNot}} (\coqref{SECF.SpecCT.BLe}{\coqdocdefinition{BLe}} \coqref{SECF.SpecCT.X}{\coqdocdefinition{X}} 4)). 
\begin{coqdoccode}
\coqdocemptyline
\coqdocnoindent
\coqdockw{Definition} \coqdef{SECF.Imp.example aexp}{example\_aexp}{\coqdocdefinition{example\_aexp}} : \coqref{SECF.Imp.aexp}{\coqdocinductive{aexp}} := \coqref{SECF.Imp.::com scope:'<x7B' x 'x7D>'}{\coqdocnotation{<\{}} 3 \coqref{SECF.Imp.:com::x '+' x}{\coqdocnotation{+}} \coqref{SECF.Imp.:com::'(' x ')'}{\coqdocnotation{(}}\coqref{SECF.Imp.X}{\coqdocdefinition{X}} \coqref{SECF.Imp.:com::x '*' x}{\coqdocnotation{\ensuremath{\times}}} 2\coqref{SECF.Imp.:com::'(' x ')'}{\coqdocnotation{)}} \coqref{SECF.Imp.::com scope:'<x7B' x 'x7D>'}{\coqdocnotation{\}>}}.\coqdoceol
\coqdocnoindent
\coqdockw{Definition} \coqdef{SECF.Imp.example bexp}{example\_bexp}{\coqdocdefinition{example\_bexp}} : \coqref{SECF.Imp.bexp}{\coqdocinductive{bexp}} := \coqref{SECF.Imp.::com scope:'<x7B' x 'x7D>'}{\coqdocnotation{<\{}} \coqref{SECF.Imp.:com::'true'}{\coqdocnotation{true}} \coqref{SECF.Imp.:com::x 'x26x26' x}{\coqdocnotation{\&\&}} \coqref{SECF.Imp.:com::'x7E' x}{\coqdocnotation{\ensuremath{\lnot}}}\coqref{SECF.Imp.:com::'(' x ')'}{\coqdocnotation{(}}\coqref{SECF.Imp.X}{\coqdocdefinition{X}} \coqref{SECF.Imp.:com::x '<=' x}{\coqdocnotation{\ensuremath{\le}}} 4\coqref{SECF.Imp.:com::'(' x ')'}{\coqdocnotation{)}} \coqref{SECF.Imp.::com scope:'<x7B' x 'x7D>'}{\coqdocnotation{\}>}}.\coqdoceol
\coqdocemptyline
\end{coqdoccode}
\subsection{Evaluation}



 The arith and boolean evaluators must now be extended to
    handle variables in the obvious way, taking a state \coqdocvar{st} as an
    extra argument: 
\begin{coqdoccode}
\coqdocemptyline
\coqdocnoindent
\coqdockw{Fixpoint} \coqdef{SECF.Imp.aeval}{aeval}{\coqdocdefinition{aeval}} (\coqdef{SECF.Imp.st:194}{st}{\coqdocbinder{st}} : \coqref{SECF.Imp.state}{\coqdocdefinition{state}}) \coqdoceol
\coqdocindent{7.50em}
(\coqdef{SECF.Imp.a:195}{a}{\coqdocbinder{a}} : \coqref{SECF.Imp.aexp}{\coqdocinductive{aexp}}) : \coqexternalref{nat}{http://rocq-prover.org/doc/V9.1.1/stdlib//Corelib.Init.Datatypes}{\coqdocinductive{nat}} :=\coqdoceol
\coqdocindent{1.00em}
\coqdockw{match} \coqref{SECF.Imp.a:195}{\coqdocvariable{a}} \coqdockw{with}\coqdoceol
\coqdocindent{1.00em}
\ensuremath{|} \coqref{SECF.Imp.ANum}{\coqdocconstructor{ANum}} \coqdocvar{n} \ensuremath{\Rightarrow} \coqdocvar{n}\coqdoceol
\coqdocindent{1.00em}
\ensuremath{|} \coqref{SECF.Imp.AId}{\coqdocconstructor{AId}} \coqdocvar{x} \ensuremath{\Rightarrow} \coqref{SECF.Imp.st:194}{\coqdocvariable{st}} \coqdocvar{x}                                \coqdoceol
\coqdocindent{1.00em}
\ensuremath{|} \coqref{SECF.Imp.::com scope:'<x7B' x 'x7D>'}{\coqdocnotation{<\{}}\coqdocvar{a1} \coqref{SECF.Imp.:com::x '+' x}{\coqdocnotation{+}} \coqdocvar{a2}\coqref{SECF.Imp.::com scope:'<x7B' x 'x7D>'}{\coqdocnotation{\}>}} \ensuremath{\Rightarrow} \coqexternalref{::nat scope:x '+' x}{http://rocq-prover.org/doc/V9.1.1/stdlib//Corelib.Init.Nat}{\coqdocnotation{(}}\coqref{SECF.Imp.aeval:196}{\coqdocdefinition{aeval}} \coqref{SECF.Imp.st:194}{\coqdocvariable{st}} \coqdocvar{a1}\coqexternalref{::nat scope:x '+' x}{http://rocq-prover.org/doc/V9.1.1/stdlib//Corelib.Init.Nat}{\coqdocnotation{)}} \coqexternalref{::nat scope:x '+' x}{http://rocq-prover.org/doc/V9.1.1/stdlib//Corelib.Init.Nat}{\coqdocnotation{+}} \coqexternalref{::nat scope:x '+' x}{http://rocq-prover.org/doc/V9.1.1/stdlib//Corelib.Init.Nat}{\coqdocnotation{(}}\coqref{SECF.Imp.aeval:196}{\coqdocdefinition{aeval}} \coqref{SECF.Imp.st:194}{\coqdocvariable{st}} \coqdocvar{a2}\coqexternalref{::nat scope:x '+' x}{http://rocq-prover.org/doc/V9.1.1/stdlib//Corelib.Init.Nat}{\coqdocnotation{)}}\coqdoceol
\coqdocindent{1.00em}
\ensuremath{|} \coqref{SECF.Imp.::com scope:'<x7B' x 'x7D>'}{\coqdocnotation{<\{}}\coqdocvar{a1} \coqref{SECF.Imp.:com::x '-' x}{\coqdocnotation{-}} \coqdocvar{a2}\coqref{SECF.Imp.::com scope:'<x7B' x 'x7D>'}{\coqdocnotation{\}>}} \ensuremath{\Rightarrow} \coqexternalref{::nat scope:x '-' x}{http://rocq-prover.org/doc/V9.1.1/stdlib//Corelib.Init.Nat}{\coqdocnotation{(}}\coqref{SECF.Imp.aeval:196}{\coqdocdefinition{aeval}} \coqref{SECF.Imp.st:194}{\coqdocvariable{st}} \coqdocvar{a1}\coqexternalref{::nat scope:x '-' x}{http://rocq-prover.org/doc/V9.1.1/stdlib//Corelib.Init.Nat}{\coqdocnotation{)}} \coqexternalref{::nat scope:x '-' x}{http://rocq-prover.org/doc/V9.1.1/stdlib//Corelib.Init.Nat}{\coqdocnotation{-}} \coqexternalref{::nat scope:x '-' x}{http://rocq-prover.org/doc/V9.1.1/stdlib//Corelib.Init.Nat}{\coqdocnotation{(}}\coqref{SECF.Imp.aeval:196}{\coqdocdefinition{aeval}} \coqref{SECF.Imp.st:194}{\coqdocvariable{st}} \coqdocvar{a2}\coqexternalref{::nat scope:x '-' x}{http://rocq-prover.org/doc/V9.1.1/stdlib//Corelib.Init.Nat}{\coqdocnotation{)}}\coqdoceol
\coqdocindent{1.00em}
\ensuremath{|} \coqref{SECF.Imp.::com scope:'<x7B' x 'x7D>'}{\coqdocnotation{<\{}}\coqdocvar{a1} \coqref{SECF.Imp.:com::x '*' x}{\coqdocnotation{\ensuremath{\times}}} \coqdocvar{a2}\coqref{SECF.Imp.::com scope:'<x7B' x 'x7D>'}{\coqdocnotation{\}>}} \ensuremath{\Rightarrow} \coqexternalref{::nat scope:x '*' x}{http://rocq-prover.org/doc/V9.1.1/stdlib//Corelib.Init.Nat}{\coqdocnotation{(}}\coqref{SECF.Imp.aeval:196}{\coqdocdefinition{aeval}} \coqref{SECF.Imp.st:194}{\coqdocvariable{st}} \coqdocvar{a1}\coqexternalref{::nat scope:x '*' x}{http://rocq-prover.org/doc/V9.1.1/stdlib//Corelib.Init.Nat}{\coqdocnotation{)}} \coqexternalref{::nat scope:x '*' x}{http://rocq-prover.org/doc/V9.1.1/stdlib//Corelib.Init.Nat}{\coqdocnotation{\ensuremath{\times}}} \coqexternalref{::nat scope:x '*' x}{http://rocq-prover.org/doc/V9.1.1/stdlib//Corelib.Init.Nat}{\coqdocnotation{(}}\coqref{SECF.Imp.aeval:196}{\coqdocdefinition{aeval}} \coqref{SECF.Imp.st:194}{\coqdocvariable{st}} \coqdocvar{a2}\coqexternalref{::nat scope:x '*' x}{http://rocq-prover.org/doc/V9.1.1/stdlib//Corelib.Init.Nat}{\coqdocnotation{)}}\coqdoceol
\coqdocindent{1.00em}
\coqdockw{end}.\coqdoceol
\coqdocemptyline
\coqdocnoindent
\coqdockw{Fixpoint} \coqdef{SECF.Imp.beval}{beval}{\coqdocdefinition{beval}} (\coqdef{SECF.Imp.st:198}{st}{\coqdocbinder{st}} : \coqref{SECF.Imp.state}{\coqdocdefinition{state}}) \coqdoceol
\coqdocindent{7.50em}
(\coqdef{SECF.Imp.b:199}{b}{\coqdocbinder{b}} : \coqref{SECF.Imp.bexp}{\coqdocinductive{bexp}}) : \coqexternalref{bool}{http://rocq-prover.org/doc/V9.1.1/stdlib//Corelib.Init.Datatypes}{\coqdocinductive{bool}} :=\coqdoceol
\coqdocindent{1.00em}
\coqdockw{match} \coqref{SECF.Imp.b:199}{\coqdocvariable{b}} \coqdockw{with}\coqdoceol
\coqdocindent{1.00em}
\ensuremath{|} \coqref{SECF.Imp.::com scope:'<x7B' x 'x7D>'}{\coqdocnotation{<\{}}\coqref{SECF.Imp.:com::'true'}{\coqdocnotation{true}}\coqref{SECF.Imp.::com scope:'<x7B' x 'x7D>'}{\coqdocnotation{\}>}}      \ensuremath{\Rightarrow} \coqref{SECF.Imp.:::'true'}{\coqdocnotation{true}}\coqdoceol
\coqdocindent{1.00em}
\ensuremath{|} \coqref{SECF.Imp.::com scope:'<x7B' x 'x7D>'}{\coqdocnotation{<\{}}\coqref{SECF.Imp.:com::'false'}{\coqdocnotation{false}}\coqref{SECF.Imp.::com scope:'<x7B' x 'x7D>'}{\coqdocnotation{\}>}}     \ensuremath{\Rightarrow} \coqref{SECF.Imp.:::'false'}{\coqdocnotation{false}}\coqdoceol
\coqdocindent{1.00em}
\ensuremath{|} \coqref{SECF.Imp.::com scope:'<x7B' x 'x7D>'}{\coqdocnotation{<\{}}\coqdocvar{a1} \coqref{SECF.Imp.:com::x '=' x}{\coqdocnotation{=}} \coqdocvar{a2}\coqref{SECF.Imp.::com scope:'<x7B' x 'x7D>'}{\coqdocnotation{\}>}}   \ensuremath{\Rightarrow} \coqexternalref{::nat scope:x '=?' x}{http://rocq-prover.org/doc/V9.1.1/stdlib//Corelib.Init.Nat}{\coqdocnotation{(}}\coqref{SECF.Imp.aeval}{\coqdocdefinition{aeval}} \coqref{SECF.Imp.st:198}{\coqdocvariable{st}} \coqdocvar{a1}\coqexternalref{::nat scope:x '=?' x}{http://rocq-prover.org/doc/V9.1.1/stdlib//Corelib.Init.Nat}{\coqdocnotation{)}} \coqexternalref{::nat scope:x '=?' x}{http://rocq-prover.org/doc/V9.1.1/stdlib//Corelib.Init.Nat}{\coqdocnotation{=?}} \coqexternalref{::nat scope:x '=?' x}{http://rocq-prover.org/doc/V9.1.1/stdlib//Corelib.Init.Nat}{\coqdocnotation{(}}\coqref{SECF.Imp.aeval}{\coqdocdefinition{aeval}} \coqref{SECF.Imp.st:198}{\coqdocvariable{st}} \coqdocvar{a2}\coqexternalref{::nat scope:x '=?' x}{http://rocq-prover.org/doc/V9.1.1/stdlib//Corelib.Init.Nat}{\coqdocnotation{)}}\coqdoceol
\coqdocindent{1.00em}
\ensuremath{|} \coqref{SECF.Imp.::com scope:'<x7B' x 'x7D>'}{\coqdocnotation{<\{}}\coqdocvar{a1} \coqref{SECF.Imp.:com::x '<>' x}{\coqdocnotation{\ensuremath{\not=}}} \coqdocvar{a2}\coqref{SECF.Imp.::com scope:'<x7B' x 'x7D>'}{\coqdocnotation{\}>}}  \ensuremath{\Rightarrow} \coqexternalref{negb}{http://rocq-prover.org/doc/V9.1.1/stdlib//Corelib.Init.Datatypes}{\coqdocdefinition{negb}} (\coqexternalref{::nat scope:x '=?' x}{http://rocq-prover.org/doc/V9.1.1/stdlib//Corelib.Init.Nat}{\coqdocnotation{(}}\coqref{SECF.Imp.aeval}{\coqdocdefinition{aeval}} \coqref{SECF.Imp.st:198}{\coqdocvariable{st}} \coqdocvar{a1}\coqexternalref{::nat scope:x '=?' x}{http://rocq-prover.org/doc/V9.1.1/stdlib//Corelib.Init.Nat}{\coqdocnotation{)}} \coqexternalref{::nat scope:x '=?' x}{http://rocq-prover.org/doc/V9.1.1/stdlib//Corelib.Init.Nat}{\coqdocnotation{=?}} \coqexternalref{::nat scope:x '=?' x}{http://rocq-prover.org/doc/V9.1.1/stdlib//Corelib.Init.Nat}{\coqdocnotation{(}}\coqref{SECF.Imp.aeval}{\coqdocdefinition{aeval}} \coqref{SECF.Imp.st:198}{\coqdocvariable{st}} \coqdocvar{a2}\coqexternalref{::nat scope:x '=?' x}{http://rocq-prover.org/doc/V9.1.1/stdlib//Corelib.Init.Nat}{\coqdocnotation{)}})\coqdoceol
\coqdocindent{1.00em}
\ensuremath{|} \coqref{SECF.Imp.::com scope:'<x7B' x 'x7D>'}{\coqdocnotation{<\{}}\coqdocvar{a1} \coqref{SECF.Imp.:com::x '<=' x}{\coqdocnotation{\ensuremath{\le}}} \coqdocvar{a2}\coqref{SECF.Imp.::com scope:'<x7B' x 'x7D>'}{\coqdocnotation{\}>}}  \ensuremath{\Rightarrow} \coqexternalref{::nat scope:x '<=?' x}{http://rocq-prover.org/doc/V9.1.1/stdlib//Corelib.Init.Nat}{\coqdocnotation{(}}\coqref{SECF.Imp.aeval}{\coqdocdefinition{aeval}} \coqref{SECF.Imp.st:198}{\coqdocvariable{st}} \coqdocvar{a1}\coqexternalref{::nat scope:x '<=?' x}{http://rocq-prover.org/doc/V9.1.1/stdlib//Corelib.Init.Nat}{\coqdocnotation{)}} \coqexternalref{::nat scope:x '<=?' x}{http://rocq-prover.org/doc/V9.1.1/stdlib//Corelib.Init.Nat}{\coqdocnotation{<=?}} \coqexternalref{::nat scope:x '<=?' x}{http://rocq-prover.org/doc/V9.1.1/stdlib//Corelib.Init.Nat}{\coqdocnotation{(}}\coqref{SECF.Imp.aeval}{\coqdocdefinition{aeval}} \coqref{SECF.Imp.st:198}{\coqdocvariable{st}} \coqdocvar{a2}\coqexternalref{::nat scope:x '<=?' x}{http://rocq-prover.org/doc/V9.1.1/stdlib//Corelib.Init.Nat}{\coqdocnotation{)}}\coqdoceol
\coqdocindent{1.00em}
\ensuremath{|} \coqref{SECF.Imp.::com scope:'<x7B' x 'x7D>'}{\coqdocnotation{<\{}}\coqdocvar{a1} \coqref{SECF.Imp.:com::x '>' x}{\coqdocnotation{>}} \coqdocvar{a2}\coqref{SECF.Imp.::com scope:'<x7B' x 'x7D>'}{\coqdocnotation{\}>}}   \ensuremath{\Rightarrow} \coqexternalref{negb}{http://rocq-prover.org/doc/V9.1.1/stdlib//Corelib.Init.Datatypes}{\coqdocdefinition{negb}} (\coqexternalref{::nat scope:x '<=?' x}{http://rocq-prover.org/doc/V9.1.1/stdlib//Corelib.Init.Nat}{\coqdocnotation{(}}\coqref{SECF.Imp.aeval}{\coqdocdefinition{aeval}} \coqref{SECF.Imp.st:198}{\coqdocvariable{st}} \coqdocvar{a1}\coqexternalref{::nat scope:x '<=?' x}{http://rocq-prover.org/doc/V9.1.1/stdlib//Corelib.Init.Nat}{\coqdocnotation{)}} \coqexternalref{::nat scope:x '<=?' x}{http://rocq-prover.org/doc/V9.1.1/stdlib//Corelib.Init.Nat}{\coqdocnotation{<=?}} \coqexternalref{::nat scope:x '<=?' x}{http://rocq-prover.org/doc/V9.1.1/stdlib//Corelib.Init.Nat}{\coqdocnotation{(}}\coqref{SECF.Imp.aeval}{\coqdocdefinition{aeval}} \coqref{SECF.Imp.st:198}{\coqdocvariable{st}} \coqdocvar{a2}\coqexternalref{::nat scope:x '<=?' x}{http://rocq-prover.org/doc/V9.1.1/stdlib//Corelib.Init.Nat}{\coqdocnotation{)}})\coqdoceol
\coqdocindent{1.00em}
\ensuremath{|} \coqref{SECF.Imp.::com scope:'<x7B' x 'x7D>'}{\coqdocnotation{<\{}}\coqref{SECF.Imp.:com::'x7E' x}{\coqdocnotation{\ensuremath{\lnot}}} \coqdocvar{b1}\coqref{SECF.Imp.::com scope:'<x7B' x 'x7D>'}{\coqdocnotation{\}>}}      \ensuremath{\Rightarrow} \coqexternalref{negb}{http://rocq-prover.org/doc/V9.1.1/stdlib//Corelib.Init.Datatypes}{\coqdocdefinition{negb}} (\coqref{SECF.Imp.beval:200}{\coqdocdefinition{beval}} \coqref{SECF.Imp.st:198}{\coqdocvariable{st}} \coqdocvar{b1})\coqdoceol
\coqdocindent{1.00em}
\ensuremath{|} \coqref{SECF.Imp.::com scope:'<x7B' x 'x7D>'}{\coqdocnotation{<\{}}\coqdocvar{b1} \coqref{SECF.Imp.:com::x 'x26x26' x}{\coqdocnotation{\&\&}} \coqdocvar{b2}\coqref{SECF.Imp.::com scope:'<x7B' x 'x7D>'}{\coqdocnotation{\}>}}  \ensuremath{\Rightarrow} \coqexternalref{andb}{http://rocq-prover.org/doc/V9.1.1/stdlib//Corelib.Init.Datatypes}{\coqdocdefinition{andb}} (\coqref{SECF.Imp.beval:200}{\coqdocdefinition{beval}} \coqref{SECF.Imp.st:198}{\coqdocvariable{st}} \coqdocvar{b1}) (\coqref{SECF.Imp.beval:200}{\coqdocdefinition{beval}} \coqref{SECF.Imp.st:198}{\coqdocvariable{st}} \coqdocvar{b2})\coqdoceol
\coqdocindent{1.00em}
\coqdockw{end}.\coqdoceol
\coqdocemptyline
\end{coqdoccode}
We can use our notation for total maps in the specific case of
    states -- i.e., we write the empty state as (\coqdocvar{\_\_} !-> 0). 
\begin{coqdoccode}
\coqdocemptyline
\coqdocnoindent
\coqdockw{Definition} \coqdef{SECF.Imp.empty st}{empty\_st}{\coqdocdefinition{empty\_st}} := (\coqref{SECF.Maps.:::'  ' '!->' x}{\coqdocnotation{\_\_}} \coqref{SECF.Maps.:::'  ' '!->' x}{\coqdocnotation{!->}} 0).\coqdoceol
\coqdocemptyline
\end{coqdoccode}
Also, we can add a notation for a ``singleton state'' with just one
    variable bound to a value. 
\begin{coqdoccode}
\coqdocnoindent
\coqdockw{Notation} \coqdef{SECF.Imp.:::x '!->' x}{"}{"}x '!->' v" := (\coqdocvar{x} \coqref{SECF.Maps.:::x '!->' x ';' x}{\coqdocnotation{!->}} \coqdocvar{v} \coqref{SECF.Maps.:::x '!->' x ';' x}{\coqdocnotation{;}} \coqref{SECF.Imp.empty st}{\coqdocdefinition{empty\_st}}) (\coqdoctac{at} \coqdockw{level} 100, \coqdoctac{right} \coqdockw{associativity}).\coqdoceol
\coqdocemptyline
\coqdocnoindent
\coqdockw{Example} \coqdef{SECF.Imp.aexp1}{aexp1}{\coqdocdefinition{aexp1}} :\coqdoceol
\coqdocindent{2.00em}
\coqref{SECF.Imp.aeval}{\coqdocdefinition{aeval}} (\coqref{SECF.Imp.X}{\coqdocdefinition{X}} \coqref{SECF.Imp.:::x '!->' x}{\coqdocnotation{!->}} 5) \coqref{SECF.Imp.::com scope:'<x7B' x 'x7D>'}{\coqdocnotation{<\{}} 3 \coqref{SECF.Imp.:com::x '+' x}{\coqdocnotation{+}} \coqref{SECF.Imp.:com::'(' x ')'}{\coqdocnotation{(}}\coqref{SECF.Imp.X}{\coqdocdefinition{X}} \coqref{SECF.Imp.:com::x '*' x}{\coqdocnotation{\ensuremath{\times}}} 2\coqref{SECF.Imp.:com::'(' x ')'}{\coqdocnotation{)}} \coqref{SECF.Imp.::com scope:'<x7B' x 'x7D>'}{\coqdocnotation{\}>}}\coqdoceol
\coqdocindent{1.00em}
\coqexternalref{::type scope:x '=' x}{http://rocq-prover.org/doc/V9.1.1/stdlib//Corelib.Init.Logic}{\coqdocnotation{=}} 13.\coqdoceol
\coqdocnoindent
\coqdockw{Proof}. \coqdoctac{reflexivity}. \coqdockw{Qed}.\coqdoceol
\coqdocnoindent
\coqdockw{Example} \coqdef{SECF.Imp.aexp2}{aexp2}{\coqdocdefinition{aexp2}} :\coqdoceol
\coqdocindent{2.00em}
\coqref{SECF.Imp.aeval}{\coqdocdefinition{aeval}} (\coqref{SECF.Imp.X}{\coqdocdefinition{X}} \coqref{SECF.Maps.:::x '!->' x ';' x}{\coqdocnotation{!->}} 5 \coqref{SECF.Maps.:::x '!->' x ';' x}{\coqdocnotation{;}} \coqref{SECF.Imp.Y}{\coqdocdefinition{Y}} \coqref{SECF.Imp.:::x '!->' x}{\coqdocnotation{!->}} 4) \coqref{SECF.Imp.::com scope:'<x7B' x 'x7D>'}{\coqdocnotation{<\{}} \coqref{SECF.Imp.Z}{\coqdocdefinition{Z}} \coqref{SECF.Imp.:com::x '+' x}{\coqdocnotation{+}} \coqref{SECF.Imp.:com::'(' x ')'}{\coqdocnotation{(}}\coqref{SECF.Imp.X}{\coqdocdefinition{X}} \coqref{SECF.Imp.:com::x '*' x}{\coqdocnotation{\ensuremath{\times}}} \coqref{SECF.Imp.Y}{\coqdocdefinition{Y}}\coqref{SECF.Imp.:com::'(' x ')'}{\coqdocnotation{)}} \coqref{SECF.Imp.::com scope:'<x7B' x 'x7D>'}{\coqdocnotation{\}>}}\coqdoceol
\coqdocindent{1.00em}
\coqexternalref{::type scope:x '=' x}{http://rocq-prover.org/doc/V9.1.1/stdlib//Corelib.Init.Logic}{\coqdocnotation{=}} 20.\coqdoceol
\coqdocnoindent
\coqdockw{Proof}. \coqdoctac{reflexivity}. \coqdockw{Qed}.\coqdoceol
\coqdocemptyline
\coqdocnoindent
\coqdockw{Example} \coqdef{SECF.Imp.bexp1}{bexp1}{\coqdocdefinition{bexp1}} :\coqdoceol
\coqdocindent{2.00em}
\coqref{SECF.Imp.beval}{\coqdocdefinition{beval}} (\coqref{SECF.Imp.X}{\coqdocdefinition{X}} \coqref{SECF.Imp.:::x '!->' x}{\coqdocnotation{!->}} 5) \coqref{SECF.Imp.::com scope:'<x7B' x 'x7D>'}{\coqdocnotation{<\{}} \coqref{SECF.Imp.:com::'true'}{\coqdocnotation{true}} \coqref{SECF.Imp.:com::x 'x26x26' x}{\coqdocnotation{\&\&}} \coqref{SECF.Imp.:com::'x7E' x}{\coqdocnotation{\ensuremath{\lnot}}}\coqref{SECF.Imp.:com::'(' x ')'}{\coqdocnotation{(}}\coqref{SECF.Imp.X}{\coqdocdefinition{X}} \coqref{SECF.Imp.:com::x '<=' x}{\coqdocnotation{\ensuremath{\le}}} 4\coqref{SECF.Imp.:com::'(' x ')'}{\coqdocnotation{)}} \coqref{SECF.Imp.::com scope:'<x7B' x 'x7D>'}{\coqdocnotation{\}>}}\coqdoceol
\coqdocindent{1.00em}
\coqexternalref{::type scope:x '=' x}{http://rocq-prover.org/doc/V9.1.1/stdlib//Corelib.Init.Logic}{\coqdocnotation{=}} \coqref{SECF.Imp.:::'true'}{\coqdocnotation{true}}.\coqdoceol
\coqdocnoindent
\coqdockw{Proof}. \coqdoctac{reflexivity}. \coqdockw{Qed}.\coqdoceol
\coqdocemptyline
\end{coqdoccode}
\section{Commands}



 Now we are ready to define the syntax and behavior of Imp
    \textit{commands} (or \textit{statements}). 
\begin{coqdoccode}
\coqdocemptyline
\end{coqdoccode}
\subsection{Syntax}



 Informally, commands \coqdocvar{c} are described by the following BNF
    grammar.


     c := skip
        | x := a
        | c ; c
        | if b then c else c end
        | while b do c end


 Here is the formal definition of the abstract syntax of
    commands: 
\begin{coqdoccode}
\coqdocemptyline
\coqdocnoindent
\coqdockw{Inductive} \coqdef{SECF.Imp.com}{com}{\coqdocinductive{com}} : \coqdockw{Type} :=\coqdoceol
\coqdocindent{1.00em}
\ensuremath{|} \coqdef{SECF.Imp.CSkip}{CSkip}{\coqdocconstructor{CSkip}}\coqdoceol
\coqdocindent{1.00em}
\ensuremath{|} \coqdef{SECF.Imp.CAsgn}{CAsgn}{\coqdocconstructor{CAsgn}} (\coqdef{SECF.Imp.x:204}{x}{\coqdocbinder{x}} : \coqdocinductive{string}) (\coqdef{SECF.Imp.a:205}{a}{\coqdocbinder{a}} : \coqref{SECF.Imp.aexp}{\coqdocinductive{aexp}})\coqdoceol
\coqdocindent{1.00em}
\ensuremath{|} \coqdef{SECF.Imp.CSeq}{CSeq}{\coqdocconstructor{CSeq}} (\coqdef{SECF.Imp.c1:206}{c1}{\coqdocbinder{c1}} \coqdef{SECF.Imp.c2:207}{c2}{\coqdocbinder{c2}} : \coqref{SECF.Imp.com:202}{\coqdocinductive{com}})\coqdoceol
\coqdocindent{1.00em}
\ensuremath{|} \coqdef{SECF.Imp.CIf}{CIf}{\coqdocconstructor{CIf}} (\coqdef{SECF.Imp.b:208}{b}{\coqdocbinder{b}} : \coqref{SECF.Imp.bexp}{\coqdocinductive{bexp}}) (\coqdef{SECF.Imp.c1:209}{c1}{\coqdocbinder{c1}} \coqdef{SECF.Imp.c2:210}{c2}{\coqdocbinder{c2}} : \coqref{SECF.Imp.com:202}{\coqdocinductive{com}})\coqdoceol
\coqdocindent{1.00em}
\ensuremath{|} \coqdef{SECF.Imp.CWhile}{CWhile}{\coqdocconstructor{CWhile}} (\coqdef{SECF.Imp.b:211}{b}{\coqdocbinder{b}} : \coqref{SECF.Imp.bexp}{\coqdocinductive{bexp}}) (\coqdef{SECF.Imp.c:212}{c}{\coqdocbinder{c}} : \coqref{SECF.Imp.com:202}{\coqdocinductive{com}}).\coqdoceol
\coqdocemptyline
\end{coqdoccode}
As we did for expressions, we can use a few \coqdockw{Notation}
    declarations to make reading and writing Imp programs more
    convenient. 
\begin{coqdoccode}
\coqdocemptyline
\coqdocnoindent
\coqdockw{Notation} \coqdef{SECF.Imp.:com:com scope:'skip'}{"}{"}'skip'"  := \coqref{SECF.Imp.CSkip}{\coqdocconstructor{CSkip}}\coqdoceol
\coqdocindent{1.00em}
(\coqdoctac{in} \coqdocvar{custom} \coqdocvar{com} \coqdoctac{at} \coqdockw{level} 0) : \coqdocvar{com\_scope}.\coqdoceol
\coqdocnoindent
\coqdockw{Notation} \coqdef{SECF.Imp.:com:com scope:x ':=' x}{"}{"}x := y"  := (\coqref{SECF.Imp.CAsgn}{\coqdocconstructor{CAsgn}} \coqdocvar{x} \coqdocvar{y})\coqdoceol
\coqdocindent{1.00em}
(\coqdoctac{in} \coqdocvar{custom} \coqdocvar{com} \coqdoctac{at} \coqdockw{level} 0, \coqdocvar{x} \coqdockw{constr} \coqdoctac{at} \coqdockw{level} 0, \coqdocvar{y} \coqdoctac{at} \coqdockw{level} 85, \coqdockw{no} \coqdockw{associativity},\coqdoceol
\coqdocindent{2.00em}
\coqdocvar{format} "x  :=  y") : \coqdocvar{com\_scope}.\coqdoceol
\coqdocnoindent
\coqdockw{Notation} \coqdef{SECF.Imp.:com:com scope:x ';' x}{"}{"}x ; y" := (\coqref{SECF.Imp.CSeq}{\coqdocconstructor{CSeq}} \coqdocvar{x} \coqdocvar{y})\coqdoceol
\coqdocindent{1.00em}
(\coqdoctac{in} \coqdocvar{custom} \coqdocvar{com} \coqdoctac{at} \coqdockw{level} 90,\coqdoceol
\coqdocindent{2.00em}
\coqdoctac{right} \coqdockw{associativity},\coqdoceol
\coqdocindent{2.00em}
\coqdocvar{format} "'[v' x ; '/' y ']'") : \coqdocvar{com\_scope}.\coqdoceol
\coqdocnoindent
\coqdockw{Notation} \coqdef{SECF.Imp.:com:com scope:'if' x 'then' x 'else' x 'end'}{"}{"}'if' x 'then' y 'else' z 'end'" := (\coqref{SECF.Imp.CIf}{\coqdocconstructor{CIf}} \coqdocvar{x} \coqdocvar{y} \coqdocvar{z})\coqdoceol
\coqdocindent{1.00em}
(\coqdoctac{in} \coqdocvar{custom} \coqdocvar{com} \coqdoctac{at} \coqdockw{level} 89, \coqdocvar{x} \coqdoctac{at} \coqdockw{level} 99, \coqdocvar{y} \coqdoctac{at} \coqdockw{level} 99, \coqdocvar{z} \coqdoctac{at} \coqdockw{level} 99,\coqdoceol
\coqdocindent{2.00em}
\coqdocvar{format} "'[v' 'if'  x  'then' '/  ' y '/' 'else' '/  ' z '/' 'end' ']'") : \coqdocvar{com\_scope}.\coqdoceol
\coqdocnoindent
\coqdockw{Notation} \coqdef{SECF.Imp.:com:com scope:'while' x 'do' x 'end'}{"}{"}'while' x 'do' y 'end'" := (\coqref{SECF.Imp.CWhile}{\coqdocconstructor{CWhile}} \coqdocvar{x} \coqdocvar{y})\coqdoceol
\coqdocindent{1.00em}
(\coqdoctac{in} \coqdocvar{custom} \coqdocvar{com} \coqdoctac{at} \coqdockw{level} 89, \coqdocvar{x} \coqdoctac{at} \coqdockw{level} 99, \coqdocvar{y} \coqdoctac{at} \coqdockw{level} 99,\coqdoceol
\coqdocindent{2.00em}
\coqdocvar{format} "'[v' 'while'  x  'do' '/  ' y '/' 'end' ']'") : \coqdocvar{com\_scope}.\coqdoceol
\coqdocemptyline
\end{coqdoccode}
For example, here is the factorial function again, written as a
    formal Rocq definition.  When this command terminates, the variable
    \coqref{SECF.SpecCT.Y}{\coqdocdefinition{Y}} will contain the factorial of the initial value of \coqref{SECF.SpecCT.X}{\coqdocdefinition{X}}. 
\begin{coqdoccode}
\coqdocemptyline
\coqdocnoindent
\coqdockw{Definition} \coqdef{SECF.Imp.fact in coq}{fact\_in\_coq}{\coqdocdefinition{fact\_in\_coq}} : \coqref{SECF.Imp.com}{\coqdocinductive{com}} :=\coqdoceol
\coqdocindent{1.00em}
\coqref{SECF.Imp.::com scope:'<x7B' x 'x7D>'}{\coqdocnotation{<\{}} \coqref{SECF.Imp.Z}{\coqdocdefinition{Z}} \coqref{SECF.Imp.:com:com scope:x ':=' x}{\coqdocnotation{:=}} \coqref{SECF.Imp.X}{\coqdocdefinition{X}}\coqref{SECF.Imp.:com:com scope:x ';' x}{\coqdocnotation{;}}\coqdoceol
\coqdocindent{2.50em}
\coqref{SECF.Imp.Y}{\coqdocdefinition{Y}} \coqref{SECF.Imp.:com:com scope:x ':=' x}{\coqdocnotation{:=}} 1\coqref{SECF.Imp.:com:com scope:x ';' x}{\coqdocnotation{;}}\coqdoceol
\coqdocindent{2.50em}
\coqref{SECF.Imp.:com:com scope:'while' x 'do' x 'end'}{\coqdocnotation{while}} \coqref{SECF.Imp.Z}{\coqdocdefinition{Z}} \coqref{SECF.Imp.:com::x '<>' x}{\coqdocnotation{\ensuremath{\not=}}} 0 \coqref{SECF.Imp.:com:com scope:'while' x 'do' x 'end'}{\coqdocnotation{do}}\coqdoceol
\coqdocindent{3.50em}
\coqref{SECF.Imp.Y}{\coqdocdefinition{Y}} \coqref{SECF.Imp.:com:com scope:x ':=' x}{\coqdocnotation{:=}} \coqref{SECF.Imp.Y}{\coqdocdefinition{Y}} \coqref{SECF.Imp.:com::x '*' x}{\coqdocnotation{\ensuremath{\times}}} \coqref{SECF.Imp.Z}{\coqdocdefinition{Z}}\coqref{SECF.Imp.:com:com scope:x ';' x}{\coqdocnotation{;}}\coqdoceol
\coqdocindent{3.50em}
\coqref{SECF.Imp.Z}{\coqdocdefinition{Z}} \coqref{SECF.Imp.:com:com scope:x ':=' x}{\coqdocnotation{:=}} \coqref{SECF.Imp.Z}{\coqdocdefinition{Z}} \coqref{SECF.Imp.:com::x '-' x}{\coqdocnotation{-}} 1\coqdoceol
\coqdocindent{2.50em}
\coqref{SECF.Imp.:com:com scope:'while' x 'do' x 'end'}{\coqdocnotation{end}} \coqref{SECF.Imp.::com scope:'<x7B' x 'x7D>'}{\coqdocnotation{\}>}}.\coqdoceol
\coqdocemptyline
\coqdocnoindent
\coqdockw{Print} \coqdocvar{fact\_in\_coq}.\coqdoceol
\coqdocemptyline
\end{coqdoccode}
\subsection{Desugaring Notations}



 Rocq offers a rich set of features to manage the increasing
    complexity of the objects we work with, such as coercions and
    notations. However, their heavy usage can make it hard to
    understand what the expressions we enter actually mean. In such
    situations it is often instructive to ``turn off'' those features to
    get a more elementary picture of things, using the following
    commands:



\begin{itemize}
\item  \coqdockw{Unset} \coqdockw{Printing} \coqdocvar{Notations} (undo with \coqdockw{Set} \coqdockw{Printing} \coqdocvar{Notations})

\item  \coqdockw{Set} \coqdockw{Printing} \coqdocvar{Coercions} (undo with \coqdockw{Unset} \coqdockw{Printing} \coqdocvar{Coercions})

\item  \coqdockw{Set} \coqdockw{Printing} \coqdockw{All} (undo with \coqdockw{Unset} \coqdockw{Printing} \coqdockw{All})

\end{itemize}


    These commands can also be used in the middle of a proof, to
    elaborate the current goal and context. 
\begin{coqdoccode}
\coqdocemptyline
\coqdocnoindent
\coqdockw{Unset} \coqdockw{Printing} \coqdocvar{Notations}.\coqdoceol
\coqdocnoindent
\coqdockw{Print} \coqdocvar{fact\_in\_coq}.\coqdoceol
\coqdocnoindent
\coqdockw{Set} \coqdockw{Printing} \coqdocvar{Notations}.\coqdoceol
\coqdocemptyline
\coqdocnoindent
\coqdockw{Print} \coqdocvar{example\_bexp}.\coqdoceol
\coqdocemptyline
\coqdocnoindent
\coqdockw{Set Printing Coercions}.\coqdoceol
\coqdocnoindent
\coqdockw{Print} \coqdocvar{example\_bexp}.\coqdoceol
\coqdocemptyline
\coqdocnoindent
\coqdockw{Print} \coqdocvar{fact\_in\_coq}.\coqdoceol
\coqdocnoindent
\coqdockw{Unset Printing Coercions}.\coqdoceol
\coqdocemptyline
\end{coqdoccode}
\subsection{\texorpdfstring{\protect\coqdockw{Locate}}{Locate} Again}


\begin{coqdoccode}
\coqdocemptyline
\end{coqdoccode}
\subsubsection{Finding identifiers}



 When used with an identifier, the \coqdockw{Locate} prints the full path to
    every value in scope with the same name.  This is useful to
    troubleshoot problems due to variable shadowing. 
\begin{coqdoccode}
\coqdocnoindent
\coqdockw{Locate} \coqdocvar{aexp}.\coqdoceol
\end{coqdoccode}
\subsubsection{Finding notations}



 When faced with an unknown notation, you can use \coqdockw{Locate} with a
    string containing one of its symbols to see its possible
    interpretations. 
\begin{coqdoccode}
\coqdocnoindent
\coqdockw{Locate} "\&\&".\coqdoceol
\coqdocnoindent
\coqdockw{Locate} ";".\coqdoceol
\coqdocemptyline
\coqdocnoindent
\coqdockw{Locate} "while".\coqdoceol
\coqdocemptyline
\end{coqdoccode}
\subsection{More Examples}


\begin{coqdoccode}
\coqdocemptyline
\end{coqdoccode}
\subsubsection{Assignment:}


\begin{coqdoccode}
\coqdocemptyline
\coqdocnoindent
\coqdockw{Definition} \coqdef{SECF.Imp.plus2}{plus2}{\coqdocdefinition{plus2}} : \coqref{SECF.Imp.com}{\coqdocinductive{com}} :=\coqdoceol
\coqdocindent{1.00em}
\coqref{SECF.Imp.::com scope:'<x7B' x 'x7D>'}{\coqdocnotation{<\{}} \coqref{SECF.Imp.X}{\coqdocdefinition{X}} \coqref{SECF.Imp.:com:com scope:x ':=' x}{\coqdocnotation{:=}} \coqref{SECF.Imp.X}{\coqdocdefinition{X}} \coqref{SECF.Imp.:com::x '+' x}{\coqdocnotation{+}} 2 \coqref{SECF.Imp.::com scope:'<x7B' x 'x7D>'}{\coqdocnotation{\}>}}.\coqdoceol
\coqdocemptyline
\coqdocnoindent
\coqdockw{Definition} \coqdef{SECF.Imp.XtimesYinZ}{XtimesYinZ}{\coqdocdefinition{XtimesYinZ}} : \coqref{SECF.Imp.com}{\coqdocinductive{com}} :=\coqdoceol
\coqdocindent{1.00em}
\coqref{SECF.Imp.::com scope:'<x7B' x 'x7D>'}{\coqdocnotation{<\{}} \coqref{SECF.Imp.Z}{\coqdocdefinition{Z}} \coqref{SECF.Imp.:com:com scope:x ':=' x}{\coqdocnotation{:=}} \coqref{SECF.Imp.X}{\coqdocdefinition{X}} \coqref{SECF.Imp.:com::x '*' x}{\coqdocnotation{\ensuremath{\times}}} \coqref{SECF.Imp.Y}{\coqdocdefinition{Y}} \coqref{SECF.Imp.::com scope:'<x7B' x 'x7D>'}{\coqdocnotation{\}>}}.\coqdoceol
\coqdocemptyline
\end{coqdoccode}
\subsubsection{Loops}


\begin{coqdoccode}
\coqdocemptyline
\coqdocnoindent
\coqdockw{Definition} \coqdef{SECF.Imp.subtract slowly body}{subtract\_slowly\_body}{\coqdocdefinition{subtract\_slowly\_body}} : \coqref{SECF.Imp.com}{\coqdocinductive{com}} :=\coqdoceol
\coqdocindent{1.00em}
\coqref{SECF.Imp.::com scope:'<x7B' x 'x7D>'}{\coqdocnotation{<\{}} \coqref{SECF.Imp.Z}{\coqdocdefinition{Z}} \coqref{SECF.Imp.:com:com scope:x ':=' x}{\coqdocnotation{:=}} \coqref{SECF.Imp.Z}{\coqdocdefinition{Z}} \coqref{SECF.Imp.:com::x '-' x}{\coqdocnotation{-}} 1 \coqref{SECF.Imp.:com:com scope:x ';' x}{\coqdocnotation{;}}\coqdoceol
\coqdocindent{2.50em}
\coqref{SECF.Imp.X}{\coqdocdefinition{X}} \coqref{SECF.Imp.:com:com scope:x ':=' x}{\coqdocnotation{:=}} \coqref{SECF.Imp.X}{\coqdocdefinition{X}} \coqref{SECF.Imp.:com::x '-' x}{\coqdocnotation{-}} 1 \coqref{SECF.Imp.::com scope:'<x7B' x 'x7D>'}{\coqdocnotation{\}>}}.\coqdoceol
\coqdocemptyline
\coqdocnoindent
\coqdockw{Definition} \coqdef{SECF.Imp.subtract slowly}{subtract\_slowly}{\coqdocdefinition{subtract\_slowly}} : \coqref{SECF.Imp.com}{\coqdocinductive{com}} :=\coqdoceol
\coqdocindent{1.00em}
\coqref{SECF.Imp.::com scope:'<x7B' x 'x7D>'}{\coqdocnotation{<\{}} \coqref{SECF.Imp.:com:com scope:'while' x 'do' x 'end'}{\coqdocnotation{while}} \coqref{SECF.Imp.X}{\coqdocdefinition{X}} \coqref{SECF.Imp.:com::x '<>' x}{\coqdocnotation{\ensuremath{\not=}}} 0 \coqref{SECF.Imp.:com:com scope:'while' x 'do' x 'end'}{\coqdocnotation{do}}\coqdoceol
\coqdocindent{3.50em}
\coqref{SECF.Imp.subtract slowly body}{\coqdocdefinition{subtract\_slowly\_body}}\coqdoceol
\coqdocindent{2.50em}
\coqref{SECF.Imp.:com:com scope:'while' x 'do' x 'end'}{\coqdocnotation{end}} \coqref{SECF.Imp.::com scope:'<x7B' x 'x7D>'}{\coqdocnotation{\}>}}.\coqdoceol
\coqdocemptyline
\coqdocnoindent
\coqdockw{Definition} \coqdef{SECF.Imp.subtract 3 from 5 slowly}{subtract\_3\_from\_5\_slowly}{\coqdocdefinition{subtract\_3\_from\_5\_slowly}} : \coqref{SECF.Imp.com}{\coqdocinductive{com}} :=\coqdoceol
\coqdocindent{1.00em}
\coqref{SECF.Imp.::com scope:'<x7B' x 'x7D>'}{\coqdocnotation{<\{}} \coqref{SECF.Imp.X}{\coqdocdefinition{X}} \coqref{SECF.Imp.:com:com scope:x ':=' x}{\coqdocnotation{:=}} 3 \coqref{SECF.Imp.:com:com scope:x ';' x}{\coqdocnotation{;}}\coqdoceol
\coqdocindent{2.50em}
\coqref{SECF.Imp.Z}{\coqdocdefinition{Z}} \coqref{SECF.Imp.:com:com scope:x ':=' x}{\coqdocnotation{:=}} 5 \coqref{SECF.Imp.:com:com scope:x ';' x}{\coqdocnotation{;}}\coqdoceol
\coqdocindent{2.50em}
\coqref{SECF.Imp.subtract slowly}{\coqdocdefinition{subtract\_slowly}} \coqref{SECF.Imp.::com scope:'<x7B' x 'x7D>'}{\coqdocnotation{\}>}}.\coqdoceol
\coqdocemptyline
\end{coqdoccode}
\subsubsection{An infinite loop:}


\begin{coqdoccode}
\coqdocemptyline
\coqdocnoindent
\coqdockw{Definition} \coqdef{SECF.Imp.loop}{loop}{\coqdocdefinition{loop}} : \coqref{SECF.Imp.com}{\coqdocinductive{com}} :=\coqdoceol
\coqdocindent{1.00em}
\coqref{SECF.Imp.::com scope:'<x7B' x 'x7D>'}{\coqdocnotation{<\{}} \coqref{SECF.Imp.:com:com scope:'while' x 'do' x 'end'}{\coqdocnotation{while}} \coqref{SECF.Imp.:com::'true'}{\coqdocnotation{true}} \coqref{SECF.Imp.:com:com scope:'while' x 'do' x 'end'}{\coqdocnotation{do}}\coqdoceol
\coqdocindent{3.50em}
\coqref{SECF.Imp.:com:com scope:'skip'}{\coqdocnotation{skip}}\coqdoceol
\coqdocindent{2.50em}
\coqref{SECF.Imp.:com:com scope:'while' x 'do' x 'end'}{\coqdocnotation{end}} \coqref{SECF.Imp.::com scope:'<x7B' x 'x7D>'}{\coqdocnotation{\}>}}.\coqdoceol
\coqdocemptyline
\end{coqdoccode}
\section{Evaluating Commands}



 Next we need to define what it means to evaluate an Imp command.
    The fact that \coqdocvar{while} loops don't necessarily terminate makes
    defining an evaluation function tricky... 
\begin{coqdoccode}
\coqdocemptyline
\end{coqdoccode}
\subsection{Evaluation as a Function (Failed Attempt)}



 Here's an attempt at defining an evaluation function for commands
    (with a bogus \coqdocvar{while} case). 
\begin{coqdoccode}
\coqdocemptyline
\coqdocnoindent
\coqdockw{Fixpoint} \coqdef{SECF.Imp.ceval fun no while}{ceval\_fun\_no\_while}{\coqdocdefinition{ceval\_fun\_no\_while}} (\coqdef{SECF.Imp.st:213}{st}{\coqdocbinder{st}} : \coqref{SECF.Imp.state}{\coqdocdefinition{state}}) (\coqdef{SECF.Imp.c:214}{c}{\coqdocbinder{c}} : \coqref{SECF.Imp.com}{\coqdocinductive{com}}) : \coqref{SECF.Imp.state}{\coqdocdefinition{state}} :=\coqdoceol
\coqdocindent{1.00em}
\coqdockw{match} \coqref{SECF.Imp.c:214}{\coqdocvariable{c}} \coqdockw{with}\coqdoceol
\coqdocindent{2.00em}
\ensuremath{|} \coqref{SECF.Imp.::com scope:'<x7B' x 'x7D>'}{\coqdocnotation{<\{}} \coqref{SECF.Imp.:com:com scope:'skip'}{\coqdocnotation{skip}} \coqref{SECF.Imp.::com scope:'<x7B' x 'x7D>'}{\coqdocnotation{\}>}} \ensuremath{\Rightarrow}\coqdoceol
\coqdocindent{4.00em}
\coqref{SECF.Imp.st:213}{\coqdocvariable{st}}\coqdoceol
\coqdocindent{2.00em}
\ensuremath{|} \coqref{SECF.Imp.::com scope:'<x7B' x 'x7D>'}{\coqdocnotation{<\{}} \coqdocvar{x} \coqref{SECF.Imp.:com:com scope:x ':=' x}{\coqdocnotation{:=}} \coqdocvar{a} \coqref{SECF.Imp.::com scope:'<x7B' x 'x7D>'}{\coqdocnotation{\}>}} \ensuremath{\Rightarrow}\coqdoceol
\coqdocindent{4.00em}
(\coqdocvar{x} \coqref{SECF.Maps.:::x '!->' x ';' x}{\coqdocnotation{!->}} \coqref{SECF.Imp.aeval}{\coqdocdefinition{aeval}} \coqref{SECF.Imp.st:213}{\coqdocvariable{st}} \coqdocvar{a} \coqref{SECF.Maps.:::x '!->' x ';' x}{\coqdocnotation{;}} \coqref{SECF.Imp.st:213}{\coqdocvariable{st}})\coqdoceol
\coqdocindent{2.00em}
\ensuremath{|} \coqref{SECF.Imp.::com scope:'<x7B' x 'x7D>'}{\coqdocnotation{<\{}} \coqdocvar{c1} \coqref{SECF.Imp.:com:com scope:x ';' x}{\coqdocnotation{;}} \coqdocvar{c2} \coqref{SECF.Imp.::com scope:'<x7B' x 'x7D>'}{\coqdocnotation{\}>}} \ensuremath{\Rightarrow}\coqdoceol
\coqdocindent{4.00em}
\coqdockw{let} \coqdef{SECF.Imp.st':217}{st'}{\coqdocbinder{st'}} := \coqref{SECF.Imp.ceval fun no while:215}{\coqdocdefinition{ceval\_fun\_no\_while}} \coqref{SECF.Imp.st:213}{\coqdocvariable{st}} \coqdocvar{c1} \coqdoctac{in}\coqdoceol
\coqdocindent{4.00em}
\coqref{SECF.Imp.ceval fun no while:215}{\coqdocdefinition{ceval\_fun\_no\_while}} \coqref{SECF.Imp.st':217}{\coqdocvariable{st'}} \coqdocvar{c2}\coqdoceol
\coqdocindent{2.00em}
\ensuremath{|} \coqref{SECF.Imp.::com scope:'<x7B' x 'x7D>'}{\coqdocnotation{<\{}} \coqref{SECF.Imp.:com:com scope:'if' x 'then' x 'else' x 'end'}{\coqdocnotation{if}} \coqdocvar{b} \coqref{SECF.Imp.:com:com scope:'if' x 'then' x 'else' x 'end'}{\coqdocnotation{then}} \coqdocvar{c1} \coqref{SECF.Imp.:com:com scope:'if' x 'then' x 'else' x 'end'}{\coqdocnotation{else}} \coqdocvar{c2} \coqref{SECF.Imp.:com:com scope:'if' x 'then' x 'else' x 'end'}{\coqdocnotation{end}}\coqref{SECF.Imp.::com scope:'<x7B' x 'x7D>'}{\coqdocnotation{\}>}} \ensuremath{\Rightarrow}\coqdoceol
\coqdocindent{4.00em}
\coqdockw{if} (\coqref{SECF.Imp.beval}{\coqdocdefinition{beval}} \coqref{SECF.Imp.st:213}{\coqdocvariable{st}} \coqdocvar{b})\coqdoceol
\coqdocindent{5.00em}
\coqdockw{then} \coqref{SECF.Imp.ceval fun no while:215}{\coqdocdefinition{ceval\_fun\_no\_while}} \coqref{SECF.Imp.st:213}{\coqdocvariable{st}} \coqdocvar{c1}\coqdoceol
\coqdocindent{5.00em}
\coqdockw{else} \coqref{SECF.Imp.ceval fun no while:215}{\coqdocdefinition{ceval\_fun\_no\_while}} \coqref{SECF.Imp.st:213}{\coqdocvariable{st}} \coqdocvar{c2}\coqdoceol
\coqdocindent{2.00em}
\ensuremath{|} \coqref{SECF.Imp.::com scope:'<x7B' x 'x7D>'}{\coqdocnotation{<\{}} \coqref{SECF.Imp.:com:com scope:'while' x 'do' x 'end'}{\coqdocnotation{while}} \coqdocvar{b} \coqref{SECF.Imp.:com:com scope:'while' x 'do' x 'end'}{\coqdocnotation{do}} \coqdocvar{c} \coqref{SECF.Imp.:com:com scope:'while' x 'do' x 'end'}{\coqdocnotation{end}} \coqref{SECF.Imp.::com scope:'<x7B' x 'x7D>'}{\coqdocnotation{\}>}} \ensuremath{\Rightarrow}\coqdoceol
\coqdocindent{4.00em}
\coqref{SECF.Imp.st:213}{\coqdocvariable{st}}  \coqdoceol
\coqdocindent{1.00em}
\coqdockw{end}.\coqdoceol
\coqdocemptyline
\end{coqdoccode}
In a more conventional functional programming language like OCaml or
    Haskell we could add the \coqdocvar{while} case as follows:


        Fixpoint ceval\_fun (st : state) (c : com) : state :=
          match c with
            ...
            | <\{ while b do c end\}> =>
                if (beval st b)
                  then ceval\_fun st <\{c ; while b do c end\}>
                  else st
          end.


    Rocq doesn't accept such a definition (``Error: Cannot guess
    decreasing argument of fix'') because the function we want to
    define is not guaranteed to terminate. Indeed, it \textit{doesn't} always
    terminate: for example, the full version of the \coqdocvar{ceval\_fun}
    function applied to the \coqref{SECF.Imp.loop}{\coqdocdefinition{loop}} program above would never
    terminate. Since Rocq aims to be not just a functional programming
    language but also a consistent logic, any potentially
    non-terminating function needs to be rejected.


    Here is an example showing what would go wrong if Rocq allowed
    non-terminating recursive functions:


         Fixpoint loop\_false (n : nat) : False := loop\_false n.


    That is, propositions like \coqexternalref{False}{http://rocq-prover.org/doc/V9.1.1/stdlib//Corelib.Init.Logic}{\coqdocinductive{False}} would become provable
    (\coqdocvar{loop\_false} 0 would be a proof of \coqexternalref{False}{http://rocq-prover.org/doc/V9.1.1/stdlib//Corelib.Init.Logic}{\coqdocinductive{False}}), which would be
    a disaster for Rocq's logical consistency.


    Thus, because it doesn't terminate on all inputs, \coqdocvar{ceval\_fun}
    cannot be written in Rocq -- at least not without additional tricks
    and workarounds (see chapter \coqdocvar{ImpCEvalFun} if you're curious
    about those). 
\begin{coqdoccode}
\coqdocemptyline
\end{coqdoccode}
\subsection{Evaluation as a Relation}



 Here's a better way: define \coqref{SECF.Hoare.HoareAssertAssume.ceval}{\coqdocinductive{ceval}} as a \textit{relation} rather than a
    \textit{function} -- i.e., make its result a \coqdockw{Prop} rather than a
    \coqref{SECF.SpecCT.state}{\coqdocdefinition{state}}, similar to what we did for \coqref{SECF.Imp.aevalR extended.aevalR}{\coqdocinductive{aevalR}} above. 

 This is an important change.  Besides freeing us from awkward
    workarounds, it gives us a ton more flexibility in the definition.
    For example, if we add nondeterministic features like \coqdocvar{any} to the
    language, we want the definition of evaluation to be
    nondeterministic -- i.e., not only will it not be total, it will
    not even be a function! 

 We'll use the notation \coqdocvar{st} =[ \coqdocvar{c} ]=> \coqdocvar{st'} for the \coqref{SECF.Hoare.HoareAssertAssume.ceval}{\coqdocinductive{ceval}} relation:
    \coqdocvar{st} =[ \coqdocvar{c} ]=> \coqdocvar{st'} means that executing program \coqdocvar{c} in a starting
    state \coqdocvar{st} results in an ending state \coqdocvar{st'}.  This can be
    pronounced ``\coqdocvar{c} takes state \coqdocvar{st} to \coqdocvar{st'}''. 
\begin{coqdoccode}
\coqdocemptyline
\end{coqdoccode}
\subsubsection{Operational Semantics}



 Here is an informal definition of evaluation, presented as inference
    rules for readability:
\par
\noindent\hrulefill\par
\noindent{}                            (E\_Skip)
                           st = \coqdocvar{skip} => st


                           aeval st a = n
\par
\noindent\hrulefill\par
\noindent{}                      (E\_Asgn)
                   st = \coqdocvar{x} := \coqdocvar{a} => (x !-> n ; st)


                           st  = \coqdocvar{c1} => st'
                           st' = \coqdocvar{c2} => st''
\par
\noindent\hrulefill\par
\noindent{}                           (E\_Seq)
                         st = \coqdocvar{c1};\coqdocvar{c2} => st''


                          beval st b = true
                           st = \coqdocvar{c1} => st'
\par
\noindent\hrulefill\par
\noindent{}               (E\_IfTrue)
                st = \coqdockw{if} \coqdocvar{b} \coqdockw{then} \coqdocvar{c1} \coqdockw{else} \coqdocvar{c2} \coqdockw{end} => st'


                         beval st b = false
                           st = \coqdocvar{c2} => st'
\par
\noindent\hrulefill\par
\noindent{}              (E\_IfFalse)
                st = \coqdockw{if} \coqdocvar{b} \coqdockw{then} \coqdocvar{c1} \coqdockw{else} \coqdocvar{c2} \coqdockw{end} => st'


                         beval st b = false
\par
\noindent\hrulefill\par
\noindent{}                 (E\_WhileFalse)
                    st = \coqdocvar{while} \coqdocvar{b} \coqdoctac{do} \coqdocvar{c} \coqdockw{end} => st


                          beval st b = true
                           st = \coqdocvar{c} => st'
                  st' = \coqdocvar{while} \coqdocvar{b} \coqdoctac{do} \coqdocvar{c} \coqdockw{end} => st''
\par
\noindent\hrulefill\par
\noindent{}                 (E\_WhileTrue)
                  st  = \coqdocvar{while} \coqdocvar{b} \coqdoctac{do} \coqdocvar{c} \coqdockw{end} => st''


 Here is the formal definition.  Make sure you understand
    how it corresponds to the inference rules. 
\begin{coqdoccode}
\coqdocemptyline
\coqdocnoindent
\coqdockw{Reserved Notation}\coqdoceol
\coqdocindent{4.50em}
"st0 '=[' c ']=>' st1"\coqdoceol
\coqdocindent{4.50em}
(\coqdoctac{at} \coqdockw{level} 40, \coqdocvar{c} \coqdocvar{custom} \coqdocvar{com} \coqdoctac{at} \coqdockw{level} 99,\coqdoceol
\coqdocindent{5.00em}
\coqdocvar{st0} \coqdockw{constr}, \coqdocvar{st1} \coqdockw{constr} \coqdoctac{at} \coqdocvar{next} \coqdockw{level},\coqdoceol
\coqdocindent{5.00em}
\coqdocvar{format} "'[hv' st0  =[ '/  ' '[' c ']' '/' ]=>  st1 ']'").\coqdoceol
\coqdocemptyline
\coqdocnoindent
\coqdockw{Inductive} \coqdef{SECF.Imp.ceval}{ceval}{\coqdocinductive{ceval}} : \coqref{SECF.Imp.com}{\coqdocinductive{com}} \coqexternalref{::type scope:x '->' x}{http://rocq-prover.org/doc/V9.1.1/stdlib//Corelib.Init.Logic}{\coqdocnotation{\ensuremath{\rightarrow}}} \coqref{SECF.Imp.state}{\coqdocdefinition{state}} \coqexternalref{::type scope:x '->' x}{http://rocq-prover.org/doc/V9.1.1/stdlib//Corelib.Init.Logic}{\coqdocnotation{\ensuremath{\rightarrow}}} \coqref{SECF.Imp.state}{\coqdocdefinition{state}} \coqexternalref{::type scope:x '->' x}{http://rocq-prover.org/doc/V9.1.1/stdlib//Corelib.Init.Logic}{\coqdocnotation{\ensuremath{\rightarrow}}} \coqdockw{Prop} :=\coqdoceol
\coqdocindent{1.00em}
\ensuremath{|} \coqdef{SECF.Imp.E Skip}{E\_Skip}{\coqdocconstructor{E\_Skip}} : \coqdockw{\ensuremath{\forall}} \coqdef{SECF.Imp.st:220}{st}{\coqdocbinder{st}},\coqdoceol
\coqdocindent{3.00em}
\coqref{SECF.Imp.st:220}{\coqdocvariable{st}} \coqref{SECF.Imp.:::x '=[' x ']=>' x}{\coqdocnotation{=[}} \coqref{SECF.Imp.:com:com scope:'skip'}{\coqdocnotation{skip}} \coqref{SECF.Imp.:::x '=[' x ']=>' x}{\coqdocnotation{]=>}} \coqref{SECF.Imp.st:220}{\coqdocvariable{st}}\coqdoceol
\coqdocindent{1.00em}
\ensuremath{|} \coqdef{SECF.Imp.E Asgn}{E\_Asgn}{\coqdocconstructor{E\_Asgn}}  : \coqdockw{\ensuremath{\forall}} \coqdef{SECF.Imp.st:221}{st}{\coqdocbinder{st}} \coqdef{SECF.Imp.a:222}{a}{\coqdocbinder{a}} \coqdef{SECF.Imp.n:223}{n}{\coqdocbinder{n}} \coqdef{SECF.Imp.x:224}{x}{\coqdocbinder{x}},\coqdoceol
\coqdocindent{3.00em}
\coqref{SECF.Imp.aeval}{\coqdocdefinition{aeval}} \coqref{SECF.Imp.st:221}{\coqdocvariable{st}} \coqref{SECF.Imp.a:222}{\coqdocvariable{a}} \coqexternalref{::type scope:x '=' x}{http://rocq-prover.org/doc/V9.1.1/stdlib//Corelib.Init.Logic}{\coqdocnotation{=}} \coqref{SECF.Imp.n:223}{\coqdocvariable{n}} \coqexternalref{::type scope:x '->' x}{http://rocq-prover.org/doc/V9.1.1/stdlib//Corelib.Init.Logic}{\coqdocnotation{\ensuremath{\rightarrow}}}\coqdoceol
\coqdocindent{3.00em}
\coqref{SECF.Imp.st:221}{\coqdocvariable{st}} \coqref{SECF.Imp.:::x '=[' x ']=>' x}{\coqdocnotation{=[}} \coqref{SECF.Imp.x:224}{\coqdocvariable{x}} \coqref{SECF.Imp.:com:com scope:x ':=' x}{\coqdocnotation{:=}} \coqref{SECF.Imp.a:222}{\coqdocvariable{a}} \coqref{SECF.Imp.:::x '=[' x ']=>' x}{\coqdocnotation{]=>}} \coqref{SECF.Imp.:::x '=[' x ']=>' x}{\coqdocnotation{(}}\coqref{SECF.Imp.x:224}{\coqdocvariable{x}} \coqref{SECF.Maps.:::x '!->' x ';' x}{\coqdocnotation{!->}} \coqref{SECF.Imp.n:223}{\coqdocvariable{n}} \coqref{SECF.Maps.:::x '!->' x ';' x}{\coqdocnotation{;}} \coqref{SECF.Imp.st:221}{\coqdocvariable{st}}\coqref{SECF.Imp.:::x '=[' x ']=>' x}{\coqdocnotation{)}}\coqdoceol
\coqdocindent{1.00em}
\ensuremath{|} \coqdef{SECF.Imp.E Seq}{E\_Seq}{\coqdocconstructor{E\_Seq}} : \coqdockw{\ensuremath{\forall}} \coqdef{SECF.Imp.c1:225}{c1}{\coqdocbinder{c1}} \coqdef{SECF.Imp.c2:226}{c2}{\coqdocbinder{c2}} \coqdef{SECF.Imp.st:227}{st}{\coqdocbinder{st}} \coqdef{SECF.Imp.st':228}{st'}{\coqdocbinder{st'}} \coqdef{SECF.Imp.st'':229}{st'{}'}{\coqdocbinder{st'{}'}},\coqdoceol
\coqdocindent{3.00em}
\coqref{SECF.Imp.st:227}{\coqdocvariable{st}}  \coqref{SECF.Imp.:::x '=[' x ']=>' x}{\coqdocnotation{=[}} \coqref{SECF.Imp.c1:225}{\coqdocvariable{c1}} \coqref{SECF.Imp.:::x '=[' x ']=>' x}{\coqdocnotation{]=>}} \coqref{SECF.Imp.st':228}{\coqdocvariable{st'}}  \coqexternalref{::type scope:x '->' x}{http://rocq-prover.org/doc/V9.1.1/stdlib//Corelib.Init.Logic}{\coqdocnotation{\ensuremath{\rightarrow}}}\coqdoceol
\coqdocindent{3.00em}
\coqref{SECF.Imp.st':228}{\coqdocvariable{st'}} \coqref{SECF.Imp.:::x '=[' x ']=>' x}{\coqdocnotation{=[}} \coqref{SECF.Imp.c2:226}{\coqdocvariable{c2}} \coqref{SECF.Imp.:::x '=[' x ']=>' x}{\coqdocnotation{]=>}} \coqref{SECF.Imp.st'':229}{\coqdocvariable{st'{}'}} \coqexternalref{::type scope:x '->' x}{http://rocq-prover.org/doc/V9.1.1/stdlib//Corelib.Init.Logic}{\coqdocnotation{\ensuremath{\rightarrow}}}\coqdoceol
\coqdocindent{3.00em}
\coqref{SECF.Imp.st:227}{\coqdocvariable{st}}  \coqref{SECF.Imp.:::x '=[' x ']=>' x}{\coqdocnotation{=[}} \coqref{SECF.Imp.c1:225}{\coqdocvariable{c1}} \coqref{SECF.Imp.:com:com scope:x ';' x}{\coqdocnotation{;}} \coqref{SECF.Imp.c2:226}{\coqdocvariable{c2}} \coqref{SECF.Imp.:::x '=[' x ']=>' x}{\coqdocnotation{]=>}} \coqref{SECF.Imp.st'':229}{\coqdocvariable{st'{}'}}\coqdoceol
\coqdocindent{1.00em}
\ensuremath{|} \coqdef{SECF.Imp.E IfTrue}{E\_IfTrue}{\coqdocconstructor{E\_IfTrue}} : \coqdockw{\ensuremath{\forall}} \coqdef{SECF.Imp.st:230}{st}{\coqdocbinder{st}} \coqdef{SECF.Imp.st':231}{st'}{\coqdocbinder{st'}} \coqdef{SECF.Imp.b:232}{b}{\coqdocbinder{b}} \coqdef{SECF.Imp.c1:233}{c1}{\coqdocbinder{c1}} \coqdef{SECF.Imp.c2:234}{c2}{\coqdocbinder{c2}},\coqdoceol
\coqdocindent{3.00em}
\coqref{SECF.Imp.beval}{\coqdocdefinition{beval}} \coqref{SECF.Imp.st:230}{\coqdocvariable{st}} \coqref{SECF.Imp.b:232}{\coqdocvariable{b}} \coqexternalref{::type scope:x '=' x}{http://rocq-prover.org/doc/V9.1.1/stdlib//Corelib.Init.Logic}{\coqdocnotation{=}} \coqref{SECF.Imp.:::'true'}{\coqdocnotation{true}} \coqexternalref{::type scope:x '->' x}{http://rocq-prover.org/doc/V9.1.1/stdlib//Corelib.Init.Logic}{\coqdocnotation{\ensuremath{\rightarrow}}}\coqdoceol
\coqdocindent{3.00em}
\coqref{SECF.Imp.st:230}{\coqdocvariable{st}} \coqref{SECF.Imp.:::x '=[' x ']=>' x}{\coqdocnotation{=[}} \coqref{SECF.Imp.c1:233}{\coqdocvariable{c1}} \coqref{SECF.Imp.:::x '=[' x ']=>' x}{\coqdocnotation{]=>}} \coqref{SECF.Imp.st':231}{\coqdocvariable{st'}} \coqexternalref{::type scope:x '->' x}{http://rocq-prover.org/doc/V9.1.1/stdlib//Corelib.Init.Logic}{\coqdocnotation{\ensuremath{\rightarrow}}}\coqdoceol
\coqdocindent{3.00em}
\coqref{SECF.Imp.st:230}{\coqdocvariable{st}} \coqref{SECF.Imp.:::x '=[' x ']=>' x}{\coqdocnotation{=[}} \coqref{SECF.Imp.:com:com scope:'if' x 'then' x 'else' x 'end'}{\coqdocnotation{if}} \coqref{SECF.Imp.b:232}{\coqdocvariable{b}} \coqref{SECF.Imp.:com:com scope:'if' x 'then' x 'else' x 'end'}{\coqdocnotation{then}} \coqref{SECF.Imp.c1:233}{\coqdocvariable{c1}} \coqref{SECF.Imp.:com:com scope:'if' x 'then' x 'else' x 'end'}{\coqdocnotation{else}} \coqref{SECF.Imp.c2:234}{\coqdocvariable{c2}} \coqref{SECF.Imp.:com:com scope:'if' x 'then' x 'else' x 'end'}{\coqdocnotation{end}}\coqref{SECF.Imp.:::x '=[' x ']=>' x}{\coqdocnotation{]=>}} \coqref{SECF.Imp.st':231}{\coqdocvariable{st'}}\coqdoceol
\coqdocindent{1.00em}
\ensuremath{|} \coqdef{SECF.Imp.E IfFalse}{E\_IfFalse}{\coqdocconstructor{E\_IfFalse}} : \coqdockw{\ensuremath{\forall}} \coqdef{SECF.Imp.st:235}{st}{\coqdocbinder{st}} \coqdef{SECF.Imp.st':236}{st'}{\coqdocbinder{st'}} \coqdef{SECF.Imp.b:237}{b}{\coqdocbinder{b}} \coqdef{SECF.Imp.c1:238}{c1}{\coqdocbinder{c1}} \coqdef{SECF.Imp.c2:239}{c2}{\coqdocbinder{c2}},\coqdoceol
\coqdocindent{3.00em}
\coqref{SECF.Imp.beval}{\coqdocdefinition{beval}} \coqref{SECF.Imp.st:235}{\coqdocvariable{st}} \coqref{SECF.Imp.b:237}{\coqdocvariable{b}} \coqexternalref{::type scope:x '=' x}{http://rocq-prover.org/doc/V9.1.1/stdlib//Corelib.Init.Logic}{\coqdocnotation{=}} \coqref{SECF.Imp.:::'false'}{\coqdocnotation{false}} \coqexternalref{::type scope:x '->' x}{http://rocq-prover.org/doc/V9.1.1/stdlib//Corelib.Init.Logic}{\coqdocnotation{\ensuremath{\rightarrow}}}\coqdoceol
\coqdocindent{3.00em}
\coqref{SECF.Imp.st:235}{\coqdocvariable{st}} \coqref{SECF.Imp.:::x '=[' x ']=>' x}{\coqdocnotation{=[}} \coqref{SECF.Imp.c2:239}{\coqdocvariable{c2}} \coqref{SECF.Imp.:::x '=[' x ']=>' x}{\coqdocnotation{]=>}} \coqref{SECF.Imp.st':236}{\coqdocvariable{st'}} \coqexternalref{::type scope:x '->' x}{http://rocq-prover.org/doc/V9.1.1/stdlib//Corelib.Init.Logic}{\coqdocnotation{\ensuremath{\rightarrow}}}\coqdoceol
\coqdocindent{3.00em}
\coqref{SECF.Imp.st:235}{\coqdocvariable{st}} \coqref{SECF.Imp.:::x '=[' x ']=>' x}{\coqdocnotation{=[}} \coqref{SECF.Imp.:com:com scope:'if' x 'then' x 'else' x 'end'}{\coqdocnotation{if}} \coqref{SECF.Imp.b:237}{\coqdocvariable{b}} \coqref{SECF.Imp.:com:com scope:'if' x 'then' x 'else' x 'end'}{\coqdocnotation{then}} \coqref{SECF.Imp.c1:238}{\coqdocvariable{c1}} \coqref{SECF.Imp.:com:com scope:'if' x 'then' x 'else' x 'end'}{\coqdocnotation{else}} \coqref{SECF.Imp.c2:239}{\coqdocvariable{c2}} \coqref{SECF.Imp.:com:com scope:'if' x 'then' x 'else' x 'end'}{\coqdocnotation{end}}\coqref{SECF.Imp.:::x '=[' x ']=>' x}{\coqdocnotation{]=>}} \coqref{SECF.Imp.st':236}{\coqdocvariable{st'}}\coqdoceol
\coqdocindent{1.00em}
\ensuremath{|} \coqdef{SECF.Imp.E WhileFalse}{E\_WhileFalse}{\coqdocconstructor{E\_WhileFalse}} : \coqdockw{\ensuremath{\forall}} \coqdef{SECF.Imp.b:240}{b}{\coqdocbinder{b}} \coqdef{SECF.Imp.st:241}{st}{\coqdocbinder{st}} \coqdef{SECF.Imp.c:242}{c}{\coqdocbinder{c}},\coqdoceol
\coqdocindent{3.00em}
\coqref{SECF.Imp.beval}{\coqdocdefinition{beval}} \coqref{SECF.Imp.st:241}{\coqdocvariable{st}} \coqref{SECF.Imp.b:240}{\coqdocvariable{b}} \coqexternalref{::type scope:x '=' x}{http://rocq-prover.org/doc/V9.1.1/stdlib//Corelib.Init.Logic}{\coqdocnotation{=}} \coqref{SECF.Imp.:::'false'}{\coqdocnotation{false}} \coqexternalref{::type scope:x '->' x}{http://rocq-prover.org/doc/V9.1.1/stdlib//Corelib.Init.Logic}{\coqdocnotation{\ensuremath{\rightarrow}}}\coqdoceol
\coqdocindent{3.00em}
\coqref{SECF.Imp.st:241}{\coqdocvariable{st}} \coqref{SECF.Imp.:::x '=[' x ']=>' x}{\coqdocnotation{=[}} \coqref{SECF.Imp.:com:com scope:'while' x 'do' x 'end'}{\coqdocnotation{while}} \coqref{SECF.Imp.b:240}{\coqdocvariable{b}} \coqref{SECF.Imp.:com:com scope:'while' x 'do' x 'end'}{\coqdocnotation{do}} \coqref{SECF.Imp.c:242}{\coqdocvariable{c}} \coqref{SECF.Imp.:com:com scope:'while' x 'do' x 'end'}{\coqdocnotation{end}} \coqref{SECF.Imp.:::x '=[' x ']=>' x}{\coqdocnotation{]=>}} \coqref{SECF.Imp.st:241}{\coqdocvariable{st}}\coqdoceol
\coqdocindent{1.00em}
\ensuremath{|} \coqdef{SECF.Imp.E WhileTrue}{E\_WhileTrue}{\coqdocconstructor{E\_WhileTrue}} : \coqdockw{\ensuremath{\forall}} \coqdef{SECF.Imp.st:243}{st}{\coqdocbinder{st}} \coqdef{SECF.Imp.st':244}{st'}{\coqdocbinder{st'}} \coqdef{SECF.Imp.st'':245}{st'{}'}{\coqdocbinder{st'{}'}} \coqdef{SECF.Imp.b:246}{b}{\coqdocbinder{b}} \coqdef{SECF.Imp.c:247}{c}{\coqdocbinder{c}},\coqdoceol
\coqdocindent{3.00em}
\coqref{SECF.Imp.beval}{\coqdocdefinition{beval}} \coqref{SECF.Imp.st:243}{\coqdocvariable{st}} \coqref{SECF.Imp.b:246}{\coqdocvariable{b}} \coqexternalref{::type scope:x '=' x}{http://rocq-prover.org/doc/V9.1.1/stdlib//Corelib.Init.Logic}{\coqdocnotation{=}} \coqref{SECF.Imp.:::'true'}{\coqdocnotation{true}} \coqexternalref{::type scope:x '->' x}{http://rocq-prover.org/doc/V9.1.1/stdlib//Corelib.Init.Logic}{\coqdocnotation{\ensuremath{\rightarrow}}}\coqdoceol
\coqdocindent{3.00em}
\coqref{SECF.Imp.st:243}{\coqdocvariable{st}}  \coqref{SECF.Imp.:::x '=[' x ']=>' x}{\coqdocnotation{=[}} \coqref{SECF.Imp.c:247}{\coqdocvariable{c}} \coqref{SECF.Imp.:::x '=[' x ']=>' x}{\coqdocnotation{]=>}} \coqref{SECF.Imp.st':244}{\coqdocvariable{st'}} \coqexternalref{::type scope:x '->' x}{http://rocq-prover.org/doc/V9.1.1/stdlib//Corelib.Init.Logic}{\coqdocnotation{\ensuremath{\rightarrow}}}\coqdoceol
\coqdocindent{3.00em}
\coqref{SECF.Imp.st':244}{\coqdocvariable{st'}} \coqref{SECF.Imp.:::x '=[' x ']=>' x}{\coqdocnotation{=[}} \coqref{SECF.Imp.:com:com scope:'while' x 'do' x 'end'}{\coqdocnotation{while}} \coqref{SECF.Imp.b:246}{\coqdocvariable{b}} \coqref{SECF.Imp.:com:com scope:'while' x 'do' x 'end'}{\coqdocnotation{do}} \coqref{SECF.Imp.c:247}{\coqdocvariable{c}} \coqref{SECF.Imp.:com:com scope:'while' x 'do' x 'end'}{\coqdocnotation{end}} \coqref{SECF.Imp.:::x '=[' x ']=>' x}{\coqdocnotation{]=>}} \coqref{SECF.Imp.st'':245}{\coqdocvariable{st'{}'}} \coqexternalref{::type scope:x '->' x}{http://rocq-prover.org/doc/V9.1.1/stdlib//Corelib.Init.Logic}{\coqdocnotation{\ensuremath{\rightarrow}}}\coqdoceol
\coqdocindent{3.00em}
\coqref{SECF.Imp.st:243}{\coqdocvariable{st}}  \coqref{SECF.Imp.:::x '=[' x ']=>' x}{\coqdocnotation{=[}} \coqref{SECF.Imp.:com:com scope:'while' x 'do' x 'end'}{\coqdocnotation{while}} \coqref{SECF.Imp.b:246}{\coqdocvariable{b}} \coqref{SECF.Imp.:com:com scope:'while' x 'do' x 'end'}{\coqdocnotation{do}} \coqref{SECF.Imp.c:247}{\coqdocvariable{c}} \coqref{SECF.Imp.:com:com scope:'while' x 'do' x 'end'}{\coqdocnotation{end}} \coqref{SECF.Imp.:::x '=[' x ']=>' x}{\coqdocnotation{]=>}} \coqref{SECF.Imp.st'':245}{\coqdocvariable{st'{}'}}\coqdoceol
\coqdocnoindent
\coqdoceol
\coqdocindent{1.00em}
\coqdockw{where} \coqdef{SECF.Imp.:::x '=[' x ']=>' x}{"}{"}st0 =[ c ]=> st1" := (\coqref{SECF.Imp.ceval:219}{\coqdocinductive{ceval}} \coqdocvar{c} \coqdocvar{st0} \coqdocvar{st1}).\coqdoceol
\coqdocemptyline
\end{coqdoccode}
The cost of defining evaluation as a relation instead of a
    function is that we now need to construct a \textit{proof} that some
    program evaluates to some result state, rather than just letting
    Rocq's computation mechanism do it for us. 
\begin{coqdoccode}
\coqdocemptyline
\coqdocnoindent
\coqdockw{Example} \coqdef{SECF.Imp.ceval example1}{ceval\_example1}{\coqdocdefinition{ceval\_example1}}:\coqdoceol
\coqdocindent{1.00em}
\coqref{SECF.Imp.empty st}{\coqdocdefinition{empty\_st}} \coqref{SECF.Imp.:::x '=[' x ']=>' x}{\coqdocnotation{=[}}\coqdoceol
\coqdocindent{2.50em}
\coqref{SECF.Imp.X}{\coqdocdefinition{X}} \coqref{SECF.Imp.:com:com scope:x ':=' x}{\coqdocnotation{:=}} 2\coqref{SECF.Imp.:com:com scope:x ';' x}{\coqdocnotation{;}}\coqdoceol
\coqdocindent{2.50em}
\coqref{SECF.Imp.:com:com scope:'if' x 'then' x 'else' x 'end'}{\coqdocnotation{if}} \coqref{SECF.Imp.:com::'(' x ')'}{\coqdocnotation{(}}\coqref{SECF.Imp.X}{\coqdocdefinition{X}} \coqref{SECF.Imp.:com::x '<=' x}{\coqdocnotation{\ensuremath{\le}}} 1\coqref{SECF.Imp.:com::'(' x ')'}{\coqdocnotation{)}}\coqdoceol
\coqdocindent{3.50em}
\coqref{SECF.Imp.:com:com scope:'if' x 'then' x 'else' x 'end'}{\coqdocnotation{then}} \coqref{SECF.Imp.Y}{\coqdocdefinition{Y}} \coqref{SECF.Imp.:com:com scope:x ':=' x}{\coqdocnotation{:=}} 3\coqdoceol
\coqdocindent{3.50em}
\coqref{SECF.Imp.:com:com scope:'if' x 'then' x 'else' x 'end'}{\coqdocnotation{else}} \coqref{SECF.Imp.Z}{\coqdocdefinition{Z}} \coqref{SECF.Imp.:com:com scope:x ':=' x}{\coqdocnotation{:=}} 4\coqdoceol
\coqdocindent{2.50em}
\coqref{SECF.Imp.:com:com scope:'if' x 'then' x 'else' x 'end'}{\coqdocnotation{end}}\coqdoceol
\coqdocindent{1.00em}
\coqref{SECF.Imp.:::x '=[' x ']=>' x}{\coqdocnotation{]=>}} \coqref{SECF.Imp.:::x '=[' x ']=>' x}{\coqdocnotation{(}}\coqref{SECF.Imp.Z}{\coqdocdefinition{Z}} \coqref{SECF.Maps.:::x '!->' x ';' x}{\coqdocnotation{!->}} 4 \coqref{SECF.Maps.:::x '!->' x ';' x}{\coqdocnotation{;}} \coqref{SECF.Imp.X}{\coqdocdefinition{X}} \coqref{SECF.Imp.:::x '!->' x}{\coqdocnotation{!->}} 2\coqref{SECF.Imp.:::x '=[' x ']=>' x}{\coqdocnotation{)}}.\coqdoceol
\coqdocnoindent
\coqdockw{Proof}.\coqdoceol
\coqdocindent{1.00em}
\coqdoctac{apply} \coqref{SECF.Imp.E Seq}{\coqdocconstructor{E\_Seq}} \coqdockw{with} (\coqref{SECF.Imp.X}{\coqdocdefinition{X}} \coqref{SECF.Imp.:::x '!->' x}{\coqdocnotation{!->}} 2).\coqdoceol
\coqdocindent{1.00em}
- \coqdoceol
\coqdocindent{2.00em}
\coqdoctac{apply} \coqref{SECF.Imp.E Asgn}{\coqdocconstructor{E\_Asgn}}. \coqdoctac{reflexivity}.\coqdoceol
\coqdocindent{1.00em}
- \coqdoceol
\coqdocindent{2.00em}
\coqdoctac{apply} \coqref{SECF.Imp.E IfFalse}{\coqdocconstructor{E\_IfFalse}}.\coqdoceol
\coqdocindent{2.00em}
+ \coqdoctac{reflexivity}.\coqdoceol
\coqdocindent{2.00em}
+ \coqdoctac{apply} \coqref{SECF.Imp.E Asgn}{\coqdocconstructor{E\_Asgn}}. \coqdoctac{reflexivity}.\coqdoceol
\coqdocnoindent
\coqdockw{Qed}.\coqdoceol
\coqdocemptyline
\end{coqdoccode}
\paragraph{Exercise: 2 stars, standard (ceval\_example2)}


\begin{coqdoccode}
\coqdocnoindent
\coqdockw{Example} \coqdef{SECF.Imp.ceval example2}{ceval\_example2}{\coqdocdefinition{ceval\_example2}}:\coqdoceol
\coqdocindent{1.00em}
\coqref{SECF.Imp.empty st}{\coqdocdefinition{empty\_st}} \coqref{SECF.Imp.:::x '=[' x ']=>' x}{\coqdocnotation{=[}}\coqdoceol
\coqdocindent{2.00em}
\coqref{SECF.Imp.X}{\coqdocdefinition{X}} \coqref{SECF.Imp.:com:com scope:x ':=' x}{\coqdocnotation{:=}} 0\coqref{SECF.Imp.:com:com scope:x ';' x}{\coqdocnotation{;}}\coqdoceol
\coqdocindent{2.00em}
\coqref{SECF.Imp.Y}{\coqdocdefinition{Y}} \coqref{SECF.Imp.:com:com scope:x ':=' x}{\coqdocnotation{:=}} 1\coqref{SECF.Imp.:com:com scope:x ';' x}{\coqdocnotation{;}}\coqdoceol
\coqdocindent{2.00em}
\coqref{SECF.Imp.Z}{\coqdocdefinition{Z}} \coqref{SECF.Imp.:com:com scope:x ':=' x}{\coqdocnotation{:=}} 2\coqdoceol
\coqdocindent{1.00em}
\coqref{SECF.Imp.:::x '=[' x ']=>' x}{\coqdocnotation{]=>}} \coqref{SECF.Imp.:::x '=[' x ']=>' x}{\coqdocnotation{(}}\coqref{SECF.Imp.Z}{\coqdocdefinition{Z}} \coqref{SECF.Maps.:::x '!->' x ';' x}{\coqdocnotation{!->}} 2 \coqref{SECF.Maps.:::x '!->' x ';' x}{\coqdocnotation{;}} \coqref{SECF.Imp.Y}{\coqdocdefinition{Y}} \coqref{SECF.Maps.:::x '!->' x ';' x}{\coqdocnotation{!->}} 1 \coqref{SECF.Maps.:::x '!->' x ';' x}{\coqdocnotation{;}} \coqref{SECF.Imp.X}{\coqdocdefinition{X}} \coqref{SECF.Imp.:::x '!->' x}{\coqdocnotation{!->}} 0\coqref{SECF.Imp.:::x '=[' x ']=>' x}{\coqdocnotation{)}}.\coqdoceol
\coqdocnoindent
\coqdockw{Proof}.\coqdoceol
 \coqdocvar{Admitted}.\coqdoceol
\end{coqdoccode}
\ensuremath{\Box} 
\begin{coqdoccode}
\coqdocemptyline
\coqdocnoindent
\coqdockw{Set} \coqdockw{Printing} \coqdockw{Implicit}.\coqdoceol
\coqdocnoindent
\coqdockw{Check} @\coqref{SECF.Imp.ceval example2}{\coqdocaxiom{ceval\_example2}}.\coqdoceol
\coqdocemptyline
\end{coqdoccode}
\paragraph{Exercise: 3 stars, standard, optional (pup\_to\_n)}




    Write an Imp program that sums the numbers from 1 to \coqref{SECF.SpecCT.X}{\coqdocdefinition{X}}
    (inclusive: 1 + 2 + ... + \coqref{SECF.SpecCT.X}{\coqdocdefinition{X}}) in the variable \coqref{SECF.SpecCT.Y}{\coqdocdefinition{Y}}.  Your program
    should update the state as shown in theorem \coqref{SECF.Imp.pup to 2 ceval}{\coqdoclemma{pup\_to\_2\_ceval}},
    which you can reverse-engineer to discover the program you should
    write.  The proof of that theorem will be somewhat lengthy. 
\begin{coqdoccode}
\coqdocemptyline
\coqdocnoindent
\coqdockw{Definition} \coqdef{SECF.Imp.pup to n}{pup\_to\_n}{\coqdocdefinition{pup\_to\_n}} : \coqref{SECF.Imp.com}{\coqdocinductive{com}}\coqdoceol
\coqdocindent{1.00em}
. \coqdocvar{Admitted}.\coqdoceol
\coqdocemptyline
\coqdocnoindent
\coqdockw{Theorem} \coqdef{SECF.Imp.pup to 2 ceval}{pup\_to\_2\_ceval}{\coqdoclemma{pup\_to\_2\_ceval}} :\coqdoceol
\coqdocindent{1.00em}
\coqref{SECF.Imp.:::x '=[' x ']=>' x}{\coqdocnotation{(}}\coqref{SECF.Imp.X}{\coqdocdefinition{X}} \coqref{SECF.Imp.:::x '!->' x}{\coqdocnotation{!->}} 2\coqref{SECF.Imp.:::x '=[' x ']=>' x}{\coqdocnotation{)}} \coqref{SECF.Imp.:::x '=[' x ']=>' x}{\coqdocnotation{=[}}\coqdoceol
\coqdocindent{2.00em}
\coqref{SECF.Imp.pup to n}{\coqdocaxiom{pup\_to\_n}}\coqdoceol
\coqdocindent{1.00em}
\coqref{SECF.Imp.:::x '=[' x ']=>' x}{\coqdocnotation{]=>}} \coqref{SECF.Imp.:::x '=[' x ']=>' x}{\coqdocnotation{(}}\coqref{SECF.Imp.X}{\coqdocdefinition{X}} \coqref{SECF.Maps.:::x '!->' x ';' x}{\coqdocnotation{!->}} 0 \coqref{SECF.Maps.:::x '!->' x ';' x}{\coqdocnotation{;}} \coqref{SECF.Imp.Y}{\coqdocdefinition{Y}} \coqref{SECF.Maps.:::x '!->' x ';' x}{\coqdocnotation{!->}} 3 \coqref{SECF.Maps.:::x '!->' x ';' x}{\coqdocnotation{;}} \coqref{SECF.Imp.X}{\coqdocdefinition{X}} \coqref{SECF.Maps.:::x '!->' x ';' x}{\coqdocnotation{!->}} 1 \coqref{SECF.Maps.:::x '!->' x ';' x}{\coqdocnotation{;}} \coqref{SECF.Imp.Y}{\coqdocdefinition{Y}} \coqref{SECF.Maps.:::x '!->' x ';' x}{\coqdocnotation{!->}} 2 \coqref{SECF.Maps.:::x '!->' x ';' x}{\coqdocnotation{;}} \coqref{SECF.Imp.Y}{\coqdocdefinition{Y}} \coqref{SECF.Maps.:::x '!->' x ';' x}{\coqdocnotation{!->}} 0 \coqref{SECF.Maps.:::x '!->' x ';' x}{\coqdocnotation{;}} \coqref{SECF.Imp.X}{\coqdocdefinition{X}} \coqref{SECF.Imp.:::x '!->' x}{\coqdocnotation{!->}} 2\coqref{SECF.Imp.:::x '=[' x ']=>' x}{\coqdocnotation{)}}.\coqdoceol
\coqdocnoindent
\coqdockw{Proof}.\coqdoceol
 \coqdocvar{Admitted}.\coqdoceol
\end{coqdoccode}
\ensuremath{\Box} 
\begin{coqdoccode}
\coqdocemptyline
\end{coqdoccode}
\subsection{Determinism of Evaluation}



 Changing from a computational to a relational definition of
    evaluation is a good move because it frees us from the artificial
    requirement that evaluation should be a total function.  But it
    also raises a question: Is the second definition of evaluation
    really a partial \textit{function}?  Or is it possible that, beginning from
    the same state \coqdocvar{st}, we could evaluate some command \coqdocvar{c} in
    different ways to reach two different output states \coqdocvar{st'} and
    \coqdocvar{st'{}'}?


    In fact, this cannot happen: \coqref{SECF.Hoare.HoareAssertAssume.ceval}{\coqdocinductive{ceval}} \textit{is} a partial function: 
\begin{coqdoccode}
\coqdocemptyline
\coqdocnoindent
\coqdockw{Theorem} \coqdef{SECF.Imp.ceval deterministic}{ceval\_deterministic}{\coqdoclemma{ceval\_deterministic}}: \coqdockw{\ensuremath{\forall}} \coqdef{SECF.Imp.c:248}{c}{\coqdocbinder{c}} \coqdef{SECF.Imp.st:249}{st}{\coqdocbinder{st}} \coqdef{SECF.Imp.st1:250}{st1}{\coqdocbinder{st1}} \coqdef{SECF.Imp.st2:251}{st2}{\coqdocbinder{st2}},\coqdoceol
\coqdocindent{2.50em}
\coqref{SECF.Imp.st:249}{\coqdocvariable{st}} \coqref{SECF.Imp.:::x '=[' x ']=>' x}{\coqdocnotation{=[}} \coqref{SECF.Imp.c:248}{\coqdocvariable{c}} \coqref{SECF.Imp.:::x '=[' x ']=>' x}{\coqdocnotation{]=>}} \coqref{SECF.Imp.st1:250}{\coqdocvariable{st1}}  \coqexternalref{::type scope:x '->' x}{http://rocq-prover.org/doc/V9.1.1/stdlib//Corelib.Init.Logic}{\coqdocnotation{\ensuremath{\rightarrow}}}\coqdoceol
\coqdocindent{2.50em}
\coqref{SECF.Imp.st:249}{\coqdocvariable{st}} \coqref{SECF.Imp.:::x '=[' x ']=>' x}{\coqdocnotation{=[}} \coqref{SECF.Imp.c:248}{\coqdocvariable{c}} \coqref{SECF.Imp.:::x '=[' x ']=>' x}{\coqdocnotation{]=>}} \coqref{SECF.Imp.st2:251}{\coqdocvariable{st2}} \coqexternalref{::type scope:x '->' x}{http://rocq-prover.org/doc/V9.1.1/stdlib//Corelib.Init.Logic}{\coqdocnotation{\ensuremath{\rightarrow}}}\coqdoceol
\coqdocindent{2.50em}
\coqref{SECF.Imp.st1:250}{\coqdocvariable{st1}} \coqexternalref{::type scope:x '=' x}{http://rocq-prover.org/doc/V9.1.1/stdlib//Corelib.Init.Logic}{\coqdocnotation{=}} \coqref{SECF.Imp.st2:251}{\coqdocvariable{st2}}.\coqdoceol
\coqdocnoindent
\coqdockw{Proof}.\coqdoceol
\coqdocindent{1.00em}
\coqdoctac{intros} \coqdocvar{c} \coqdocvar{st} \coqdocvar{st1} \coqdocvar{st2} \coqdocvar{E1} \coqdocvar{E2}.\coqdoceol
\coqdocindent{1.00em}
\coqdoctac{generalize} \coqdoctac{dependent} \coqdocvar{st2}.\coqdoceol
\coqdocindent{1.00em}
\coqdoctac{induction} \coqdocvar{E1}; \coqdoctac{intros} \coqdocvar{st2} \coqdocvar{E2}; \coqdoctac{inversion} \coqdocvar{E2}; \coqdoctac{subst}.\coqdoceol
\coqdocindent{1.00em}
-  \coqdoctac{reflexivity}.\coqdoceol
\coqdocindent{1.00em}
-  \coqdoctac{reflexivity}.\coqdoceol
\coqdocindent{1.00em}
- \coqdoceol
\coqdocindent{2.00em}
\coqdoctac{rewrite} (\coqdocvar{IHE1\_1} \coqdocvar{st'0} \coqdocvar{H1}) \coqdoctac{in} *.\coqdoceol
\coqdocindent{2.00em}
\coqdoctac{apply} \coqdocvar{IHE1\_2}. \coqdoctac{assumption}.\coqdoceol
\coqdocindent{1.00em}
- \coqdoceol
\coqdocindent{3.00em}
\coqdoctac{apply} \coqdocvar{IHE1}. \coqdoctac{assumption}.\coqdoceol
\coqdocindent{1.00em}
- \coqdoceol
\coqdocindent{3.00em}
\coqdoctac{rewrite} \coqdocvar{H} \coqdoctac{in} \coqdocvar{H5}. \coqdoctac{discriminate}.\coqdoceol
\coqdocindent{1.00em}
- \coqdoceol
\coqdocindent{3.00em}
\coqdoctac{rewrite} \coqdocvar{H} \coqdoctac{in} \coqdocvar{H5}. \coqdoctac{discriminate}.\coqdoceol
\coqdocindent{1.00em}
- \coqdoceol
\coqdocindent{3.00em}
\coqdoctac{apply} \coqdocvar{IHE1}. \coqdoctac{assumption}.\coqdoceol
\coqdocindent{1.00em}
- \coqdoceol
\coqdocindent{2.00em}
\coqdoctac{reflexivity}.\coqdoceol
\coqdocindent{1.00em}
- \coqdoceol
\coqdocindent{2.00em}
\coqdoctac{rewrite} \coqdocvar{H} \coqdoctac{in} \coqdocvar{H2}. \coqdoctac{discriminate}.\coqdoceol
\coqdocindent{1.00em}
- \coqdoceol
\coqdocindent{2.00em}
\coqdoctac{rewrite} \coqdocvar{H} \coqdoctac{in} \coqdocvar{H4}. \coqdoctac{discriminate}.\coqdoceol
\coqdocindent{1.00em}
- \coqdoceol
\coqdocindent{2.00em}
\coqdoctac{rewrite} (\coqdocvar{IHE1\_1} \coqdocvar{st'0} \coqdocvar{H3}) \coqdoctac{in} *.\coqdoceol
\coqdocindent{2.00em}
\coqdoctac{apply} \coqdocvar{IHE1\_2}. \coqdoctac{assumption}. \coqdockw{Qed}.\coqdoceol
\coqdocemptyline
\end{coqdoccode}
\section{Reasoning About Imp Programs}



 We'll get into more systematic and powerful techniques for
    reasoning about Imp programs in \textit{Programming Language
    Foundations}, but we can already do a few things (albeit in a
    somewhat low-level way) just by working with the bare definitions.
    This section explores some examples. 
\begin{coqdoccode}
\coqdocemptyline
\coqdocnoindent
\coqdockw{Theorem} \coqdef{SECF.Imp.plus2 spec}{plus2\_spec}{\coqdoclemma{plus2\_spec}} : \coqdockw{\ensuremath{\forall}} \coqdef{SECF.Imp.st:252}{st}{\coqdocbinder{st}} \coqdef{SECF.Imp.n:253}{n}{\coqdocbinder{n}} \coqdef{SECF.Imp.st':254}{st'}{\coqdocbinder{st'}},\coqdoceol
\coqdocindent{1.00em}
\coqref{SECF.Imp.st:252}{\coqdocvariable{st}} \coqref{SECF.Imp.X}{\coqdocdefinition{X}} \coqexternalref{::type scope:x '=' x}{http://rocq-prover.org/doc/V9.1.1/stdlib//Corelib.Init.Logic}{\coqdocnotation{=}} \coqref{SECF.Imp.n:253}{\coqdocvariable{n}} \coqexternalref{::type scope:x '->' x}{http://rocq-prover.org/doc/V9.1.1/stdlib//Corelib.Init.Logic}{\coqdocnotation{\ensuremath{\rightarrow}}}\coqdoceol
\coqdocindent{1.00em}
\coqref{SECF.Imp.st:252}{\coqdocvariable{st}} \coqref{SECF.Imp.:::x '=[' x ']=>' x}{\coqdocnotation{=[}} \coqref{SECF.Imp.plus2}{\coqdocdefinition{plus2}} \coqref{SECF.Imp.:::x '=[' x ']=>' x}{\coqdocnotation{]=>}} \coqref{SECF.Imp.st':254}{\coqdocvariable{st'}} \coqexternalref{::type scope:x '->' x}{http://rocq-prover.org/doc/V9.1.1/stdlib//Corelib.Init.Logic}{\coqdocnotation{\ensuremath{\rightarrow}}}\coqdoceol
\coqdocindent{1.00em}
\coqref{SECF.Imp.st':254}{\coqdocvariable{st'}} \coqref{SECF.Imp.X}{\coqdocdefinition{X}} \coqexternalref{::type scope:x '=' x}{http://rocq-prover.org/doc/V9.1.1/stdlib//Corelib.Init.Logic}{\coqdocnotation{=}} \coqref{SECF.Imp.n:253}{\coqdocvariable{n}} \coqexternalref{::nat scope:x '+' x}{http://rocq-prover.org/doc/V9.1.1/stdlib//Corelib.Init.Nat}{\coqdocnotation{+}} 2.\coqdoceol
\coqdocnoindent
\coqdockw{Proof}.\coqdoceol
\coqdocindent{1.00em}
\coqdoctac{intros} \coqdocvar{st} \coqdocvar{n} \coqdocvar{st'} \coqdocvar{HX} \coqdocvar{Heval}.\coqdoceol
\coqdocemptyline
\end{coqdoccode}
Inverting \coqdocvar{Heval} essentially forces Rocq to expand one step of
      the \coqref{SECF.Hoare.HoareAssertAssume.ceval}{\coqdocinductive{ceval}} computation -- in this case revealing that \coqdocvar{st'}
      must be \coqdocvar{st} extended with the new value of \coqref{SECF.SpecCT.X}{\coqdocdefinition{X}}, since \coqref{SECF.Imp.plus2}{\coqdocdefinition{plus2}}
      is an assignment. 
\begin{coqdoccode}
\coqdocemptyline
\coqdocindent{1.00em}
\coqdoctac{inversion} \coqdocvar{Heval}. \coqdoctac{subst}. \coqdoctac{clear} \coqdocvar{Heval}. \coqdoctac{simpl}.\coqdoceol
\coqdocindent{1.00em}
\coqdoctac{apply} \coqref{SECF.Maps.t update eq}{\coqdocaxiom{t\_update\_eq}}. \coqdockw{Qed}.\coqdoceol
\coqdocemptyline
\end{coqdoccode}
\paragraph{Exercise: 3 stars, standard, optional (XtimesYinZ\_spec)}




    State and prove a specification of \coqref{SECF.Imp.XtimesYinZ}{\coqdocdefinition{XtimesYinZ}}. 
\begin{coqdoccode}
\coqdocemptyline
\coqdocemptyline
\coqdocnoindent
\coqdockw{Definition} \coqdef{SECF.Imp.manual grade for XtimesYinZ spec}{manual\_grade\_for\_XtimesYinZ\_spec}{\coqdocdefinition{manual\_grade\_for\_XtimesYinZ\_spec}} : \coqexternalref{option}{http://rocq-prover.org/doc/V9.1.1/stdlib//Corelib.Init.Datatypes}{\coqdocinductive{option}} (\coqexternalref{nat}{http://rocq-prover.org/doc/V9.1.1/stdlib//Corelib.Init.Datatypes}{\coqdocinductive{nat}}\coqexternalref{::type scope:x '*' x}{http://rocq-prover.org/doc/V9.1.1/stdlib//Corelib.Init.Datatypes}{\coqdocnotation{\ensuremath{\times}}}\coqdocinductive{string}) := \coqexternalref{None}{http://rocq-prover.org/doc/V9.1.1/stdlib//Corelib.Init.Datatypes}{\coqdocconstructor{None}}.\coqdoceol
\end{coqdoccode}
\ensuremath{\Box} 

\paragraph{Exercise: 3 stars, standard, especially useful (loop\_never\_stops)}


\begin{coqdoccode}
\coqdocnoindent
\coqdockw{Theorem} \coqdef{SECF.Imp.loop never stops}{loop\_never\_stops}{\coqdoclemma{loop\_never\_stops}} : \coqdockw{\ensuremath{\forall}} \coqdef{SECF.Imp.st:255}{st}{\coqdocbinder{st}} \coqdef{SECF.Imp.st':256}{st'}{\coqdocbinder{st'}},\coqdoceol
\coqdocindent{1.00em}
\coqexternalref{::type scope:'x7E' x}{http://rocq-prover.org/doc/V9.1.1/stdlib//Corelib.Init.Logic}{\coqdocnotation{\~{}(}}\coqref{SECF.Imp.st:255}{\coqdocvariable{st}} \coqref{SECF.Imp.:::x '=[' x ']=>' x}{\coqdocnotation{=[}} \coqref{SECF.Imp.loop}{\coqdocdefinition{loop}} \coqref{SECF.Imp.:::x '=[' x ']=>' x}{\coqdocnotation{]=>}} \coqref{SECF.Imp.st':256}{\coqdocvariable{st'}}\coqexternalref{::type scope:'x7E' x}{http://rocq-prover.org/doc/V9.1.1/stdlib//Corelib.Init.Logic}{\coqdocnotation{)}}.\coqdoceol
\coqdocnoindent
\coqdockw{Proof}.\coqdoceol
\coqdocindent{1.00em}
\coqdoctac{intros} \coqdocvar{st} \coqdocvar{st'} \coqdocvar{contra}. \coqdoctac{unfold} \coqref{SECF.Imp.loop}{\coqdocdefinition{loop}} \coqdoctac{in} \coqdocvar{contra}.\coqdoceol
\coqdocindent{1.00em}
\coqdocvar{remember} \coqref{SECF.Imp.::com scope:'<x7B' x 'x7D>'}{\coqdocnotation{<\{}} \coqref{SECF.Imp.:com:com scope:'while' x 'do' x 'end'}{\coqdocnotation{while}} \coqref{SECF.Imp.:com::'true'}{\coqdocnotation{true}} \coqref{SECF.Imp.:com:com scope:'while' x 'do' x 'end'}{\coqdocnotation{do}} \coqref{SECF.Imp.:com:com scope:'skip'}{\coqdocnotation{skip}} \coqref{SECF.Imp.:com:com scope:'while' x 'do' x 'end'}{\coqdocnotation{end}} \coqref{SECF.Imp.::com scope:'<x7B' x 'x7D>'}{\coqdocnotation{\}>}} \coqdockw{as} \coqdocvar{loopdef}\coqdoceol
\coqdocindent{5.50em}
\coqdocvar{eqn}:\coqdocvar{Heqloopdef}.\coqdoceol
\coqdocemptyline
\end{coqdoccode}
Proceed by induction on the assumed derivation showing that
      \coqdocvar{loopdef} terminates.  Most of the cases are immediately
      contradictory and so can be solved in one step with
      \coqdoctac{discriminate}. 
\begin{coqdoccode}
\coqdocemptyline
 \coqdocvar{Admitted}.\coqdoceol
\end{coqdoccode}
\ensuremath{\Box} 

\paragraph{Exercise: 3 stars, standard (no\_whiles\_eqv)}




    Consider the following function: 
\begin{coqdoccode}
\coqdocemptyline
\coqdocnoindent
\coqdockw{Fixpoint} \coqdef{SECF.Imp.no whiles}{no\_whiles}{\coqdocdefinition{no\_whiles}} (\coqdef{SECF.Imp.c:257}{c}{\coqdocbinder{c}} : \coqref{SECF.Imp.com}{\coqdocinductive{com}}) : \coqexternalref{bool}{http://rocq-prover.org/doc/V9.1.1/stdlib//Corelib.Init.Datatypes}{\coqdocinductive{bool}} :=\coqdoceol
\coqdocindent{1.00em}
\coqdockw{match} \coqref{SECF.Imp.c:257}{\coqdocvariable{c}} \coqdockw{with}\coqdoceol
\coqdocindent{1.00em}
\ensuremath{|} \coqref{SECF.Imp.::com scope:'<x7B' x 'x7D>'}{\coqdocnotation{<\{}} \coqref{SECF.Imp.:com:com scope:'skip'}{\coqdocnotation{skip}} \coqref{SECF.Imp.::com scope:'<x7B' x 'x7D>'}{\coqdocnotation{\}>}} \ensuremath{\Rightarrow}\coqdoceol
\coqdocindent{3.00em}
\coqref{SECF.Imp.:::'true'}{\coqdocnotation{true}}\coqdoceol
\coqdocindent{1.00em}
\ensuremath{|} \coqref{SECF.Imp.::com scope:'<x7B' x 'x7D>'}{\coqdocnotation{<\{}} \coqdocvar{\_} \coqref{SECF.Imp.:com:com scope:x ':=' x}{\coqdocnotation{:=}} \coqdocvar{\_} \coqref{SECF.Imp.::com scope:'<x7B' x 'x7D>'}{\coqdocnotation{\}>}} \ensuremath{\Rightarrow}\coqdoceol
\coqdocindent{3.00em}
\coqref{SECF.Imp.:::'true'}{\coqdocnotation{true}}\coqdoceol
\coqdocindent{1.00em}
\ensuremath{|} \coqref{SECF.Imp.::com scope:'<x7B' x 'x7D>'}{\coqdocnotation{<\{}} \coqdocvar{c1} \coqref{SECF.Imp.:com:com scope:x ';' x}{\coqdocnotation{;}} \coqdocvar{c2} \coqref{SECF.Imp.::com scope:'<x7B' x 'x7D>'}{\coqdocnotation{\}>}} \ensuremath{\Rightarrow}\coqdoceol
\coqdocindent{3.00em}
\coqexternalref{andb}{http://rocq-prover.org/doc/V9.1.1/stdlib//Corelib.Init.Datatypes}{\coqdocdefinition{andb}} (\coqref{SECF.Imp.no whiles:258}{\coqdocdefinition{no\_whiles}} \coqdocvar{c1}) (\coqref{SECF.Imp.no whiles:258}{\coqdocdefinition{no\_whiles}} \coqdocvar{c2})\coqdoceol
\coqdocindent{1.00em}
\ensuremath{|} \coqref{SECF.Imp.::com scope:'<x7B' x 'x7D>'}{\coqdocnotation{<\{}} \coqref{SECF.Imp.:com:com scope:'if' x 'then' x 'else' x 'end'}{\coqdocnotation{if}} \coqdocvar{\_} \coqref{SECF.Imp.:com:com scope:'if' x 'then' x 'else' x 'end'}{\coqdocnotation{then}} \coqdocvar{ct} \coqref{SECF.Imp.:com:com scope:'if' x 'then' x 'else' x 'end'}{\coqdocnotation{else}} \coqdocvar{cf} \coqref{SECF.Imp.:com:com scope:'if' x 'then' x 'else' x 'end'}{\coqdocnotation{end}} \coqref{SECF.Imp.::com scope:'<x7B' x 'x7D>'}{\coqdocnotation{\}>}} \ensuremath{\Rightarrow}\coqdoceol
\coqdocindent{3.00em}
\coqexternalref{andb}{http://rocq-prover.org/doc/V9.1.1/stdlib//Corelib.Init.Datatypes}{\coqdocdefinition{andb}} (\coqref{SECF.Imp.no whiles:258}{\coqdocdefinition{no\_whiles}} \coqdocvar{ct}) (\coqref{SECF.Imp.no whiles:258}{\coqdocdefinition{no\_whiles}} \coqdocvar{cf})\coqdoceol
\coqdocindent{1.00em}
\ensuremath{|} \coqref{SECF.Imp.::com scope:'<x7B' x 'x7D>'}{\coqdocnotation{<\{}} \coqref{SECF.Imp.:com:com scope:'while' x 'do' x 'end'}{\coqdocnotation{while}} \coqdocvar{\_} \coqref{SECF.Imp.:com:com scope:'while' x 'do' x 'end'}{\coqdocnotation{do}} \coqdocvar{\_} \coqref{SECF.Imp.:com:com scope:'while' x 'do' x 'end'}{\coqdocnotation{end}} \coqref{SECF.Imp.::com scope:'<x7B' x 'x7D>'}{\coqdocnotation{\}>}}  \ensuremath{\Rightarrow}\coqdoceol
\coqdocindent{3.00em}
\coqref{SECF.Imp.:::'false'}{\coqdocnotation{false}}\coqdoceol
\coqdocindent{1.00em}
\coqdockw{end}.\coqdoceol
\coqdocemptyline
\end{coqdoccode}
This predicate yields \coqexternalref{true}{http://rocq-prover.org/doc/V9.1.1/stdlib//Corelib.Init.Datatypes}{\coqdocconstructor{true}} just on programs that have no while
    loops.  Using \coqdockw{Inductive}, write a property \coqref{SECF.Imp.no whilesR}{\coqdocinductive{no\_whilesR}} such that
    \coqref{SECF.Imp.no whilesR}{\coqdocinductive{no\_whilesR}} \coqdocvar{c} is provable exactly when \coqdocvar{c} is a program with no
    while loops.  Then prove its equivalence with \coqref{SECF.Imp.no whiles}{\coqdocdefinition{no\_whiles}}. 
\begin{coqdoccode}
\coqdocemptyline
\coqdocnoindent
\coqdockw{Inductive} \coqdef{SECF.Imp.no whilesR}{no\_whilesR}{\coqdocinductive{no\_whilesR}}: \coqref{SECF.Imp.com}{\coqdocinductive{com}} \coqexternalref{::type scope:x '->' x}{http://rocq-prover.org/doc/V9.1.1/stdlib//Corelib.Init.Logic}{\coqdocnotation{\ensuremath{\rightarrow}}} \coqdockw{Prop} :=\coqdoceol
\coqdocindent{0.50em}
\coqdoceol
\coqdocnoindent
.\coqdoceol
\coqdocemptyline
\coqdocnoindent
\coqdockw{Theorem} \coqdef{SECF.Imp.no whiles eqv}{no\_whiles\_eqv}{\coqdoclemma{no\_whiles\_eqv}}:\coqdoceol
\coqdocindent{1.00em}
\coqdockw{\ensuremath{\forall}} \coqdef{SECF.Imp.c:262}{c}{\coqdocbinder{c}}, \coqref{SECF.Imp.no whiles}{\coqdocdefinition{no\_whiles}} \coqref{SECF.Imp.c:262}{\coqdocvariable{c}} \coqexternalref{::type scope:x '=' x}{http://rocq-prover.org/doc/V9.1.1/stdlib//Corelib.Init.Logic}{\coqdocnotation{=}} \coqref{SECF.Imp.:::'true'}{\coqdocnotation{true}} \coqexternalref{::type scope:x '<->' x}{http://rocq-prover.org/doc/V9.1.1/stdlib//Corelib.Init.Logic}{\coqdocnotation{\ensuremath{\leftrightarrow}}} \coqref{SECF.Imp.no whilesR}{\coqdocinductive{no\_whilesR}} \coqref{SECF.Imp.c:262}{\coqdocvariable{c}}.\coqdoceol
\coqdocnoindent
\coqdockw{Proof}.\coqdoceol
 \coqdocvar{Admitted}.\coqdoceol
\end{coqdoccode}
\ensuremath{\Box} 

\paragraph{Exercise: 4 stars, standard (no\_whiles\_terminating)}




    Imp programs that don't involve while loops always terminate.
    State and prove a theorem \coqdocvar{no\_whiles\_terminating} that says this.


    Use either \coqref{SECF.Imp.no whiles}{\coqdocdefinition{no\_whiles}} or \coqref{SECF.Imp.no whilesR}{\coqdocinductive{no\_whilesR}}, as you prefer. 
\begin{coqdoccode}
\coqdocemptyline
\coqdocemptyline
\coqdocnoindent
\coqdockw{Definition} \coqdef{SECF.Imp.manual grade for no whiles terminating}{manual\_grade\_for\_no\_whiles\_terminating}{\coqdocdefinition{manual\_grade\_for\_no\_whiles\_terminating}} : \coqexternalref{option}{http://rocq-prover.org/doc/V9.1.1/stdlib//Corelib.Init.Datatypes}{\coqdocinductive{option}} (\coqexternalref{nat}{http://rocq-prover.org/doc/V9.1.1/stdlib//Corelib.Init.Datatypes}{\coqdocinductive{nat}}\coqexternalref{::type scope:x '*' x}{http://rocq-prover.org/doc/V9.1.1/stdlib//Corelib.Init.Datatypes}{\coqdocnotation{\ensuremath{\times}}}\coqdocinductive{string}) := \coqexternalref{None}{http://rocq-prover.org/doc/V9.1.1/stdlib//Corelib.Init.Datatypes}{\coqdocconstructor{None}}.\coqdoceol
\end{coqdoccode}
\ensuremath{\Box} 
\begin{coqdoccode}
\coqdocemptyline
\end{coqdoccode}
\section{Additional Exercises}



\paragraph{Exercise: 3 stars, standard (stack\_compiler)}




    Old HP Calculators, programming languages like Forth and Postscript,
    and abstract machines like the Java Virtual Machine all evaluate
    arithmetic expressions using a \textit{stack}. For instance, the expression


      (2*3)+(3*(4-2))


   would be written as


      2 3 * 3 4 2 - * +


   and evaluated like this (where we show the program being evaluated
   on the right and the contents of the stack on the left):


                  |    2 3 * 3 4 2 - * +
      2           |    3 * 3 4 2 - * +
      3, 2        |    * 3 4 2 - * +
      6           |    3 4 2 - * +
      3, 6        |    4 2 - * +
      4, 3, 6     |    2 - * +
      2, 4, 3, 6  |    - * +
      2, 3, 6     |    * +
      6, 6        |    +
      12          |


  The goal of this exercise is to write a small compiler that
  translates \coqref{SECF.Imp.aexp}{\coqdocinductive{aexp}}s into stack machine instructions.


  The instruction set for our stack language will consist of the
  following instructions:

\begin{itemize}
\item  \coqref{SECF.Imp.SPush}{\coqdocconstructor{SPush}} \coqdocvar{n}: Push the number \coqdocvar{n} on the stack.

\item  \coqref{SECF.Imp.SLoad}{\coqdocconstructor{SLoad}} \coqdocvar{x}: Load the identifier \coqdocvar{x} from the store and push it
                  on the stack

\item  \coqref{SECF.Imp.SPlus}{\coqdocconstructor{SPlus}}:   Pop the two top numbers from the stack, add them, and
                  push the result onto the stack.

\item  \coqref{SECF.Imp.SMinus}{\coqdocconstructor{SMinus}}:  Similar, but subtract the first number from the second.

\item  \coqref{SECF.Imp.SMult}{\coqdocconstructor{SMult}}:   Similar, but multiply. 

\end{itemize}
\begin{coqdoccode}
\coqdocemptyline
\coqdocnoindent
\coqdockw{Inductive} \coqdef{SECF.Imp.sinstr}{sinstr}{\coqdocinductive{sinstr}} : \coqdockw{Type} :=\coqdoceol
\coqdocnoindent
\ensuremath{|} \coqdef{SECF.Imp.SPush}{SPush}{\coqdocconstructor{SPush}} (\coqdef{SECF.Imp.n:265}{n}{\coqdocbinder{n}} : \coqexternalref{nat}{http://rocq-prover.org/doc/V9.1.1/stdlib//Corelib.Init.Datatypes}{\coqdocinductive{nat}})\coqdoceol
\coqdocnoindent
\ensuremath{|} \coqdef{SECF.Imp.SLoad}{SLoad}{\coqdocconstructor{SLoad}} (\coqdef{SECF.Imp.x:266}{x}{\coqdocbinder{x}} : \coqdocinductive{string})\coqdoceol
\coqdocnoindent
\ensuremath{|} \coqdef{SECF.Imp.SPlus}{SPlus}{\coqdocconstructor{SPlus}}\coqdoceol
\coqdocnoindent
\ensuremath{|} \coqdef{SECF.Imp.SMinus}{SMinus}{\coqdocconstructor{SMinus}}\coqdoceol
\coqdocnoindent
\ensuremath{|} \coqdef{SECF.Imp.SMult}{SMult}{\coqdocconstructor{SMult}}.\coqdoceol
\coqdocemptyline
\end{coqdoccode}
Write a function to evaluate programs in the stack language. It
    should take as input a state, a stack represented as a list of
    numbers (top stack item is the head of the list), and a program
    represented as a list of instructions, and it should return the
    stack after executing the program.  Test your function on the
    examples below.


    Note that it is unspecified what to do when encountering an
    \coqref{SECF.Imp.SPlus}{\coqdocconstructor{SPlus}}, \coqref{SECF.Imp.SMinus}{\coqdocconstructor{SMinus}}, or \coqref{SECF.Imp.SMult}{\coqdocconstructor{SMult}} instruction if the stack contains
    fewer than two elements.  In a sense, it is immaterial what we do,
    since a correct compiler will never emit such a malformed program.
    But for sake of later exercises, it would be best to skip the
    offending instruction and continue with the next one.  
\begin{coqdoccode}
\coqdocemptyline
\coqdocnoindent
\coqdockw{Fixpoint} \coqdef{SECF.Imp.s execute}{s\_execute}{\coqdocdefinition{s\_execute}} (\coqdef{SECF.Imp.st:267}{st}{\coqdocbinder{st}} : \coqref{SECF.Imp.state}{\coqdocdefinition{state}}) (\coqdef{SECF.Imp.stack:268}{stack}{\coqdocbinder{stack}} : \coqexternalref{list}{http://rocq-prover.org/doc/V9.1.1/stdlib//Corelib.Init.Datatypes}{\coqdocinductive{list}} \coqexternalref{nat}{http://rocq-prover.org/doc/V9.1.1/stdlib//Corelib.Init.Datatypes}{\coqdocinductive{nat}})\coqdoceol
\coqdocindent{9.50em}
(\coqdef{SECF.Imp.prog:269}{prog}{\coqdocbinder{prog}} : \coqexternalref{list}{http://rocq-prover.org/doc/V9.1.1/stdlib//Corelib.Init.Datatypes}{\coqdocinductive{list}} \coqref{SECF.Imp.sinstr}{\coqdocinductive{sinstr}})\coqdoceol
\coqdocindent{8.50em}
: \coqexternalref{list}{http://rocq-prover.org/doc/V9.1.1/stdlib//Corelib.Init.Datatypes}{\coqdocinductive{list}} \coqexternalref{nat}{http://rocq-prover.org/doc/V9.1.1/stdlib//Corelib.Init.Datatypes}{\coqdocinductive{nat}}\coqdoceol
\coqdocindent{1.00em}
. \coqdocvar{Admitted}.\coqdoceol
\coqdocemptyline
\coqdocnoindent
\coqdockw{Check} \coqref{SECF.Imp.s execute}{\coqdocaxiom{s\_execute}}.\coqdoceol
\coqdocemptyline
\coqdocnoindent
\coqdockw{Example} \coqdef{SECF.Imp.s execute1}{s\_execute1}{\coqdocdefinition{s\_execute1}} :\coqdoceol
\coqdocindent{2.50em}
\coqref{SECF.Imp.s execute}{\coqdocaxiom{s\_execute}} \coqref{SECF.Imp.empty st}{\coqdocdefinition{empty\_st}} \coqdocnotation{[]}\coqdoceol
\coqdocindent{3.50em}
\coqdocnotation{[}\coqref{SECF.Imp.SPush}{\coqdocconstructor{SPush}} 5\coqdocnotation{;} \coqref{SECF.Imp.SPush}{\coqdocconstructor{SPush}} 3\coqdocnotation{;} \coqref{SECF.Imp.SPush}{\coqdocconstructor{SPush}} 1\coqdocnotation{;} \coqref{SECF.Imp.SMinus}{\coqdocconstructor{SMinus}}\coqdocnotation{]}\coqdoceol
\coqdocindent{1.50em}
\coqexternalref{::type scope:x '=' x}{http://rocq-prover.org/doc/V9.1.1/stdlib//Corelib.Init.Logic}{\coqdocnotation{=}} \coqdocnotation{[}2\coqdocnotation{;} 5\coqdocnotation{]}.\coqdoceol
 \coqdocvar{Admitted}.\coqdoceol
\coqdocemptyline
\coqdocnoindent
\coqdockw{Example} \coqdef{SECF.Imp.s execute2}{s\_execute2}{\coqdocdefinition{s\_execute2}} :\coqdoceol
\coqdocindent{2.50em}
\coqref{SECF.Imp.s execute}{\coqdocaxiom{s\_execute}} (\coqref{SECF.Imp.X}{\coqdocdefinition{X}} \coqref{SECF.Imp.:::x '!->' x}{\coqdocnotation{!->}} 3) \coqdocnotation{[}3\coqdocnotation{;}4\coqdocnotation{]}\coqdoceol
\coqdocindent{3.50em}
\coqdocnotation{[}\coqref{SECF.Imp.SPush}{\coqdocconstructor{SPush}} 4\coqdocnotation{;} \coqref{SECF.Imp.SLoad}{\coqdocconstructor{SLoad}} \coqref{SECF.Imp.X}{\coqdocdefinition{X}}\coqdocnotation{;} \coqref{SECF.Imp.SMult}{\coqdocconstructor{SMult}}\coqdocnotation{;} \coqref{SECF.Imp.SPlus}{\coqdocconstructor{SPlus}}\coqdocnotation{]}\coqdoceol
\coqdocindent{1.50em}
\coqexternalref{::type scope:x '=' x}{http://rocq-prover.org/doc/V9.1.1/stdlib//Corelib.Init.Logic}{\coqdocnotation{=}} \coqdocnotation{[}15\coqdocnotation{;} 4\coqdocnotation{]}.\coqdoceol
 \coqdocvar{Admitted}.\coqdoceol
\coqdocemptyline
\end{coqdoccode}
Next, write a function that compiles an \coqref{SECF.Imp.aexp}{\coqdocinductive{aexp}} into a stack
    machine program. The effect of running the program should be the
    same as pushing the value of the expression on the stack. 
\begin{coqdoccode}
\coqdocemptyline
\coqdocnoindent
\coqdockw{Fixpoint} \coqdef{SECF.Imp.s compile}{s\_compile}{\coqdocdefinition{s\_compile}} (\coqdef{SECF.Imp.e:271}{e}{\coqdocbinder{e}} : \coqref{SECF.Imp.aexp}{\coqdocinductive{aexp}}) : \coqexternalref{list}{http://rocq-prover.org/doc/V9.1.1/stdlib//Corelib.Init.Datatypes}{\coqdocinductive{list}} \coqref{SECF.Imp.sinstr}{\coqdocinductive{sinstr}}\coqdoceol
\coqdocindent{1.00em}
. \coqdocvar{Admitted}.\coqdoceol
\coqdocemptyline
\end{coqdoccode}
After you've defined \coqref{SECF.Imp.s compile}{\coqdocdefinition{s\_compile}}, prove the following to test
    that it works. 
\begin{coqdoccode}
\coqdocemptyline
\coqdocnoindent
\coqdockw{Example} \coqdef{SECF.Imp.s compile1}{s\_compile1}{\coqdocdefinition{s\_compile1}} :\coqdoceol
\coqdocindent{1.00em}
\coqref{SECF.Imp.s compile}{\coqdocaxiom{s\_compile}} \coqref{SECF.Imp.::com scope:'<x7B' x 'x7D>'}{\coqdocnotation{<\{}} \coqref{SECF.Imp.X}{\coqdocdefinition{X}} \coqref{SECF.Imp.:com::x '-' x}{\coqdocnotation{-}} \coqref{SECF.Imp.:com::'(' x ')'}{\coqdocnotation{(}}2 \coqref{SECF.Imp.:com::x '*' x}{\coqdocnotation{\ensuremath{\times}}} \coqref{SECF.Imp.Y}{\coqdocdefinition{Y}}\coqref{SECF.Imp.:com::'(' x ')'}{\coqdocnotation{)}} \coqref{SECF.Imp.::com scope:'<x7B' x 'x7D>'}{\coqdocnotation{\}>}}\coqdoceol
\coqdocindent{1.00em}
\coqexternalref{::type scope:x '=' x}{http://rocq-prover.org/doc/V9.1.1/stdlib//Corelib.Init.Logic}{\coqdocnotation{=}} \coqdocnotation{[}\coqref{SECF.Imp.SLoad}{\coqdocconstructor{SLoad}} \coqref{SECF.Imp.X}{\coqdocdefinition{X}}\coqdocnotation{;} \coqref{SECF.Imp.SPush}{\coqdocconstructor{SPush}} 2\coqdocnotation{;} \coqref{SECF.Imp.SLoad}{\coqdocconstructor{SLoad}} \coqref{SECF.Imp.Y}{\coqdocdefinition{Y}}\coqdocnotation{;} \coqref{SECF.Imp.SMult}{\coqdocconstructor{SMult}}\coqdocnotation{;} \coqref{SECF.Imp.SMinus}{\coqdocconstructor{SMinus}}\coqdocnotation{]}.\coqdoceol
 \coqdocvar{Admitted}.\coqdoceol
\end{coqdoccode}
\ensuremath{\Box} 

\paragraph{Exercise: 3 stars, standard (execute\_app)}



 Execution can be decomposed in the following sense: executing
    stack program \coqref{SECF.Equiv.Himp.p1}{\coqdocdefinition{p1}} ++ \coqref{SECF.Equiv.Himp.p2}{\coqdocdefinition{p2}} is the same as executing \coqref{SECF.Equiv.Himp.p1}{\coqdocdefinition{p1}}, taking
    the resulting stack, and executing \coqref{SECF.Equiv.Himp.p2}{\coqdocdefinition{p2}} from that stack. Prove
    that fact. 
\begin{coqdoccode}
\coqdocemptyline
\coqdocnoindent
\coqdockw{Theorem} \coqdef{SECF.Imp.execute app}{execute\_app}{\coqdoclemma{execute\_app}} : \coqdockw{\ensuremath{\forall}} \coqdef{SECF.Imp.st:273}{st}{\coqdocbinder{st}} \coqdef{SECF.Imp.p1:274}{p1}{\coqdocbinder{p1}} \coqdef{SECF.Imp.p2:275}{p2}{\coqdocbinder{p2}} \coqdef{SECF.Imp.stack:276}{stack}{\coqdocbinder{stack}},\coqdoceol
\coqdocindent{1.00em}
\coqref{SECF.Imp.s execute}{\coqdocaxiom{s\_execute}} \coqref{SECF.Imp.st:273}{\coqdocvariable{st}} \coqref{SECF.Imp.stack:276}{\coqdocvariable{stack}} (\coqref{SECF.Imp.p1:274}{\coqdocvariable{p1}} \coqexternalref{::list scope:x '++' x}{http://rocq-prover.org/doc/V9.1.1/stdlib//Corelib.Init.Datatypes}{\coqdocnotation{++}} \coqref{SECF.Imp.p2:275}{\coqdocvariable{p2}}) \coqexternalref{::type scope:x '=' x}{http://rocq-prover.org/doc/V9.1.1/stdlib//Corelib.Init.Logic}{\coqdocnotation{=}} \coqref{SECF.Imp.s execute}{\coqdocaxiom{s\_execute}} \coqref{SECF.Imp.st:273}{\coqdocvariable{st}} (\coqref{SECF.Imp.s execute}{\coqdocaxiom{s\_execute}} \coqref{SECF.Imp.st:273}{\coqdocvariable{st}} \coqref{SECF.Imp.stack:276}{\coqdocvariable{stack}} \coqref{SECF.Imp.p1:274}{\coqdocvariable{p1}}) \coqref{SECF.Imp.p2:275}{\coqdocvariable{p2}}.\coqdoceol
\coqdocnoindent
\coqdockw{Proof}.\coqdoceol
 \coqdocvar{Admitted}.\coqdoceol
\coqdocemptyline
\end{coqdoccode}
\ensuremath{\Box} 

\paragraph{Exercise: 3 stars, standard (stack\_compiler\_correct)}



 Now we'll prove the correctness of the compiler implemented in the
    previous exercise.  Begin by proving the following lemma. If it
    becomes difficult, consider whether your implementation of
    \coqref{SECF.Imp.s execute}{\coqdocdefinition{s\_execute}} or \coqref{SECF.Imp.s compile}{\coqdocdefinition{s\_compile}} could be simplified. 
\begin{coqdoccode}
\coqdocemptyline
\coqdocnoindent
\coqdockw{Lemma} \coqdef{SECF.Imp.s compile correct aux}{s\_compile\_correct\_aux}{\coqdoclemma{s\_compile\_correct\_aux}} : \coqdockw{\ensuremath{\forall}} \coqdef{SECF.Imp.st:277}{st}{\coqdocbinder{st}} \coqdef{SECF.Imp.e:278}{e}{\coqdocbinder{e}} \coqdef{SECF.Imp.stack:279}{stack}{\coqdocbinder{stack}},\coqdoceol
\coqdocindent{1.00em}
\coqref{SECF.Imp.s execute}{\coqdocaxiom{s\_execute}} \coqref{SECF.Imp.st:277}{\coqdocvariable{st}} \coqref{SECF.Imp.stack:279}{\coqdocvariable{stack}} (\coqref{SECF.Imp.s compile}{\coqdocaxiom{s\_compile}} \coqref{SECF.Imp.e:278}{\coqdocvariable{e}}) \coqexternalref{::type scope:x '=' x}{http://rocq-prover.org/doc/V9.1.1/stdlib//Corelib.Init.Logic}{\coqdocnotation{=}} \coqref{SECF.Imp.aeval}{\coqdocdefinition{aeval}} \coqref{SECF.Imp.st:277}{\coqdocvariable{st}} \coqref{SECF.Imp.e:278}{\coqdocvariable{e}} \coqexternalref{::list scope:x '::' x}{http://rocq-prover.org/doc/V9.1.1/stdlib//Corelib.Init.Datatypes}{\coqdocnotation{::}} \coqref{SECF.Imp.stack:279}{\coqdocvariable{stack}}.\coqdoceol
\coqdocnoindent
\coqdockw{Proof}.\coqdoceol
 \coqdocvar{Admitted}.\coqdoceol
\coqdocemptyline
\end{coqdoccode}
The main theorem should be a very easy corollary of that lemma. 
\begin{coqdoccode}
\coqdocemptyline
\coqdocnoindent
\coqdockw{Theorem} \coqdef{SECF.Imp.s compile correct}{s\_compile\_correct}{\coqdoclemma{s\_compile\_correct}} : \coqdockw{\ensuremath{\forall}} (\coqdef{SECF.Imp.st:280}{st}{\coqdocbinder{st}} : \coqref{SECF.Imp.state}{\coqdocdefinition{state}}) (\coqdef{SECF.Imp.e:281}{e}{\coqdocbinder{e}} : \coqref{SECF.Imp.aexp}{\coqdocinductive{aexp}}),\coqdoceol
\coqdocindent{1.00em}
\coqref{SECF.Imp.s execute}{\coqdocaxiom{s\_execute}} \coqref{SECF.Imp.st:280}{\coqdocvariable{st}} \coqdocnotation{[]} (\coqref{SECF.Imp.s compile}{\coqdocaxiom{s\_compile}} \coqref{SECF.Imp.e:281}{\coqdocvariable{e}}) \coqexternalref{::type scope:x '=' x}{http://rocq-prover.org/doc/V9.1.1/stdlib//Corelib.Init.Logic}{\coqdocnotation{=}} \coqdocnotation{[} \coqref{SECF.Imp.aeval}{\coqdocdefinition{aeval}} \coqref{SECF.Imp.st:280}{\coqdocvariable{st}} \coqref{SECF.Imp.e:281}{\coqdocvariable{e}} \coqdocnotation{]}.\coqdoceol
\coqdocnoindent
\coqdockw{Proof}.\coqdoceol
 \coqdocvar{Admitted}.\coqdoceol
\coqdocemptyline
\end{coqdoccode}
\ensuremath{\Box} 

\paragraph{Exercise: 3 stars, standard, optional (short\_circuit)}




    Most modern programming languages use a ``short-circuit'' evaluation
    rule for boolean \coqexternalref{and}{http://rocq-prover.org/doc/V9.1.1/stdlib//Corelib.Init.Logic}{\coqdocinductive{and}}: to evaluate \coqref{SECF.SpecCT.BAnd}{\coqdocdefinition{BAnd}} \coqdocvar{b1} \coqdocvar{b2}, first evaluate
    \coqdocvar{b1}.  If it evaluates to \coqexternalref{false}{http://rocq-prover.org/doc/V9.1.1/stdlib//Corelib.Init.Datatypes}{\coqdocconstructor{false}}, then the entire \coqref{SECF.SpecCT.BAnd}{\coqdocdefinition{BAnd}}
    expression evaluates to \coqexternalref{false}{http://rocq-prover.org/doc/V9.1.1/stdlib//Corelib.Init.Datatypes}{\coqdocconstructor{false}} immediately, without evaluating
    \coqdocvar{b2}.  Otherwise, \coqdocvar{b2} is evaluated to determine the result of the
    \coqref{SECF.SpecCT.BAnd}{\coqdocdefinition{BAnd}} expression.


    Write an alternate version of \coqref{SECF.Imp.beval}{\coqdocdefinition{beval}} that performs short-circuit
    evaluation of \coqref{SECF.SpecCT.BAnd}{\coqdocdefinition{BAnd}} in this manner, and prove that it is
    equivalent to \coqref{SECF.Imp.beval}{\coqdocdefinition{beval}}.  (N.b. This is only true because expression
    evaluation in Imp is rather simple.  In a bigger language where
    evaluating an expression might diverge, the short-circuiting \coqref{SECF.SpecCT.BAnd}{\coqdocdefinition{BAnd}}
    would \textit{not} be equivalent to the original, since it would make more
    programs terminate.) 
\begin{coqdoccode}
\coqdocemptyline
\coqdocemptyline
\coqdocnoindent
\coqdockw{Module} \coqdef{SECF.Imp.BreakImp}{BreakImp}{\coqdocmodule{BreakImp}}.\coqdoceol
\end{coqdoccode}
\paragraph{Exercise: 4 stars, standard, optional (break\_imp)}




    Imperative languages like C and Java often include a \coqdocvar{break} or
    similar statement for interrupting the execution of loops. In this
    exercise we consider how to add \coqdocvar{break} to Imp.  First, we need to
    enrich the language of commands with an additional case. 
\begin{coqdoccode}
\coqdocemptyline
\coqdocnoindent
\coqdockw{Inductive} \coqdef{SECF.Imp.BreakImp.com}{com}{\coqdocinductive{com}} : \coqdockw{Type} :=\coqdoceol
\coqdocindent{1.00em}
\ensuremath{|} \coqdef{SECF.Imp.BreakImp.CSkip}{CSkip}{\coqdocconstructor{CSkip}}\coqdoceol
\coqdocindent{1.00em}
\ensuremath{|} \coqdef{SECF.Imp.BreakImp.CBreak}{CBreak}{\coqdocconstructor{CBreak}}                        \coqdoceol
\coqdocindent{1.00em}
\ensuremath{|} \coqdef{SECF.Imp.BreakImp.CAsgn}{CAsgn}{\coqdocconstructor{CAsgn}} (\coqdef{SECF.Imp.x:284}{x}{\coqdocbinder{x}} : \coqdocinductive{string}) (\coqdef{SECF.Imp.a:285}{a}{\coqdocbinder{a}} : \coqref{SECF.Imp.aexp}{\coqdocinductive{aexp}})\coqdoceol
\coqdocindent{1.00em}
\ensuremath{|} \coqdef{SECF.Imp.BreakImp.CSeq}{CSeq}{\coqdocconstructor{CSeq}} (\coqdef{SECF.Imp.c1:286}{c1}{\coqdocbinder{c1}} \coqdef{SECF.Imp.c2:287}{c2}{\coqdocbinder{c2}} : \coqref{SECF.Imp.com:282}{\coqdocinductive{com}})\coqdoceol
\coqdocindent{1.00em}
\ensuremath{|} \coqdef{SECF.Imp.BreakImp.CIf}{CIf}{\coqdocconstructor{CIf}} (\coqdef{SECF.Imp.b:288}{b}{\coqdocbinder{b}} : \coqref{SECF.Imp.bexp}{\coqdocinductive{bexp}}) (\coqdef{SECF.Imp.c1:289}{c1}{\coqdocbinder{c1}} \coqdef{SECF.Imp.c2:290}{c2}{\coqdocbinder{c2}} : \coqref{SECF.Imp.com:282}{\coqdocinductive{com}})\coqdoceol
\coqdocindent{1.00em}
\ensuremath{|} \coqdef{SECF.Imp.BreakImp.CWhile}{CWhile}{\coqdocconstructor{CWhile}} (\coqdef{SECF.Imp.b:291}{b}{\coqdocbinder{b}} : \coqref{SECF.Imp.bexp}{\coqdocinductive{bexp}}) (\coqdef{SECF.Imp.c:292}{c}{\coqdocbinder{c}} : \coqref{SECF.Imp.com:282}{\coqdocinductive{com}}).\coqdoceol
\coqdocemptyline
\coqdocnoindent
\coqdockw{Notation} \coqdef{SECF.Imp.BreakImp.:com:com scope:'break'}{"}{"}'break'" := \coqref{SECF.Imp.BreakImp.CBreak}{\coqdocconstructor{CBreak}} (\coqdoctac{in} \coqdocvar{custom} \coqdocvar{com} \coqdoctac{at} \coqdockw{level} 0) : \coqdocvar{com\_scope}.\coqdoceol
\coqdocnoindent
\coqdockw{Notation} \coqdef{SECF.Imp.BreakImp.:com:com scope:'skip'}{"}{"}'skip'"  := \coqref{SECF.Imp.BreakImp.CSkip}{\coqdocconstructor{CSkip}}\coqdoceol
\coqdocindent{1.00em}
(\coqdoctac{in} \coqdocvar{custom} \coqdocvar{com} \coqdoctac{at} \coqdockw{level} 0) : \coqdocvar{com\_scope}.\coqdoceol
\coqdocnoindent
\coqdockw{Notation} \coqdef{SECF.Imp.BreakImp.:com:com scope:x ':=' x}{"}{"}x := y"  := (\coqref{SECF.Imp.BreakImp.CAsgn}{\coqdocconstructor{CAsgn}} \coqdocvar{x} \coqdocvar{y})\coqdoceol
\coqdocindent{1.00em}
(\coqdoctac{in} \coqdocvar{custom} \coqdocvar{com} \coqdoctac{at} \coqdockw{level} 0, \coqdocvar{x} \coqdockw{constr} \coqdoctac{at} \coqdockw{level} 0, \coqdocvar{y} \coqdoctac{at} \coqdockw{level} 85, \coqdockw{no} \coqdockw{associativity},\coqdoceol
\coqdocindent{2.00em}
\coqdocvar{format} "x  :=  y") : \coqdocvar{com\_scope}.\coqdoceol
\coqdocnoindent
\coqdockw{Notation} \coqdef{SECF.Imp.BreakImp.:com:com scope:x ';' x}{"}{"}x ; y" := (\coqref{SECF.Imp.BreakImp.CSeq}{\coqdocconstructor{CSeq}} \coqdocvar{x} \coqdocvar{y})\coqdoceol
\coqdocindent{1.00em}
(\coqdoctac{in} \coqdocvar{custom} \coqdocvar{com} \coqdoctac{at} \coqdockw{level} 90,\coqdoceol
\coqdocindent{2.00em}
\coqdoctac{right} \coqdockw{associativity},\coqdoceol
\coqdocindent{2.00em}
\coqdocvar{format} "'[v' x ; '/' y ']'") : \coqdocvar{com\_scope}.\coqdoceol
\coqdocnoindent
\coqdockw{Notation} \coqdef{SECF.Imp.BreakImp.:com:com scope:'if' x 'then' x 'else' x 'end'}{"}{"}'if' x 'then' y 'else' z 'end'" := (\coqref{SECF.Imp.BreakImp.CIf}{\coqdocconstructor{CIf}} \coqdocvar{x} \coqdocvar{y} \coqdocvar{z})\coqdoceol
\coqdocindent{1.00em}
(\coqdoctac{in} \coqdocvar{custom} \coqdocvar{com} \coqdoctac{at} \coqdockw{level} 89, \coqdocvar{x} \coqdoctac{at} \coqdockw{level} 99, \coqdocvar{y} \coqdoctac{at} \coqdockw{level} 99, \coqdocvar{z} \coqdoctac{at} \coqdockw{level} 99,\coqdoceol
\coqdocindent{2.00em}
\coqdocvar{format} "'[v' 'if'  x  'then' '/  ' y '/' 'else' '/  ' z '/' 'end' ']'") : \coqdocvar{com\_scope}.\coqdoceol
\coqdocnoindent
\coqdockw{Notation} \coqdef{SECF.Imp.BreakImp.:com:com scope:'while' x 'do' x 'end'}{"}{"}'while' x 'do' y 'end'" := (\coqref{SECF.Imp.BreakImp.CWhile}{\coqdocconstructor{CWhile}} \coqdocvar{x} \coqdocvar{y})\coqdoceol
\coqdocindent{1.00em}
(\coqdoctac{in} \coqdocvar{custom} \coqdocvar{com} \coqdoctac{at} \coqdockw{level} 89, \coqdocvar{x} \coqdoctac{at} \coqdockw{level} 99, \coqdocvar{y} \coqdoctac{at} \coqdockw{level} 99,\coqdoceol
\coqdocindent{2.00em}
\coqdocvar{format} "'[v' 'while'  x  'do' '/  ' y '/' 'end' ']'") : \coqdocvar{com\_scope}.\coqdoceol
\coqdocemptyline
\end{coqdoccode}
Next, we need to define the behavior of \coqdocvar{break}.  Informally,
    whenever \coqdocvar{break} is executed in a sequence of commands, it stops
    the execution of that sequence and signals that the innermost
    enclosing loop should terminate.  (If there aren't any
    enclosing loops, then the whole program simply terminates.)  The
    final state should be the same as the one in which the \coqdocvar{break}
    statement was executed.


    One important point is what to do when there are multiple loops
    enclosing a given \coqdocvar{break}. In those cases, \coqdocvar{break} should only
    terminate the \textit{innermost} loop. Thus, after executing the
    following...


       X := 0;
       Y := 1;
       while 0 <> Y do
         while true do
           break
         end;
         X := 1;
         Y := Y - 1
       end


    ... the value of \coqref{SECF.SpecCT.X}{\coqdocdefinition{X}} should be 1, and not 0.


    One way of expressing this behavior is to add another parameter to
    the evaluation relation that specifies whether evaluation of a
    command executes a \coqdocvar{break} statement: 
\begin{coqdoccode}
\coqdocemptyline
\coqdocnoindent
\coqdockw{Inductive} \coqdef{SECF.Imp.BreakImp.result}{result}{\coqdocinductive{result}} : \coqdockw{Type} :=\coqdoceol
\coqdocindent{1.00em}
\ensuremath{|} \coqdef{SECF.Imp.BreakImp.SContinue}{SContinue}{\coqdocconstructor{SContinue}}\coqdoceol
\coqdocindent{1.00em}
\ensuremath{|} \coqdef{SECF.Imp.BreakImp.SBreak}{SBreak}{\coqdocconstructor{SBreak}}.\coqdoceol
\coqdocemptyline
\coqdocnoindent
\coqdockw{Reserved Notation}\coqdoceol
\coqdocindent{4.50em}
"st0 '=[' c ']=>' st1 '/' s"\coqdoceol
\coqdocindent{4.50em}
(\coqdoctac{at} \coqdockw{level} 40, \coqdocvar{c} \coqdocvar{custom} \coqdocvar{com} \coqdoctac{at} \coqdockw{level} 99,\coqdoceol
\coqdocindent{5.00em}
\coqdocvar{st0} \coqdockw{constr}, \coqdocvar{st1} \coqdockw{constr} \coqdoctac{at} \coqdocvar{next} \coqdockw{level},\coqdoceol
\coqdocindent{5.00em}
\coqdocvar{format} "'[hv' st0  =[ '/  ' '[' c ']' '/' ]=>  st1 / s ']'").\coqdoceol
\coqdocemptyline
\end{coqdoccode}
Intuitively, \coqdocvar{st} =[ \coqdocvar{c} ]=> \coqdocvar{st'} / \coqdocvar{s} means that, if \coqdocvar{c} is started in
    state \coqdocvar{st}, then it terminates in state \coqdocvar{st'} and either signals
    that the innermost surrounding loop (or the whole program) should
    exit immediately (\coqdocvar{s} = \coqref{SECF.Imp.BreakImp.SBreak}{\coqdocconstructor{SBreak}}) or that execution should continue
    normally (\coqdocvar{s} = \coqref{SECF.Imp.BreakImp.SContinue}{\coqdocconstructor{SContinue}}).


    The definition of the ``\coqdocvar{st} =[ \coqdocvar{c} ]=> \coqdocvar{st'} / \coqdocvar{s}'' relation is very
    similar to the one we gave above for the regular evaluation
    relation (\coqdocvar{st} =[ \coqdocvar{c} ]=> \coqdocvar{st'}) -- we just need to handle the
    termination signals appropriately:



\begin{itemize}
\item  If the command is \coqdocvar{skip}, then the state doesn't change and
      execution of any enclosing loop can continue normally.



\item  If the command is \coqdocvar{break}, the state stays unchanged but we
      signal a \coqref{SECF.Imp.BreakImp.SBreak}{\coqdocconstructor{SBreak}}.



\item  If the command is an assignment, then we update the binding for
      that variable in the state accordingly and signal that execution
      can continue normally.



\item  If the command is of the form \coqdockw{if} \coqdocvar{b} \coqdockw{then} \coqdocvar{c1} \coqdockw{else} \coqdocvar{c2} \coqdockw{end}, then
      the state is updated as in the original semantics of Imp, except
      that we also propagate the signal from the execution of
      whichever branch was taken.



\item  If the command is a sequence \coqdocvar{c1} ; \coqdocvar{c2}, we first execute
      \coqdocvar{c1}.  If this yields a \coqref{SECF.Imp.BreakImp.SBreak}{\coqdocconstructor{SBreak}}, we skip the execution of \coqdocvar{c2}
      and propagate the \coqref{SECF.Imp.BreakImp.SBreak}{\coqdocconstructor{SBreak}} signal to the surrounding context;
      the resulting state is the same as the one obtained by
      executing \coqdocvar{c1} alone. Otherwise, we execute \coqdocvar{c2} on the state
      obtained after executing \coqdocvar{c1}, and propagate the signal
      generated there.



\item  Finally, for a loop of the form \coqdocvar{while} \coqdocvar{b} \coqdoctac{do} \coqdocvar{c} \coqdockw{end}, the
      semantics is almost the same as before. The only difference is
      that, when \coqdocvar{b} evaluates to true, we execute \coqdocvar{c} and check the
      signal that it raises.  If that signal is \coqref{SECF.Imp.BreakImp.SContinue}{\coqdocconstructor{SContinue}}, then the
      execution proceeds as in the original semantics. Otherwise, we
      stop the execution of the loop, and the resulting state is the
      same as the one resulting from the execution of the current
      iteration.  In either case, since \coqdocvar{break} only terminates the
      innermost loop, \coqdocvar{while} signals \coqref{SECF.Imp.BreakImp.SContinue}{\coqdocconstructor{SContinue}}. 
\end{itemize}


 Based on the above description, complete the definition of the
    \coqref{SECF.Hoare.HoareAssertAssume.ceval}{\coqdocinductive{ceval}} relation. 
\begin{coqdoccode}
\coqdocemptyline
\coqdocnoindent
\coqdockw{Inductive} \coqdef{SECF.Imp.BreakImp.ceval}{ceval}{\coqdocinductive{ceval}} : \coqref{SECF.Imp.BreakImp.com}{\coqdocinductive{com}} \coqexternalref{::type scope:x '->' x}{http://rocq-prover.org/doc/V9.1.1/stdlib//Corelib.Init.Logic}{\coqdocnotation{\ensuremath{\rightarrow}}} \coqref{SECF.Imp.state}{\coqdocdefinition{state}} \coqexternalref{::type scope:x '->' x}{http://rocq-prover.org/doc/V9.1.1/stdlib//Corelib.Init.Logic}{\coqdocnotation{\ensuremath{\rightarrow}}} \coqref{SECF.Imp.BreakImp.result}{\coqdocinductive{result}} \coqexternalref{::type scope:x '->' x}{http://rocq-prover.org/doc/V9.1.1/stdlib//Corelib.Init.Logic}{\coqdocnotation{\ensuremath{\rightarrow}}} \coqref{SECF.Imp.state}{\coqdocdefinition{state}} \coqexternalref{::type scope:x '->' x}{http://rocq-prover.org/doc/V9.1.1/stdlib//Corelib.Init.Logic}{\coqdocnotation{\ensuremath{\rightarrow}}} \coqdockw{Prop} :=\coqdoceol
\coqdocindent{1.00em}
\ensuremath{|} \coqdef{SECF.Imp.BreakImp.E Skip}{E\_Skip}{\coqdocconstructor{E\_Skip}} : \coqdockw{\ensuremath{\forall}} \coqdef{SECF.Imp.st:297}{st}{\coqdocbinder{st}},\coqdoceol
\coqdocindent{3.00em}
\coqref{SECF.Imp.st:297}{\coqdocvariable{st}} \coqref{SECF.Imp.BreakImp.:::x '=[' x ']=>' x '/' x}{\coqdocnotation{=[}} \coqref{SECF.Imp.BreakImp.CSkip}{\coqdocconstructor{CSkip}} \coqref{SECF.Imp.BreakImp.:::x '=[' x ']=>' x '/' x}{\coqdocnotation{]=>}} \coqref{SECF.Imp.st:297}{\coqdocvariable{st}} \coqref{SECF.Imp.BreakImp.:::x '=[' x ']=>' x '/' x}{\coqdocnotation{/}} \coqref{SECF.Imp.BreakImp.SContinue}{\coqdocconstructor{SContinue}}\coqdoceol
\coqdocindent{1.00em}
\coqdoceol
\coqdocnoindent
\coqdoceol
\coqdocindent{1.00em}
\coqdockw{where} \coqdef{SECF.Imp.BreakImp.:::x '=[' x ']=>' x '/' x}{"}{"}st '=[' c ']=>' st' '/' s" := (\coqref{SECF.Imp.ceval:296}{\coqdocinductive{ceval}} \coqdocvar{c} \coqdocvar{st} \coqdocvar{s} \coqdocvar{st'}).\coqdoceol
\coqdocemptyline
\end{coqdoccode}
Now prove the following properties of your definition of \coqref{SECF.Hoare.HoareAssertAssume.ceval}{\coqdocinductive{ceval}}: 
\begin{coqdoccode}
\coqdocemptyline
\coqdocnoindent
\coqdockw{Theorem} \coqdef{SECF.Imp.BreakImp.break ignore}{break\_ignore}{\coqdoclemma{break\_ignore}} : \coqdockw{\ensuremath{\forall}} \coqdef{SECF.Imp.c:298}{c}{\coqdocbinder{c}} \coqdef{SECF.Imp.st:299}{st}{\coqdocbinder{st}} \coqdef{SECF.Imp.st':300}{st'}{\coqdocbinder{st'}} \coqdef{SECF.Imp.s:301}{s}{\coqdocbinder{s}},\coqdoceol
\coqdocindent{2.50em}
\coqref{SECF.Imp.st:299}{\coqdocvariable{st}} \coqref{SECF.Imp.BreakImp.:::x '=[' x ']=>' x '/' x}{\coqdocnotation{=[}} \coqref{SECF.Imp.BreakImp.:com:com scope:'break'}{\coqdocnotation{break}}\coqref{SECF.Imp.BreakImp.:com:com scope:x ';' x}{\coqdocnotation{;}} \coqref{SECF.Imp.c:298}{\coqdocvariable{c}} \coqref{SECF.Imp.BreakImp.:::x '=[' x ']=>' x '/' x}{\coqdocnotation{]=>}} \coqref{SECF.Imp.st':300}{\coqdocvariable{st'}} \coqref{SECF.Imp.BreakImp.:::x '=[' x ']=>' x '/' x}{\coqdocnotation{/}} \coqref{SECF.Imp.s:301}{\coqdocvariable{s}} \coqexternalref{::type scope:x '->' x}{http://rocq-prover.org/doc/V9.1.1/stdlib//Corelib.Init.Logic}{\coqdocnotation{\ensuremath{\rightarrow}}}\coqdoceol
\coqdocindent{2.50em}
\coqref{SECF.Imp.st:299}{\coqdocvariable{st}} \coqexternalref{::type scope:x '=' x}{http://rocq-prover.org/doc/V9.1.1/stdlib//Corelib.Init.Logic}{\coqdocnotation{=}} \coqref{SECF.Imp.st':300}{\coqdocvariable{st'}}.\coqdoceol
\coqdocnoindent
\coqdockw{Proof}.\coqdoceol
 \coqdocvar{Admitted}.\coqdoceol
\coqdocemptyline
\coqdocnoindent
\coqdockw{Theorem} \coqdef{SECF.Imp.BreakImp.while continue}{while\_continue}{\coqdoclemma{while\_continue}} : \coqdockw{\ensuremath{\forall}} \coqdef{SECF.Imp.b:302}{b}{\coqdocbinder{b}} \coqdef{SECF.Imp.c:303}{c}{\coqdocbinder{c}} \coqdef{SECF.Imp.st:304}{st}{\coqdocbinder{st}} \coqdef{SECF.Imp.st':305}{st'}{\coqdocbinder{st'}} \coqdef{SECF.Imp.s:306}{s}{\coqdocbinder{s}},\coqdoceol
\coqdocindent{1.00em}
\coqref{SECF.Imp.st:304}{\coqdocvariable{st}} \coqref{SECF.Imp.BreakImp.:::x '=[' x ']=>' x '/' x}{\coqdocnotation{=[}} \coqref{SECF.Imp.BreakImp.:com:com scope:'while' x 'do' x 'end'}{\coqdocnotation{while}} \coqref{SECF.Imp.b:302}{\coqdocvariable{b}} \coqref{SECF.Imp.BreakImp.:com:com scope:'while' x 'do' x 'end'}{\coqdocnotation{do}} \coqref{SECF.Imp.c:303}{\coqdocvariable{c}} \coqref{SECF.Imp.BreakImp.:com:com scope:'while' x 'do' x 'end'}{\coqdocnotation{end}} \coqref{SECF.Imp.BreakImp.:::x '=[' x ']=>' x '/' x}{\coqdocnotation{]=>}} \coqref{SECF.Imp.st':305}{\coqdocvariable{st'}} \coqref{SECF.Imp.BreakImp.:::x '=[' x ']=>' x '/' x}{\coqdocnotation{/}} \coqref{SECF.Imp.s:306}{\coqdocvariable{s}} \coqexternalref{::type scope:x '->' x}{http://rocq-prover.org/doc/V9.1.1/stdlib//Corelib.Init.Logic}{\coqdocnotation{\ensuremath{\rightarrow}}}\coqdoceol
\coqdocindent{1.00em}
\coqref{SECF.Imp.s:306}{\coqdocvariable{s}} \coqexternalref{::type scope:x '=' x}{http://rocq-prover.org/doc/V9.1.1/stdlib//Corelib.Init.Logic}{\coqdocnotation{=}} \coqref{SECF.Imp.BreakImp.SContinue}{\coqdocconstructor{SContinue}}.\coqdoceol
\coqdocnoindent
\coqdockw{Proof}.\coqdoceol
 \coqdocvar{Admitted}.\coqdoceol
\coqdocemptyline
\coqdocnoindent
\coqdockw{Theorem} \coqdef{SECF.Imp.BreakImp.while stops on break}{while\_stops\_on\_break}{\coqdoclemma{while\_stops\_on\_break}} : \coqdockw{\ensuremath{\forall}} \coqdef{SECF.Imp.b:307}{b}{\coqdocbinder{b}} \coqdef{SECF.Imp.c:308}{c}{\coqdocbinder{c}} \coqdef{SECF.Imp.st:309}{st}{\coqdocbinder{st}} \coqdef{SECF.Imp.st':310}{st'}{\coqdocbinder{st'}},\coqdoceol
\coqdocindent{1.00em}
\coqref{SECF.Imp.beval}{\coqdocdefinition{beval}} \coqref{SECF.Imp.st:309}{\coqdocvariable{st}} \coqref{SECF.Imp.b:307}{\coqdocvariable{b}} \coqexternalref{::type scope:x '=' x}{http://rocq-prover.org/doc/V9.1.1/stdlib//Corelib.Init.Logic}{\coqdocnotation{=}} \coqref{SECF.Imp.:::'true'}{\coqdocnotation{true}} \coqexternalref{::type scope:x '->' x}{http://rocq-prover.org/doc/V9.1.1/stdlib//Corelib.Init.Logic}{\coqdocnotation{\ensuremath{\rightarrow}}}\coqdoceol
\coqdocindent{1.00em}
\coqref{SECF.Imp.st:309}{\coqdocvariable{st}} \coqref{SECF.Imp.BreakImp.:::x '=[' x ']=>' x '/' x}{\coqdocnotation{=[}} \coqref{SECF.Imp.c:308}{\coqdocvariable{c}} \coqref{SECF.Imp.BreakImp.:::x '=[' x ']=>' x '/' x}{\coqdocnotation{]=>}} \coqref{SECF.Imp.st':310}{\coqdocvariable{st'}} \coqref{SECF.Imp.BreakImp.:::x '=[' x ']=>' x '/' x}{\coqdocnotation{/}} \coqref{SECF.Imp.BreakImp.SBreak}{\coqdocconstructor{SBreak}} \coqexternalref{::type scope:x '->' x}{http://rocq-prover.org/doc/V9.1.1/stdlib//Corelib.Init.Logic}{\coqdocnotation{\ensuremath{\rightarrow}}}\coqdoceol
\coqdocindent{1.00em}
\coqref{SECF.Imp.st:309}{\coqdocvariable{st}} \coqref{SECF.Imp.BreakImp.:::x '=[' x ']=>' x '/' x}{\coqdocnotation{=[}} \coqref{SECF.Imp.BreakImp.:com:com scope:'while' x 'do' x 'end'}{\coqdocnotation{while}} \coqref{SECF.Imp.b:307}{\coqdocvariable{b}} \coqref{SECF.Imp.BreakImp.:com:com scope:'while' x 'do' x 'end'}{\coqdocnotation{do}} \coqref{SECF.Imp.c:308}{\coqdocvariable{c}} \coqref{SECF.Imp.BreakImp.:com:com scope:'while' x 'do' x 'end'}{\coqdocnotation{end}} \coqref{SECF.Imp.BreakImp.:::x '=[' x ']=>' x '/' x}{\coqdocnotation{]=>}} \coqref{SECF.Imp.st':310}{\coqdocvariable{st'}} \coqref{SECF.Imp.BreakImp.:::x '=[' x ']=>' x '/' x}{\coqdocnotation{/}} \coqref{SECF.Imp.BreakImp.SContinue}{\coqdocconstructor{SContinue}}.\coqdoceol
\coqdocnoindent
\coqdockw{Proof}.\coqdoceol
 \coqdocvar{Admitted}.\coqdoceol
\coqdocemptyline
\coqdocnoindent
\coqdockw{Theorem} \coqdef{SECF.Imp.BreakImp.seq continue}{seq\_continue}{\coqdoclemma{seq\_continue}} : \coqdockw{\ensuremath{\forall}} \coqdef{SECF.Imp.c1:311}{c1}{\coqdocbinder{c1}} \coqdef{SECF.Imp.c2:312}{c2}{\coqdocbinder{c2}} \coqdef{SECF.Imp.st:313}{st}{\coqdocbinder{st}} \coqdef{SECF.Imp.st':314}{st'}{\coqdocbinder{st'}} \coqdef{SECF.Imp.st'':315}{st'{}'}{\coqdocbinder{st'{}'}},\coqdoceol
\coqdocindent{1.00em}
\coqref{SECF.Imp.st:313}{\coqdocvariable{st}} \coqref{SECF.Imp.BreakImp.:::x '=[' x ']=>' x '/' x}{\coqdocnotation{=[}} \coqref{SECF.Imp.c1:311}{\coqdocvariable{c1}} \coqref{SECF.Imp.BreakImp.:::x '=[' x ']=>' x '/' x}{\coqdocnotation{]=>}} \coqref{SECF.Imp.st':314}{\coqdocvariable{st'}} \coqref{SECF.Imp.BreakImp.:::x '=[' x ']=>' x '/' x}{\coqdocnotation{/}} \coqref{SECF.Imp.BreakImp.SContinue}{\coqdocconstructor{SContinue}} \coqexternalref{::type scope:x '->' x}{http://rocq-prover.org/doc/V9.1.1/stdlib//Corelib.Init.Logic}{\coqdocnotation{\ensuremath{\rightarrow}}}\coqdoceol
\coqdocindent{1.00em}
\coqref{SECF.Imp.st':314}{\coqdocvariable{st'}} \coqref{SECF.Imp.BreakImp.:::x '=[' x ']=>' x '/' x}{\coqdocnotation{=[}} \coqref{SECF.Imp.c2:312}{\coqdocvariable{c2}} \coqref{SECF.Imp.BreakImp.:::x '=[' x ']=>' x '/' x}{\coqdocnotation{]=>}} \coqref{SECF.Imp.st'':315}{\coqdocvariable{st'{}'}} \coqref{SECF.Imp.BreakImp.:::x '=[' x ']=>' x '/' x}{\coqdocnotation{/}} \coqref{SECF.Imp.BreakImp.SContinue}{\coqdocconstructor{SContinue}} \coqexternalref{::type scope:x '->' x}{http://rocq-prover.org/doc/V9.1.1/stdlib//Corelib.Init.Logic}{\coqdocnotation{\ensuremath{\rightarrow}}}\coqdoceol
\coqdocindent{1.00em}
\coqref{SECF.Imp.st:313}{\coqdocvariable{st}} \coqref{SECF.Imp.BreakImp.:::x '=[' x ']=>' x '/' x}{\coqdocnotation{=[}} \coqref{SECF.Imp.c1:311}{\coqdocvariable{c1}} \coqref{SECF.Imp.BreakImp.:com:com scope:x ';' x}{\coqdocnotation{;}} \coqref{SECF.Imp.c2:312}{\coqdocvariable{c2}} \coqref{SECF.Imp.BreakImp.:::x '=[' x ']=>' x '/' x}{\coqdocnotation{]=>}} \coqref{SECF.Imp.st'':315}{\coqdocvariable{st'{}'}} \coqref{SECF.Imp.BreakImp.:::x '=[' x ']=>' x '/' x}{\coqdocnotation{/}} \coqref{SECF.Imp.BreakImp.SContinue}{\coqdocconstructor{SContinue}}.\coqdoceol
\coqdocnoindent
\coqdockw{Proof}.\coqdoceol
 \coqdocvar{Admitted}.\coqdoceol
\coqdocemptyline
\coqdocnoindent
\coqdockw{Theorem} \coqdef{SECF.Imp.BreakImp.seq stops on break}{seq\_stops\_on\_break}{\coqdoclemma{seq\_stops\_on\_break}} : \coqdockw{\ensuremath{\forall}} \coqdef{SECF.Imp.c1:316}{c1}{\coqdocbinder{c1}} \coqdef{SECF.Imp.c2:317}{c2}{\coqdocbinder{c2}} \coqdef{SECF.Imp.st:318}{st}{\coqdocbinder{st}} \coqdef{SECF.Imp.st':319}{st'}{\coqdocbinder{st'}},\coqdoceol
\coqdocindent{1.00em}
\coqref{SECF.Imp.st:318}{\coqdocvariable{st}} \coqref{SECF.Imp.BreakImp.:::x '=[' x ']=>' x '/' x}{\coqdocnotation{=[}} \coqref{SECF.Imp.c1:316}{\coqdocvariable{c1}} \coqref{SECF.Imp.BreakImp.:::x '=[' x ']=>' x '/' x}{\coqdocnotation{]=>}} \coqref{SECF.Imp.st':319}{\coqdocvariable{st'}} \coqref{SECF.Imp.BreakImp.:::x '=[' x ']=>' x '/' x}{\coqdocnotation{/}} \coqref{SECF.Imp.BreakImp.SBreak}{\coqdocconstructor{SBreak}} \coqexternalref{::type scope:x '->' x}{http://rocq-prover.org/doc/V9.1.1/stdlib//Corelib.Init.Logic}{\coqdocnotation{\ensuremath{\rightarrow}}}\coqdoceol
\coqdocindent{1.00em}
\coqref{SECF.Imp.st:318}{\coqdocvariable{st}} \coqref{SECF.Imp.BreakImp.:::x '=[' x ']=>' x '/' x}{\coqdocnotation{=[}} \coqref{SECF.Imp.c1:316}{\coqdocvariable{c1}} \coqref{SECF.Imp.BreakImp.:com:com scope:x ';' x}{\coqdocnotation{;}} \coqref{SECF.Imp.c2:317}{\coqdocvariable{c2}} \coqref{SECF.Imp.BreakImp.:::x '=[' x ']=>' x '/' x}{\coqdocnotation{]=>}} \coqref{SECF.Imp.st':319}{\coqdocvariable{st'}} \coqref{SECF.Imp.BreakImp.:::x '=[' x ']=>' x '/' x}{\coqdocnotation{/}} \coqref{SECF.Imp.BreakImp.SBreak}{\coqdocconstructor{SBreak}}.\coqdoceol
\coqdocnoindent
\coqdockw{Proof}.\coqdoceol
 \coqdocvar{Admitted}.\coqdoceol
\end{coqdoccode}
\ensuremath{\Box} 

\paragraph{Exercise: 3 stars, advanced, optional (while\_break\_true)}


\begin{coqdoccode}
\coqdocnoindent
\coqdockw{Theorem} \coqdef{SECF.Imp.BreakImp.while break true}{while\_break\_true}{\coqdoclemma{while\_break\_true}} : \coqdockw{\ensuremath{\forall}} \coqdef{SECF.Imp.b:320}{b}{\coqdocbinder{b}} \coqdef{SECF.Imp.c:321}{c}{\coqdocbinder{c}} \coqdef{SECF.Imp.st:322}{st}{\coqdocbinder{st}} \coqdef{SECF.Imp.st':323}{st'}{\coqdocbinder{st'}},\coqdoceol
\coqdocindent{1.00em}
\coqref{SECF.Imp.st:322}{\coqdocvariable{st}} \coqref{SECF.Imp.BreakImp.:::x '=[' x ']=>' x '/' x}{\coqdocnotation{=[}} \coqref{SECF.Imp.BreakImp.:com:com scope:'while' x 'do' x 'end'}{\coqdocnotation{while}} \coqref{SECF.Imp.b:320}{\coqdocvariable{b}} \coqref{SECF.Imp.BreakImp.:com:com scope:'while' x 'do' x 'end'}{\coqdocnotation{do}} \coqref{SECF.Imp.c:321}{\coqdocvariable{c}} \coqref{SECF.Imp.BreakImp.:com:com scope:'while' x 'do' x 'end'}{\coqdocnotation{end}} \coqref{SECF.Imp.BreakImp.:::x '=[' x ']=>' x '/' x}{\coqdocnotation{]=>}} \coqref{SECF.Imp.st':323}{\coqdocvariable{st'}} \coqref{SECF.Imp.BreakImp.:::x '=[' x ']=>' x '/' x}{\coqdocnotation{/}} \coqref{SECF.Imp.BreakImp.SContinue}{\coqdocconstructor{SContinue}} \coqexternalref{::type scope:x '->' x}{http://rocq-prover.org/doc/V9.1.1/stdlib//Corelib.Init.Logic}{\coqdocnotation{\ensuremath{\rightarrow}}}\coqdoceol
\coqdocindent{1.00em}
\coqref{SECF.Imp.beval}{\coqdocdefinition{beval}} \coqref{SECF.Imp.st':323}{\coqdocvariable{st'}} \coqref{SECF.Imp.b:320}{\coqdocvariable{b}} \coqexternalref{::type scope:x '=' x}{http://rocq-prover.org/doc/V9.1.1/stdlib//Corelib.Init.Logic}{\coqdocnotation{=}} \coqref{SECF.Imp.:::'true'}{\coqdocnotation{true}} \coqexternalref{::type scope:x '->' x}{http://rocq-prover.org/doc/V9.1.1/stdlib//Corelib.Init.Logic}{\coqdocnotation{\ensuremath{\rightarrow}}}\coqdoceol
\coqdocindent{1.00em}
\coqexternalref{::type scope:'exists' x '..' x ',' x}{http://rocq-prover.org/doc/V9.1.1/stdlib//Corelib.Init.Logic}{\coqdocnotation{\ensuremath{\exists}}} \coqdef{SECF.Imp.st'':324}{st'{}'}{\coqdocbinder{st'{}'}}\coqexternalref{::type scope:'exists' x '..' x ',' x}{http://rocq-prover.org/doc/V9.1.1/stdlib//Corelib.Init.Logic}{\coqdocnotation{,}} \coqref{SECF.Imp.st'':324}{\coqdocvariable{st'{}'}} \coqref{SECF.Imp.BreakImp.:::x '=[' x ']=>' x '/' x}{\coqdocnotation{=[}} \coqref{SECF.Imp.c:321}{\coqdocvariable{c}} \coqref{SECF.Imp.BreakImp.:::x '=[' x ']=>' x '/' x}{\coqdocnotation{]=>}} \coqref{SECF.Imp.st':323}{\coqdocvariable{st'}} \coqref{SECF.Imp.BreakImp.:::x '=[' x ']=>' x '/' x}{\coqdocnotation{/}} \coqref{SECF.Imp.BreakImp.SBreak}{\coqdocconstructor{SBreak}}.\coqdoceol
\coqdocnoindent
\coqdockw{Proof}.\coqdoceol
 \coqdocvar{Admitted}.\coqdoceol
\end{coqdoccode}
\ensuremath{\Box} 

\paragraph{Exercise: 4 stars, advanced, optional (ceval\_deterministic)}


\begin{coqdoccode}
\coqdocnoindent
\coqdockw{Theorem} \coqdef{SECF.Imp.BreakImp.ceval deterministic}{ceval\_deterministic}{\coqdoclemma{ceval\_deterministic}}: \coqdockw{\ensuremath{\forall}} (\coqdef{SECF.Imp.c:325}{c}{\coqdocbinder{c}}:\coqref{SECF.Imp.BreakImp.com}{\coqdocinductive{com}}) \coqdef{SECF.Imp.st:326}{st}{\coqdocbinder{st}} \coqdef{SECF.Imp.st1:327}{st1}{\coqdocbinder{st1}} \coqdef{SECF.Imp.st2:328}{st2}{\coqdocbinder{st2}} \coqdef{SECF.Imp.s1:329}{s1}{\coqdocbinder{s1}} \coqdef{SECF.Imp.s2:330}{s2}{\coqdocbinder{s2}},\coqdoceol
\coqdocindent{2.50em}
\coqref{SECF.Imp.st:326}{\coqdocvariable{st}} \coqref{SECF.Imp.BreakImp.:::x '=[' x ']=>' x '/' x}{\coqdocnotation{=[}} \coqref{SECF.Imp.c:325}{\coqdocvariable{c}} \coqref{SECF.Imp.BreakImp.:::x '=[' x ']=>' x '/' x}{\coqdocnotation{]=>}} \coqref{SECF.Imp.st1:327}{\coqdocvariable{st1}} \coqref{SECF.Imp.BreakImp.:::x '=[' x ']=>' x '/' x}{\coqdocnotation{/}} \coqref{SECF.Imp.s1:329}{\coqdocvariable{s1}} \coqexternalref{::type scope:x '->' x}{http://rocq-prover.org/doc/V9.1.1/stdlib//Corelib.Init.Logic}{\coqdocnotation{\ensuremath{\rightarrow}}}\coqdoceol
\coqdocindent{2.50em}
\coqref{SECF.Imp.st:326}{\coqdocvariable{st}} \coqref{SECF.Imp.BreakImp.:::x '=[' x ']=>' x '/' x}{\coqdocnotation{=[}} \coqref{SECF.Imp.c:325}{\coqdocvariable{c}} \coqref{SECF.Imp.BreakImp.:::x '=[' x ']=>' x '/' x}{\coqdocnotation{]=>}} \coqref{SECF.Imp.st2:328}{\coqdocvariable{st2}} \coqref{SECF.Imp.BreakImp.:::x '=[' x ']=>' x '/' x}{\coqdocnotation{/}} \coqref{SECF.Imp.s2:330}{\coqdocvariable{s2}} \coqexternalref{::type scope:x '->' x}{http://rocq-prover.org/doc/V9.1.1/stdlib//Corelib.Init.Logic}{\coqdocnotation{\ensuremath{\rightarrow}}}\coqdoceol
\coqdocindent{2.50em}
\coqref{SECF.Imp.st1:327}{\coqdocvariable{st1}} \coqexternalref{::type scope:x '=' x}{http://rocq-prover.org/doc/V9.1.1/stdlib//Corelib.Init.Logic}{\coqdocnotation{=}} \coqref{SECF.Imp.st2:328}{\coqdocvariable{st2}} \coqexternalref{::type scope:x '/x5C' x}{http://rocq-prover.org/doc/V9.1.1/stdlib//Corelib.Init.Logic}{\coqdocnotation{\ensuremath{\land}}} \coqref{SECF.Imp.s1:329}{\coqdocvariable{s1}} \coqexternalref{::type scope:x '=' x}{http://rocq-prover.org/doc/V9.1.1/stdlib//Corelib.Init.Logic}{\coqdocnotation{=}} \coqref{SECF.Imp.s2:330}{\coqdocvariable{s2}}.\coqdoceol
\coqdocnoindent
\coqdockw{Proof}.\coqdoceol
 \coqdocvar{Admitted}.\coqdoceol
\coqdocemptyline
\end{coqdoccode}
\ensuremath{\Box} 
\begin{coqdoccode}
\coqdocnoindent
\coqdockw{End} \coqref{SECF.Imp.BreakImp}{\coqdocmodule{BreakImp}}.\coqdoceol
\coqdocemptyline
\end{coqdoccode}
\paragraph{Exercise: 4 stars, standard, optional (add\_for\_loop)}




    Add C-style \coqdockw{for} loops to the language of commands, update the
    \coqref{SECF.Hoare.HoareAssertAssume.ceval}{\coqdocinductive{ceval}} definition to define the semantics of \coqdockw{for} loops, and add
    cases for \coqdockw{for} loops as needed so that all the proofs in this
    file are accepted by Rocq.


    A \coqdockw{for} loop should be parameterized by (a) a statement executed
    initially, (b) a test that is run on each iteration of the loop to
    determine whether the loop should continue, (c) a statement
    executed at the end of each loop iteration, and (d) a statement
    that makes up the body of the loop.  (You don't need to worry
    about making up a concrete Notation for \coqdockw{for} loops, but feel free
    to play with this too if you like.) 
\begin{coqdoccode}
\coqdocemptyline
\coqdocemptyline
\end{coqdoccode}
\end{document}
